% Options for packages loaded elsewhere
\PassOptionsToPackage{unicode}{hyperref}
\PassOptionsToPackage{hyphens}{url}
\documentclass[
  12pt,
  letterpaper,
  oneside]{report}
\usepackage{xcolor}
\usepackage[letterpaper,left=1.5in,right=1in,top=1in,bottom=1in]{geometry}
\usepackage{amsmath,amssymb}
\setcounter{secnumdepth}{3}
\usepackage{iftex}
\ifPDFTeX
  \usepackage[T1]{fontenc}
  \usepackage[utf8]{inputenc}
  \usepackage{textcomp} % provide euro and other symbols
\else % if luatex or xetex
  \usepackage{unicode-math} % this also loads fontspec
  \defaultfontfeatures{Scale=MatchLowercase}
  \defaultfontfeatures[\rmfamily]{Ligatures=TeX,Scale=1}
\fi
\usepackage{lmodern}
\ifPDFTeX\else
  % xetex/luatex font selection
\fi
% Use upquote if available, for straight quotes in verbatim environments
\IfFileExists{upquote.sty}{\usepackage{upquote}}{}
\IfFileExists{microtype.sty}{% use microtype if available
  \usepackage[]{microtype}
  \UseMicrotypeSet[protrusion]{basicmath} % disable protrusion for tt fonts
}{}
\makeatletter
\@ifundefined{KOMAClassName}{% if non-KOMA class
  \IfFileExists{parskip.sty}{%
    \usepackage{parskip}
  }{% else
    \setlength{\parindent}{0pt}
    \setlength{\parskip}{6pt plus 2pt minus 1pt}}
}{% if KOMA class
  \KOMAoptions{parskip=half}}
\makeatother
\setlength{\parindent}{0.5in}
\setlength{\parskip}{0pt}
\usepackage{color}
\usepackage{fancyvrb}
\newcommand{\VerbBar}{|}
\newcommand{\VERB}{\Verb[commandchars=\\\{\}]}
\DefineVerbatimEnvironment{Highlighting}{Verbatim}{commandchars=\\\{\}}
% Add ',fontsize=\small' for more characters per line
\newenvironment{Shaded}{}{}
\newcommand{\AlertTok}[1]{\textcolor[rgb]{1.00,0.00,0.00}{\textbf{#1}}}
\newcommand{\AnnotationTok}[1]{\textcolor[rgb]{0.38,0.63,0.69}{\textbf{\textit{#1}}}}
\newcommand{\AttributeTok}[1]{\textcolor[rgb]{0.49,0.56,0.16}{#1}}
\newcommand{\BaseNTok}[1]{\textcolor[rgb]{0.25,0.63,0.44}{#1}}
\newcommand{\BuiltInTok}[1]{\textcolor[rgb]{0.00,0.50,0.00}{#1}}
\newcommand{\CharTok}[1]{\textcolor[rgb]{0.25,0.44,0.63}{#1}}
\newcommand{\CommentTok}[1]{\textcolor[rgb]{0.38,0.63,0.69}{\textit{#1}}}
\newcommand{\CommentVarTok}[1]{\textcolor[rgb]{0.38,0.63,0.69}{\textbf{\textit{#1}}}}
\newcommand{\ConstantTok}[1]{\textcolor[rgb]{0.53,0.00,0.00}{#1}}
\newcommand{\ControlFlowTok}[1]{\textcolor[rgb]{0.00,0.44,0.13}{\textbf{#1}}}
\newcommand{\DataTypeTok}[1]{\textcolor[rgb]{0.56,0.13,0.00}{#1}}
\newcommand{\DecValTok}[1]{\textcolor[rgb]{0.25,0.63,0.44}{#1}}
\newcommand{\DocumentationTok}[1]{\textcolor[rgb]{0.73,0.13,0.13}{\textit{#1}}}
\newcommand{\ErrorTok}[1]{\textcolor[rgb]{1.00,0.00,0.00}{\textbf{#1}}}
\newcommand{\ExtensionTok}[1]{#1}
\newcommand{\FloatTok}[1]{\textcolor[rgb]{0.25,0.63,0.44}{#1}}
\newcommand{\FunctionTok}[1]{\textcolor[rgb]{0.02,0.16,0.49}{#1}}
\newcommand{\ImportTok}[1]{\textcolor[rgb]{0.00,0.50,0.00}{\textbf{#1}}}
\newcommand{\InformationTok}[1]{\textcolor[rgb]{0.38,0.63,0.69}{\textbf{\textit{#1}}}}
\newcommand{\KeywordTok}[1]{\textcolor[rgb]{0.00,0.44,0.13}{\textbf{#1}}}
\newcommand{\NormalTok}[1]{#1}
\newcommand{\OperatorTok}[1]{\textcolor[rgb]{0.40,0.40,0.40}{#1}}
\newcommand{\OtherTok}[1]{\textcolor[rgb]{0.00,0.44,0.13}{#1}}
\newcommand{\PreprocessorTok}[1]{\textcolor[rgb]{0.74,0.48,0.00}{#1}}
\newcommand{\RegionMarkerTok}[1]{#1}
\newcommand{\SpecialCharTok}[1]{\textcolor[rgb]{0.25,0.44,0.63}{#1}}
\newcommand{\SpecialStringTok}[1]{\textcolor[rgb]{0.73,0.40,0.53}{#1}}
\newcommand{\StringTok}[1]{\textcolor[rgb]{0.25,0.44,0.63}{#1}}
\newcommand{\VariableTok}[1]{\textcolor[rgb]{0.10,0.09,0.49}{#1}}
\newcommand{\VerbatimStringTok}[1]{\textcolor[rgb]{0.25,0.44,0.63}{#1}}
\newcommand{\WarningTok}[1]{\textcolor[rgb]{0.38,0.63,0.69}{\textbf{\textit{#1}}}}
\usepackage{longtable,booktabs,array}
\usepackage{setspace}
\usepackage[backend=biber,style=numeric,sorting=none]{biblatex}
\addbibresource{references.bib}
\newcounter{none} % for unnumbered tables
\usepackage{calc} % for calculating minipage widths
% Correct order of tables after \paragraph or \subparagraph
\usepackage{etoolbox}
\makeatletter
\patchcmd\longtable{\par}{\if@noskipsec\mbox{}\fi\par}{}{}
\makeatother
% Allow footnotes in longtable head/foot
\IfFileExists{footnotehyper.sty}{\usepackage{footnotehyper}}{\usepackage{footnote}}
\makesavenoteenv{longtable}
\setlength{\emergencystretch}{3em} % prevent overfull lines
\providecommand{\tightlist}{%
  \setlength{\itemsep}{0pt}\setlength{\parskip}{0pt}}
\usepackage{hyperref}
\usepackage{bookmark}
\IfFileExists{xurl.sty}{\usepackage{xurl}}{} % add URL line breaks if available
\urlstyle{same}
\hypersetup{
  pdftitle={A Unified Framework for Cross-Layer Security Analysis of Deep-Learning Software Ecosystems},
  pdfauthor={Tony Arslan},
  hidelinks,
  pdfcreator={LaTeX via pandoc}}

\title{A Unified Framework for Cross-Layer Security Analysis of Deep-Learning Software Ecosystems}
\author{Tony Arslan\\University of Nebraska--Lincoln\\Advisor: Dr.~Witawas Srisa-an}
\date{November 2026}

\begin{document}
\maketitle
\doublespacing

\pagenumbering{roman}
\chapter*{Research Status}
\addcontentsline{toc}{chapter}{Research Status}
\begin{quote}
\textbf{Working manuscript, not a completed-results dissertation.}
The April 2026 preliminary pipeline is the only completed end-to-end research run described here; the later protocol dry run validates file and checksum procedures, not UMSAF. Model inspection, runtime tracing, targeted reverse engineering, cross-layer integration, and final evaluation remain planned. Interpret proposed methods and expected contributions accordingly.
\end{quote}
\chapter*{Abstract}
\addcontentsline{toc}{chapter}{Abstract}

Deep-learning applications depend on software ecosystems extending far
beyond high-level application source code, including serialized model
artifacts, computational graphs, third-party dependencies, native
binaries, runtime environments, hardware-acceleration layers, and
complex software supply chains. Frameworks and libraries such as
TensorFlow, Keras, and NumPy rely extensively on native C/C++ components
beneath high-level Python abstractions. Consequently, vulnerabilities
and security-relevant behaviors originating within native binaries,
serialized models, dependency chains, and runtime layers may remain
undetected when analysis is limited to individual parts of the stack.

Prior research has identified limitations in existing security analysis
techniques for ML/DL software ecosystems and risks involving vulnerable
dependencies, silent computational failures, malicious or manipulated
model artifacts, hidden execution paths, software supply-chain
weaknesses, and runtime vulnerabilities. This
dissertation\textquotesingle s empirical focus is a controlled
deep-learning workload; broader ML/DL ecosystems provide research
context, not a claim that the proposed framework has been validated
across them.

This research proposes a \textbf{Unified Multilayer Security Analysis
Framework (UMSAF)} for cross-layer security analysis of deep-learning
software ecosystems. It integrates software composition analysis,
Software Bill of Materials (SBOM) generation, vulnerability
intelligence, binary reverse engineering, serialized-model inspection,
runtime behavior analysis, cross-layer interaction modeling, and risk
assessment. UMSAF is designed to correlate evidence from high-level
Python application code through dependencies, model artifacts, native
components, and runtime execution. Its methods may inform broader ML/DL
ecosystem analysis, but the primary evaluation is bounded to the
selected deep-learning case study.

The framework will use interpretable rule-based risk scoring and assess
whether defensible labeled data supports a separate
\textbf{machine-learning-based risk-prioritization method}, using
features from dependency graphs, vulnerability metadata, binary
structures, API interactions, and runtime behaviors. This possible
analysis technique is distinct from the deep-learning workload being
analyzed. Its evaluation is conditional on data adequacy and does not
replace the rule-based and severity-only comparison. Binary
reverse-engineering tools such as Ghidra will be used to inspect
selected native components associated with TensorFlow, Keras, NumPy, and
related dependencies.

The framework will be developed and evaluated using a \textbf{design
science and empirical evaluation methodology} centered on a controlled
\textbf{deep-learning image-classification application implemented as a
Convolutional Neural Network (CNN) with TensorFlow/Keras and CIFAR-10}.
This workload provides a reproducible case for examining dependencies,
native components, serialized model artifacts, runtime behavior, and
cross-layer interactions. Additional workloads may be included when
needed to evaluate a specific capability, but one CNN case cannot
establish general effectiveness across ML/DL systems.

The objective is to develop and empirically evaluate a \textbf{scalable,
automated, and holistic cross-layer security analysis framework} for
deep-learning software ecosystems, with conclusions bounded by the
workloads and evidence actually evaluated.

\textbf{Keywords:} Deep Learning, Deep-Learning Software Security,
Machine-Learning-Based Risk Prioritization, AI Software Assurance,
Software Composition Analysis, Reverse Engineering, Binary Analysis,
TensorFlow, Keras, NumPy, Software Bill of Materials (SBOM), Software
Supply Chain Security, Cross-Layer Security Analysis, Vulnerability
Detection, Serialized Model Security, Runtime Analysis

\clearpage
\tableofcontents
\clearpage
\pagenumbering{arabic}

\chapter{Introduction}\label{1-introduction}

The rapid adoption of \textbf{machine learning (ML), particularly deep
learning (DL),} has transformed modern software systems across domains
such as healthcare, finance, energy, and cybersecurity. This
dissertation focuses empirically on a Python-based deep-learning
application and its supporting ecosystem. Such applications depend on
frameworks and libraries including TensorFlow, Keras, and NumPy, where
high-level Python interfaces interact with native C/C++ components,
serialized model artifacts, and runtime environments. The broader ML/DL
context motivates the work, but is not itself the scope of empirical
validation.

Although Python provides memory-safety characteristics that reduce
certain classes of low-level vulnerabilities, the underlying native
C/C++ components do not provide the same protections. These components
may introduce security risks such as buffer overflows, memory
corruption, improper memory management, unsafe system interactions, and
vulnerabilities inherited from outdated or compromised dependencies. In
addition, the growing complexity of ML/DL software supply chains
introduces risks involving vulnerable third-party packages, dependency
confusion, malicious package insertion, compromised build artifacts, and
tampering with binary or serialized model artifacts.

Existing software security approaches commonly focus on individual
aspects of the software stack, such as source code analysis, dependency
scanning, vulnerability identification, or runtime monitoring. Such
approaches may provide limited visibility into third-party libraries,
native binaries, serialized models, and interactions that occur across
software layers. This limitation is particularly significant in ML/DL
systems, where substantial computational functionality may be delegated
from high-level Python APIs to underlying native C/C++ components and
runtime execution environments. When these layers are analyzed
independently, vulnerabilities, hidden execution paths, malicious model
behaviors, and security-relevant cross-layer interactions may remain
unidentified or insufficiently characterized.

Unlike approaches that focus primarily on isolated vulnerability
detection, dependency analysis, or binary inspection, this research
positions \textbf{ML/DL security as a holistic software ecosystem
assurance problem} involving application code, third-party dependencies,
serialized model artifacts, native binaries, runtime environments,
software supply chains, and interactions among these layers. To address
this problem, this research proposes the \textbf{Unified Multilayer
Security Analysis Framework (UMSAF)}, which integrates software
composition analysis, Software Bill of Materials (SBOM) generation,
vulnerability intelligence, binary reverse engineering, serialized model
inspection, runtime behavior analysis, cross-layer interaction modeling,
and integrated risk assessment. By correlating security evidence across
multiple layers of the ML/DL software stack, UMSAF aims to provide more
comprehensive visibility into how vulnerabilities and security-relevant
behaviors originate, propagate, interact, and manifest throughout the
software ecosystem.

\chapter{Problem Statement}\label{2-problem-statement}

Modern \textbf{machine learning (ML) and deep learning (DL) software
systems} operate across multiple interconnected layers, including
high-level Python application code, third-party dependencies, serialized
model artifacts, computational frameworks, native C/C++ components,
runtime environments, and hardware acceleration layers. Existing
security analysis approaches commonly examine these components
independently and therefore may not provide comprehensive visibility
into vulnerabilities, security-relevant behaviors, and interactions
across the complete ML/DL software stack.

Specifically:

\begin{itemize}
\tightlist
\item
  Source-level analysis tools may provide limited visibility into
  vulnerabilities contained within compiled native C/C++ components.
\item
  Software composition analysis and dependency scanning tools can
  identify known vulnerable dependencies but generally do not examine
  the internal behavior of native binaries or their runtime execution.
\item
  Binary analysis and reverse engineering tools can inspect compiled
  components but are typically not integrated with dependency
  intelligence, serialized model analysis, or application-level context.
\item
  Serialized model artifacts and computational graphs may introduce
  security risks that are not adequately represented by conventional
  source code, dependency, or binary analysis alone.
\item
  Runtime analysis can reveal execution behaviors that are not
  observable through static analysis but may lack correlation with
  findings from application, dependency, model, and binary layers.
\item
  Cross-layer interactions and vulnerability propagation paths are not
  consistently modeled, correlated, and evaluated as part of a unified
  security analysis process.
\end{itemize}

This fragmentation creates \textbf{security visibility gaps} in which
vulnerabilities and security-relevant behaviors may remain undetected,
insufficiently characterized, or disconnected from their broader
execution context. A vulnerability originating in one layer may
propagate through dependencies, native components, serialized models, or
runtime execution and manifest in another layer, making isolated
analysis insufficient for understanding the security posture of the
complete system.

Therefore, there is a need for a \textbf{unified, cross-layer security
analysis framework} capable of examining ML/DL applications as
interconnected software ecosystems rather than collections of
independently analyzed components. Such a framework should correlate
security evidence across application code, third-party dependencies,
serialized model artifacts, native binaries, and runtime environments to
provide end-to-end visibility into vulnerabilities, cross-layer
interactions, vulnerability propagation, and software supply-chain
risks.

\chapter{Research Objectives}\label{3-research-objectives}

The primary objective of this research is to design, implement, and
empirically evaluate the \textbf{Unified Multilayer Security Analysis
Framework (UMSAF)} for cross-layer security analysis, with its primary
empirical evaluation conducted on a controlled TensorFlow/Keras
deep-learning software ecosystem.

The specific research objectives are to:

\begin{enumerate}
\def\labelenumi{\arabic{enumi}.}
\item
  \textbf{Identify and characterize complete dependency structures} of
  Python-based ML/DL applications, including direct, transitive, and
  native dependencies.
\item
  \textbf{Identify and correlate known vulnerabilities} using
  established vulnerability intelligence sources, including CVE, OSV,
  and NVD data.
\item
  \textbf{Generate and analyze Software Bills of Materials (SBOMs)} to
  improve visibility into software components, dependencies, versions,
  and associated supply-chain risks within ML/DL software ecosystems.
\item
  \textbf{Extract and analyze native C/C++ binary components} associated
  with ML/DL frameworks and their dependencies to identify
  security-relevant characteristics that may not be visible at the
  Python application layer.
\item
  \textbf{Apply binary reverse engineering and static analysis
  techniques} to inspect compiled libraries, native interfaces, symbols,
  control structures, and potentially vulnerable execution paths.
\item
  \textbf{Inspect serialized model artifacts and computational
  structures} to identify security-relevant characteristics, potentially
  unsafe behaviors, and relationships between model artifacts and
  underlying framework components.
\item
  \textbf{Analyze runtime behavior and execution characteristics} to
  identify security-relevant activities and behaviors that may not be
  observable through static analysis alone.
\item
  \textbf{Model and correlate cross-layer interactions} among Python
  application code, third-party dependencies, serialized model
  artifacts, native C/C++ components, and runtime execution
  environments.
\item
  \textbf{Develop automated risk classification and vulnerability
  prioritization mechanisms} using features derived from dependency
  intelligence, vulnerability metadata, binary structures, API
  interactions, model artifacts, and runtime behaviors.
\item
  \textbf{Empirically evaluate the effectiveness of UMSAF} in providing
  integrated cross-layer security visibility and identifying security
  risks that may be missed or insufficiently characterized by isolated
  analysis approaches.
\end{enumerate}

\chapter{Scope of Research}\label{4-scope-of-research}

This research focuses empirically on a Python-based
\textbf{deep-learning software ecosystem} using frameworks and
applications that rely on native C/C++ computational components. The
broader ML/DL ecosystem is the conceptual context for UMSAF; this study
does not claim that its single primary workload validates general
effectiveness across all ML/DL systems. Within the selected case, the
research investigates how vulnerabilities, security-relevant behaviors,
and attack surfaces originate, propagate, interact, and manifest across:

\begin{itemize}
\tightlist
\item
  application code
\item
  third-party libraries and dependencies
\item
  machine learning and deep learning frameworks
\item
  serialized model artifacts and computational graphs
\item
  native C/C++ binaries
\item
  runtime environments
\item
  hardware acceleration and supporting infrastructure components
\end{itemize}

The primary case study encompasses a deep-learning software ecosystem
utilizing:

\begin{itemize}
\tightlist
\item
  TensorFlow
\item
  Keras
\item
  NumPy
\item
  and related Python and native dependencies that support ML/DL
  execution.
\end{itemize}

The experimental analysis emphasizes software environments in which
high-level Python APIs interact with lower-level native C/C++
implementations through compiled libraries, framework execution engines,
runtime components, and hardware acceleration interfaces.

The proposed research specifically investigates in that case study:

\begin{itemize}
\tightlist
\item
  software composition and dependency analysis
\item
  Software Bill of Materials (SBOM) generation and analysis
\item
  software supply-chain security
\item
  vulnerability identification and correlation
\item
  native binary extraction and analysis
\item
  reverse engineering of native ML/DL components
\item
  serialized model and computational graph inspection
\item
  runtime behavior and execution analysis
\item
  cross-layer interaction modeling
\item
  vulnerability propagation across software layers
\item
  and integrated risk classification and vulnerability prioritization.
\end{itemize}

The research emphasizes correlation of security evidence across
interconnected layers of the selected deep-learning software ecosystem
rather than analyzing individual software components in isolation. The
framework is designed with broader ML/DL applicability in mind, but
conclusions will be limited to the cases and evidence actually
evaluated.

The dissertation does not attempt to address all aspects of ML/DL
security. In particular, the following topics are outside the primary
scope of this research:

\begin{itemize}
\tightlist
\item
  adversarial robustness optimization
\item
  model accuracy or performance improvement
\item
  federated learning security
\item
  cryptographic privacy-preserving machine learning
\item
  side-channel attacks against hardware accelerators
\item
  and formal verification of neural network correctness.
\end{itemize}

Although adversarial machine learning literature is considered where
relevant to the broader security context, the primary emphasis of this
research is \textbf{software ecosystem security, cross-layer security
analysis, and multilayer software assurance}, rather than adversarial
perturbation generation or defense techniques.

The experimental implementation is limited to \textbf{controlled
research environments}, centered on a \textbf{deep-learning
image-classification application implemented as a CNN with
TensorFlow/Keras using CIFAR-10}, to develop and empirically evaluate
UMSAF. Additional workloads or framework components may be incorporated
when necessary to test a specific capability. The research does not
claim exhaustive coverage or general effectiveness across ML/DL
frameworks, application types, hardware platforms, or deployment
environments.

Despite these defined boundaries, the proposed research aims to
establish a \textbf{scalable and extensible foundation for holistic,
cross-layer security assurance} within modern ML/DL software ecosystems.

\chapter{Research Questions}\label{5-research-questions}

The following research questions guide the design, implementation, and
empirical evaluation of the proposed Unified Multilayer Security
Analysis Framework (UMSAF):

\subsubsection{RQ1 --- Dependency and Supply-Chain
Analysis}\label{rq1--dependency-and-supply-chain-analysis}

How can dependencies, native binary components, and associated software
supply-chain risks within the selected Python-based deep-learning
software ecosystem be systematically identified, analyzed, and
correlated?

\subsubsection{RQ2 --- Effectiveness of Existing Security Analysis
Approaches}\label{rq2--effectiveness-of-existing-security-analysis-approaches}

How effective are existing software security analysis techniques when
applied to the selected deep-learning application, its framework,
dependencies, and native components, and what security visibility gaps
remain when these techniques are applied independently?

\subsubsection{RQ3 --- Native Binary Analysis and Reverse
Engineering}\label{rq3--native-binary-analysis-and-reverse-engineering}

To what extent can binary analysis and reverse engineering techniques
recover security-relevant structural and behavioral information from
native C/C++ components underlying the selected Python-based
deep-learning framework and application?

\subsubsection{RQ4 --- Serialized Model and Computational Graph
Security}\label{rq4--serialized-model-and-computational-graph-security}

How can the selected deep-learning model artifact and computational
graph be systematically inspected to identify security-relevant
characteristics, potentially malicious functionality, unsafe operations,
and interactions with underlying framework components?

\subsubsection{RQ5 --- Runtime Security
Analysis}\label{rq5--runtime-security-analysis}

What security-relevant behaviors and execution characteristics can be
identified through runtime analysis that may not be observable through
source-level, dependency, serialized model, or binary analysis alone?

\subsubsection{RQ6 --- Cross-Layer Security
Analysis}\label{rq6--cross-layer-security-analysis}

To what extent can cross-layer interaction modeling and correlation of
security evidence improve vulnerability identification,
characterization, and prioritization across application code,
dependencies, serialized model artifacts, native binaries, and runtime
environments?

\subsubsection{RQ7 --- Integrated Framework
Effectiveness}\label{rq7--integrated-framework-effectiveness}

To what extent does UMSAF provide more comprehensive security visibility
and risk identification in the evaluated deep-learning software
ecosystem than selected isolated security-analysis approaches?

\chapter{Threat Model for Deep-Learning Software Ecosystems}\label{6-threat-model-for-deep-learning-software-ecosystems}

\section{Overview}\label{61-overview}

The threat model draws on the broader \textbf{machine learning (ML) and
deep learning (DL) software ecosystem} literature. The
dissertation\textquotesingle s primary empirical focus is a
deep-learning workload and its supporting ecosystem; the seven-layer
model is a conceptual analysis framework, not evidence that every layer
or threat has been validated across all ML/DL systems. Such ecosystems
may include application code, third-party dependencies, ML/DL
frameworks, serialized model artifacts and computational graphs, native
C/C++ components, runtime environments, hardware acceleration, and
operating-system or infrastructure layers.

ML/DL systems frequently rely on:

\begin{itemize}
\tightlist
\item
  dynamically loaded native libraries
\item
  serialized model artifacts and computational graphs
\item
  framework execution engines
\item
  hardware-specific execution paths
\item
  third-party package repositories
\item
  external software dependencies
\item
  and externally distributed pretrained models.
\end{itemize}

As a result, vulnerabilities and security-relevant behaviors may
originate within one or more software layers and propagate across
abstraction and trust boundaries through dependency relationships, model
execution, native library invocation, and runtime interactions.

This research adopts a \textbf{holistic, cross-layer threat-modeling
approach} to characterize threats affecting ML/DL software ecosystems
and to guide the design and empirical evaluation of the proposed
\textbf{Unified Multilayer Security Analysis Framework (UMSAF)}.

Accordingly, ML/DL software security is treated as an
\textbf{ecosystem-level software assurance problem} rather than a
collection of isolated, independently analyzed security concerns.

\section{ML/DL Ecosystem Layer Model}\label{62-mldl-ecosystem-layer-model}

ML/DL software ecosystems operate across multiple interconnected
software and runtime layers. Modern ML/DL environments rely extensively
on interactions among:

\begin{itemize}
\tightlist
\item
  high-level application code
\item
  third-party libraries and dependencies
\item
  ML/DL frameworks
\item
  serialized model artifacts and computational graphs
\item
  native C/C++ binary components
\item
  runtime and hardware acceleration environments
\item
  and operating system and infrastructure components.
\end{itemize}

These multilayer interactions introduce distinct attack surfaces, trust
boundaries, dependency relationships, and vulnerability propagation
paths. Vulnerabilities or security-relevant behaviors originating within
one layer may affect other layers through dependency relationships,
serialized model execution, framework interactions, native library
invocation, or runtime execution.

This research therefore adopts a \textbf{multilayer ecosystem
perspective} for threat modeling and security analysis rather than
treating ML/DL systems as isolated software applications or
independently analyzed components.

A more detailed conceptual layer model and associated vulnerability
propagation analysis are presented in Section 6.10.

\section{Threat Actors}\label{63-threat-actors}

The threat model considers several categories of adversaries and threat
sources that may affect ML/DL software ecosystems.

\subsection{Malicious Package Maintainers}\label{631-malicious-package-maintainers}

Attackers may publish or distribute:

\begin{itemize}
\tightlist
\item
  compromised Python packages
\item
  malicious wheel distributions
\item
  tampered dependency updates
\item
  or malicious native components embedded within software packages.
\end{itemize}

These attacks may target:

\begin{itemize}
\tightlist
\item
  package repositories such as PyPI
\item
  dependency resolution mechanisms
\item
  direct and transitive dependency chains
\item
  and package distribution mechanisms.
\end{itemize}

\subsection{Software Supply-Chain Attackers}\label{632-software-supply-chain-attackers}

Adversaries may compromise:

\begin{itemize}
\tightlist
\item
  dependency repositories
\item
  CI/CD pipelines
\item
  build environments
\item
  package distribution mechanisms
\item
  model-sharing platforms
\item
  or software artifacts.
\end{itemize}

Potential impacts include:

\begin{itemize}
\tightlist
\item
  insertion of malicious binaries
\item
  dependency poisoning
\item
  artifact tampering
\item
  compromised build outputs
\item
  and hidden backdoors.
\end{itemize}

\subsection{Malicious Model Providers}\label{633-malicious-model-providers}

Attackers may distribute:

\begin{itemize}
\tightlist
\item
  malicious or tampered serialized model artifacts
\item
  poisoned pretrained models
\item
  manipulated computational graphs
\item
  or models containing hidden or unexpected execution behavior.
\end{itemize}

Such artifacts may attempt to invoke unauthorized framework or runtime
functionality, access system resources, or trigger unexpected execution
paths when loaded or executed.

\subsection{Adversarial Users}\label{634-adversarial-users}

Attackers may craft malicious or unexpected inputs intended to:

\begin{itemize}
\tightlist
\item
  trigger vulnerable native operations
\item
  exploit memory-safety weaknesses
\item
  manipulate runtime behavior
\item
  cause denial-of-service conditions
\item
  or influence inference behavior.
\end{itemize}

\subsection{Insider Threats}\label{635-insider-threats}

Internal developers, administrators, or other trusted users may
intentionally or unintentionally:

\begin{itemize}
\tightlist
\item
  introduce vulnerable or compromised dependencies
\item
  weaken or disable security controls
\item
  modify software or model artifacts
\item
  deploy untrusted models
\item
  or alter runtime and infrastructure configurations.
\end{itemize}

\section{Threat Surfaces}\label{64-threat-surfaces}

The proposed threat model identifies the following primary attack
surfaces.

\subsection{Python Application and Dependency Ecosystem}\label{641-python-application-and-dependency-ecosystem}

ML/DL applications depend extensively on frameworks and libraries such
as TensorFlow, Keras, NumPy, and related direct and transitive
dependencies.

Potential threats include:

\begin{itemize}
\tightlist
\item
  dependency confusion
\item
  typosquatting
\item
  compromised packages
\item
  vulnerable transitive dependencies
\item
  malicious package updates
\item
  and dependency tampering.
\end{itemize}

\subsection{Native C/C++ Binary Components}\label{642-native-cc-binary-components}

ML/DL frameworks rely extensively on:

\begin{itemize}
\tightlist
\item
  compiled C/C++ libraries
\item
  dynamically linked libraries
\item
  native extensions
\item
  numerical computation components
\item
  GPU kernels
\item
  and hardware acceleration modules.
\end{itemize}

Potential threats include:

\begin{itemize}
\tightlist
\item
  buffer overflows
\item
  integer overflows
\item
  use-after-free vulnerabilities
\item
  memory corruption
\item
  unsafe memory operations
\item
  and vulnerabilities inherited from native dependencies.
\end{itemize}

\subsection{Serialized Models and Computational Graphs}\label{643-serialized-models-and-computational-graphs}

Modern ML/DL frameworks use serialized model artifacts and computational
graph representations to support model persistence, portability,
deployment, interoperability, and runtime execution.

These artifacts may contain:

\begin{itemize}
\tightlist
\item
  model architectures
\item
  computational graphs
\item
  tensor operations
\item
  operators
\item
  metadata
\item
  weights
\item
  execution descriptors
\item
  signatures
\item
  and runtime configuration information.
\end{itemize}

Because serialized model artifacts may influence or participate directly
in runtime execution, they represent an important security boundary
between model data, framework execution, native components, and the
underlying runtime environment.

\subsection{Serialized Model Security Implications}\label{644-serialized-model-security-implications}

Serialized ML/DL artifacts may introduce security risks involving:

\begin{itemize}
\tightlist
\item
  malicious or manipulated graph execution
\item
  hidden or unexpected operators
\item
  unauthorized runtime invocation
\item
  model tampering
\item
  poisoned pretrained models
\item
  unsafe deserialization
\item
  abuse of framework functionality
\item
  and hidden execution paths.
\end{itemize}

Externally sourced model artifacts therefore represent a significant
trust boundary. Effective security analysis requires examination not
only of application code and software dependencies but also of model
structures, operators, metadata relationships, runtime interactions, and
execution paths associated with serialized artifacts.

For these reasons, UMSAF incorporates \textbf{serialized model
inspection and computational graph analysis} as core components of its
multilayer security analysis methodology.

\subsection{Runtime Execution Environment}\label{645-runtime-execution-environment}

Security-relevant behavior may emerge during:

\begin{itemize}
\tightlist
\item
  model loading
\item
  training
\item
  inference
\item
  native library execution
\item
  GPU or accelerator execution
\item
  and distributed computation.
\end{itemize}

Potential threats and security concerns include:

\begin{itemize}
\tightlist
\item
  runtime manipulation
\item
  execution hijacking
\item
  unexpected native invocation
\item
  resource abuse
\item
  denial-of-service conditions
\item
  and silent computational failures.
\end{itemize}

\subsection{Operating System and Infrastructure Layer}\label{646-operating-system-and-infrastructure-layer}

Underlying operating system and infrastructure components may introduce
risks involving:

\begin{itemize}
\tightlist
\item
  vulnerable system libraries
\item
  vulnerable or compromised drivers
\item
  insecure container configurations
\item
  container escape vulnerabilities
\item
  cloud or infrastructure misconfigurations
\item
  excessive privileges
\item
  and compromised runtime environments.
\end{itemize}

\section{Threat Categories}\label{65-threat-categories}

The threat model categorizes security threats into the following primary
classes.

\subsection{Software Supply-Chain Threats}\label{651-software-supply-chain-threats}

Examples include:

\begin{itemize}
\tightlist
\item
  compromised dependencies
\item
  malicious package artifacts
\item
  dependency poisoning
\item
  compromised build artifacts
\item
  and vulnerable transitive components.
\end{itemize}

\subsection{Native Memory-Safety and Binary Vulnerabilities}\label{652-native-memory-safety-and-binary-vulnerabilities}

Examples include:

\begin{itemize}
\tightlist
\item
  heap overflows
\item
  stack corruption
\item
  use-after-free conditions
\item
  integer overflows
\item
  unsafe memory operations
\item
  and other native-code vulnerabilities.
\end{itemize}

\subsection{Malicious or Compromised Model Artifacts}\label{653-malicious-or-compromised-model-artifacts}

Examples include:

\begin{itemize}
\tightlist
\item
  poisoned models
\item
  malicious or tampered serialized models
\item
  manipulated computational graphs
\item
  hidden execution behavior
\item
  and unauthorized framework or runtime invocation.
\end{itemize}

\subsection{Silent Computational and Runtime Failures}\label{654-silent-computational-and-runtime-failures}

Examples include:

\begin{itemize}
\tightlist
\item
  incorrect inference results
\item
  hidden training corruption
\item
  numerical instability
\item
  unexpected runtime behavior
\item
  and silent framework or native-component failures.
\end{itemize}

\subsection{Adversarial ML/DL Threats}\label{655-adversarial-mldl-threats}

Examples include:

\begin{itemize}
\tightlist
\item
  adversarial inputs
\item
  model extraction
\item
  inference manipulation
\item
  and evasion attacks.
\end{itemize}

Although these threats are relevant to the broader ML/DL security
landscape, adversarial robustness and adversarial perturbation defense
are not primary research areas of this dissertation.

\subsection{Runtime and Infrastructure Threats}\label{656-runtime-and-infrastructure-threats}

Examples include:

\begin{itemize}
\tightlist
\item
  runtime exploitation
\item
  GPU or accelerator runtime vulnerabilities
\item
  insecure distributed execution
\item
  container or operating system compromise
\item
  and cloud infrastructure compromise.
\end{itemize}

\section{Trust Boundaries}\label{66-trust-boundaries}

The framework identifies several critical trust boundaries:

{\def\LTcaptype{none} % do not increment counter
\begin{longtable}[]{@{}ll@{}}
\toprule\noalign{}
\textbf{Boundary} & \textbf{Description} \\
\midrule\noalign{}
\endhead
\bottomrule\noalign{}
\endlastfoot
Application ↔ External Dependencies & Trust placed in third-party and
transitive software components \\
Python/Application ↔ Native Binary & Transition from high-level
application execution to memory-unsafe native C/C++ components \\
Serialized Model ↔ Framework/Runtime & Loading and execution of imported
or externally obtained model artifacts \\
Framework ↔ Native Components & Transition from framework abstractions
to compiled native implementations \\
Framework/Runtime ↔ Hardware Layer & Interaction with GPU, accelerator,
driver, and kernel components \\
Runtime ↔ Operating System/Infrastructure & Interaction between ML/DL
execution environments and underlying system resources \\
Local Environment ↔ External/Cloud Infrastructure & Trust boundary
associated with distributed or externally hosted execution
environments \\
\end{longtable}
}

These boundaries represent locations where:

\begin{itemize}
\tightlist
\item
  trust assumptions may fail
\item
  control or privilege transitions may occur
\item
  untrusted artifacts may enter the system
\item
  and vulnerabilities or security-relevant behaviors may propagate
  across software layers.
\end{itemize}

\section{Threat and Vulnerability Propagation Across Layers}\label{67-threat-and-vulnerability-propagation-across-layers}

A central premise of this research is that security risks within ML/DL
software ecosystems may propagate across multiple software layers rather
than remaining isolated within the component in which they originate.

An example propagation path is:

\textbf{Malicious Python Package} ↓ \textbf{Compromised Native Binary} ↓
\textbf{Unsafe Runtime Execution} ↓ \textbf{Operating System /
Infrastructure Impact}

Another example is:

\textbf{Malicious or Tampered Serialized Model} ↓ \textbf{Framework /
Computational Graph Execution} ↓ \textbf{Native or Runtime Component
Invocation} ↓ \textbf{Unauthorized System Interaction}

These propagation paths illustrate why isolated security analysis may
provide incomplete visibility into the security implications of
vulnerabilities and malicious behaviors that cross software-layer
boundaries.

\section{Security Objectives}\label{68-security-objectives}

The proposed framework seeks to improve the following security
properties:

{\def\LTcaptype{none} % do not increment counter
\begin{longtable}[]{@{}ll@{}}
\toprule\noalign{}
\textbf{Objective} & \textbf{Description} \\
\midrule\noalign{}
\endhead
\bottomrule\noalign{}
\endlastfoot
Integrity & Identify tampered, compromised, or vulnerable dependencies,
binaries, and model artifacts \\
Confidentiality & Identify execution paths or behaviors that may expose
protected information or system resources \\
Availability & Identify vulnerabilities and behaviors capable of
disrupting ML/DL services or runtime execution \\
Transparency & Improve visibility into dependencies, model execution,
native components, and runtime behavior \\
Traceability & Correlate security findings and vulnerability propagation
paths across software layers \\
Trustworthiness & Improve confidence in the security posture of ML/DL
software ecosystems \\
\end{longtable}
}

\section{Threat Model Implications for UMSAF}\label{69-threat-model-implications-for-umsaf}

The threat model directly informs the architecture and analytical
capabilities of the proposed UMSAF.

Specifically, UMSAF will:

\begin{itemize}
\tightlist
\item
  identify and analyze dependency chains
\item
  generate and analyze SBOM information
\item
  correlate known vulnerability intelligence
\item
  analyze package and artifact integrity
\item
  extract and inspect native binary components
\item
  apply reverse engineering techniques to native ML/DL components
\item
  inspect serialized model artifacts and computational graphs
\item
  analyze runtime behavior
\item
  model and correlate cross-layer interactions
\item
  identify vulnerability and threat propagation paths
\item
  and support integrated risk classification and vulnerability
  prioritization.
\end{itemize}

These capabilities are intended to provide comprehensive cross-layer
security visibility into ML/DL software ecosystems and identify security
relationships that may not be apparent when individual layers are
analyzed independently.

\section{ML/DL Software Ecosystem Layer Model}\label{610-mldl-software-ecosystem-layer-model}

The proposed research models ML/DL software ecosystems as multilayer
architectures composed of interconnected application, dependency,
framework, serialization, native binary, runtime, hardware, and
infrastructure components. This conceptual model provides the foundation
for analyzing how vulnerabilities and security-relevant behaviors
originate, propagate, interact, and manifest across multiple abstraction
layers.

The conceptual architecture consists of the following layers:

\textbf{Application Layer} ↓ \textbf{Dependency / Software Supply-Chain
Layer} ↓ \textbf{ML/DL Framework Layer} ↓ \textbf{Serialization /
Computational Graph Layer} ↓ \textbf{Native Binary Layer} ↓
\textbf{Runtime / Hardware Acceleration Layer} ↓ \textbf{Operating
System / Infrastructure Layer}

Each layer introduces distinct attack surfaces, trust boundaries,
dependency relationships, and potential vulnerability propagation paths.

\subsection{Application Layer}\label{application-layer}

The \textbf{Application Layer} contains developer-written Python code,
external interfaces, training scripts, inference logic, data-processing
logic, and application functionality. Although this layer is highly
visible to developers and conventional source-analysis tools,
substantial ML/DL execution may occur within lower layers.

\subsection{Dependency / Software Supply-Chain
Layer}\label{dependency--software-supply-chain-layer}

The \textbf{Dependency / Software Supply-Chain Layer} includes direct
and transitive Python packages, native dependencies, package
repositories, distribution artifacts, and associated dependency
metadata. This layer introduces security concerns involving vulnerable
components, dependency confusion, package tampering, malicious updates,
compromised software artifacts, and other software supply-chain risks.

\subsection{ML/DL Framework Layer}\label{mldl-framework-layer}

The \textbf{ML/DL Framework Layer} includes frameworks and libraries
such as TensorFlow, Keras, NumPy, and related components. These
technologies provide high-level abstractions for tensor computation,
model construction, training, inference, optimization, and execution
while frequently relying on underlying native implementations.

\subsection{Serialization / Computational Graph
Layer}\label{serialization--computational-graph-layer}

The \textbf{Serialization / Computational Graph Layer} includes
serialized model artifacts, computational graphs, operators, metadata,
execution descriptors, and related model representations. These
artifacts may influence runtime execution and introduce attack surfaces
involving model tampering, unsafe serialization or deserialization,
hidden operators, and unexpected execution paths.

\subsection{Native Binary Layer}\label{native-binary-layer}

The \textbf{Native Binary Layer} contains compiled C/C++ libraries,
dynamically linked modules, numerical computation engines, tensor
operation implementations, native extensions, and hardware interaction
components. Because this layer contains memory-unsafe native code,
vulnerabilities such as buffer overflows, integer overflows,
use-after-free conditions, and other memory-corruption weaknesses may
arise.

\subsection{Runtime / Hardware Acceleration
Layer}\label{runtime--hardware-acceleration-layer}

The \textbf{Runtime / Hardware Acceleration Layer} includes framework
execution engines, GPU runtimes, accelerator libraries, distributed
execution components, thread schedulers, memory managers, and hardware
optimization mechanisms. Security-relevant behavior at this layer may
include unexpected execution paths, runtime manipulation, resource
abuse, hardware-specific vulnerabilities, and silent computational
failures.

\subsection{Operating System / Infrastructure
Layer}\label{operating-system--infrastructure-layer}

The \textbf{Operating System / Infrastructure Layer} includes operating
system libraries, drivers, container environments, cloud infrastructure,
orchestration platforms, and deployment configurations. Vulnerabilities
or misconfigurations at this layer may affect ML/DL framework execution
and application-level behavior.

A central premise of this research is that vulnerabilities and
security-relevant behaviors within ML/DL ecosystems may cross multiple
layers rather than remaining isolated within a single component. For
example, a compromised Python dependency may introduce a malicious
native binary that subsequently affects runtime execution and underlying
system resources. Similarly, a malicious serialized model artifact may
trigger unexpected framework, native, or runtime operations that are not
directly visible at the application layer.

This layered conceptual model therefore provides the foundation for
UMSAF by enabling systematic analysis and correlation of dependencies,
vulnerability intelligence, serialized model artifacts, native binaries,
runtime behaviors, and cross-layer propagation paths across modern ML/DL
software ecosystems.

\chapter{Research Timeline}\label{7-research-timeline}

This research targets \textbf{August 2027 graduation}, conditional on
completing the comprehensive exam by the applicable seven-month cutoff
and meeting candidacy requirements. The provisional internal stretch
target is \textbf{March 31, 2027} for research/manuscript completion and
two paper submissions, not defense or final deposit. Review 2 revisions
are scheduled through March 21; the advisor-approved dissertation must
reach the Reading Committee by March 22 to allow the required four weeks
before an earliest early-defense window of April 19-22. Official
August-term deadlines are July 29 for the defense and July 30 for final
dissertation actions; the exact exam cutoff, department deadlines,
advisor/committee availability, and venue dates remain to be confirmed.
The working calendar follows
\url{../writing/Dissertation_Completion_Schedule.md}.

The maximum commitment is \textbf{10 hours/week} (one hour each weekday
and five weekend hours). With the schedule\textquotesingle s two
provisional lighter weeks, approximately \textbf{203 hours} are
available through March 31, including reserve. Retaining the full
research scope is high-risk within that capacity; feasibility must be
reassessed at the schedule\textquotesingle s evidence-based checkpoints
rather than assuming success or silently dropping requirements. Paper
targets are submissions, not guaranteed acceptance or publication.

\subsection{Phase 1 --- Research Contract and Baseline
Refresh}\label{phase-1--research-contract-and-baseline-refresh}

\textbf{November 2-22, 2026}

\begin{itemize}
\tightlist
\item
  Confirm retained RQ1-RQ7/H1-H7, claims, metrics, ground truth, and
  experimental protocol
\item
  Select and rebuild the Windows baseline; record a verified environment
  manifest
\item
  Add baseline tests, rerun Steps 1-6, and reconcile April results and
  skipped capabilities
\item
  Confirm deadlines, reviewer availability, and two distinct paper
  contributions/venues
\item
  Reassess remaining effort using actual hours by November 22
\end{itemize}

\subsection{Phase 2 --- Missing Evidence
Producers}\label{phase-2--missing-evidence-producers}

\textbf{November 23-December 27, 2026}

\begin{itemize}
\tightlist
\item
  Implement and test serialized-model inspection and controlled model
  scenarios
\item
  Capture reproducible Windows runtime observations for control/hardened
  workloads
\item
  Perform targeted PE reverse engineering with saved Ghidra evidence
\item
  Map selected Python APIs, framework/native components, and runtime
  evidence with confidence labels
\item
  Write methods and paper sections incrementally; validate all producers
  before integration
\end{itemize}

\subsection{Phase 3 --- Integrated UMSAF
Artifact}\label{phase-3--integrated-umsaf-artifact}

\textbf{December 28, 2026-January 17, 2027}

\begin{itemize}
\tightlist
\item
  Validate common evidence/component schemas and the cross-layer
  evidence graph
\item
  Implement confidence-aware propagation analysis and negative-control
  tests
\item
  Implement interpretable risk scoring, severity-only comparison, and
  sensitivity analysis
\item
  Complete the data-adequacy decision and ML evaluation if justified
\item
  Validate integrated reports, provenance, and new-module regression
  tests
\end{itemize}

\subsection{Phase 4 --- Empirical Evaluation and
Results}\label{phase-4--empirical-evaluation-and-results}

\textbf{January 18-31, 2027}

\begin{itemize}
\tightlist
\item
  Run repeated controlled workloads and verified security cases
\item
  Compare UMSAF with isolated-analysis baselines using equivalent
  evidence/resources
\item
  Calculate predefined metrics and complete
  scalability/cross-configuration analysis
\item
  Evaluate H1-H7 using saved evidence, including negative/inconclusive
  results
\item
  Freeze results and conduct the validity/reproducibility audit by
  January 31
\end{itemize}

\subsection{Phase 5 --- Complete Draft, Archive, and First
Review}\label{phase-5--complete-draft-archive-and-first-review}

\textbf{February 1-28, 2027}

\begin{itemize}
\tightlist
\item
  Assemble the incrementally written dissertation and reproducibility
  archive by \textbf{February 7}
\item
  Review 1: \textbf{February 8-21}; address feedback \textbf{February
  22-28}
\item
  Submit Paper 1 by \textbf{February 19}, after paper-specific review
  and venue checks
\item
  Send the revised complete dissertation for Review 2 by
  \textbf{February 28}
\end{itemize}

\subsection{Phase 6 --- Second Review, Second Paper, and Defense
Readiness}\label{phase-6--second-review-second-paper-and-defense-readiness}

\textbf{March 1-21, 2027}

\begin{itemize}
\tightlist
\item
  Review 2: \textbf{March 1-14}; address feedback \textbf{March 15-21}
\item
  Complete Paper 2 draft by \textbf{March 7} and submit by \textbf{March
  19}
\item
  Obtain required approvals, finalize slides, and rehearse within the
  weekly limit
\end{itemize}

\subsection{Phase 7 --- Reading Committee Review, Defense, Corrections,
and Final
Deposit}\label{phase-7--reading-committee-review-defense-corrections-and-final-deposit}

\textbf{March 22-July 30, 2027}

\begin{itemize}
\tightlist
\item
  By \textbf{March 22}, send the advisor-approved dissertation to the
  Reading Committee; the required four-week review period means an early
  defense cannot occur before \textbf{April 19}.
\item
  March 31 remains the provisional research/manuscript and
  paper-submission stretch target, not a defense or deposit deadline.
\item
  Target an early defense during \textbf{April 19-22 or later}, subject
  to eligibility, committee availability, application/format-review lead
  times, and approval.
\item
  Complete corrections and final deposit after the defense, no later
  than the August-term deadlines of \textbf{July 29} (latest defense)
  and \textbf{July 30} (final dissertation actions); confirm earlier
  department requirements.
\item
  April-July is protected committee-review, defense, deposit, and
  contingency time---not permission to delay required research or expand
  scope.
\end{itemize}

\chapter{Related Work}\label{8-related-work}

Research in \textbf{machine learning (ML) and deep learning (DL)
security} has evolved from an early emphasis on adversarial examples and
model-level attacks toward broader concerns involving software
vulnerabilities, serialized model security, runtime behavior, software
supply chains, native binary components, and lifecycle-oriented security
assurance. This evolution reflects increasing recognition that modern
ML/DL systems operate as complex, multilayer software ecosystems in
which security risks may originate, propagate, interact, and manifest
across application code, dependencies, frameworks, serialized model
artifacts, native components, runtime environments, hardware
acceleration layers, and supporting infrastructure.
\cite{xiao2018security,lai2025security,tan2026securing,meiklejohn2025modelsupplychain}

The literature reviewed in this section establishes the foundation for
the proposed \textbf{Unified Multilayer Security Analysis Framework
(UMSAF)} and identifies limitations in existing security-analysis
approaches when applied independently to individual components or layers
of ML/DL software ecosystems.

\section{Reverse Engineering and Binary Analysis}\label{81-reverse-engineering-and-binary-analysis}

Reverse engineering is a well-established technique for analyzing
compiled software, recovering structural information, investigating
execution behavior, and identifying potential vulnerabilities in binary
components. Tools such as Ghidra, IDA Pro, and Binary Ninja provide
capabilities including disassembly, decompilation, control-flow graph
reconstruction, symbol analysis, and related binary inspection
techniques.

Traditionally, reverse engineering has been applied extensively to
malware analysis, vulnerability research, software recovery, and
standalone binary inspection. Its integration into holistic security
analysis of \textbf{ML/DL software ecosystems}, however, remains
comparatively limited. This gap is significant because modern ML/DL
frameworks and supporting numerical libraries rely extensively on
compiled native components written in C and C++. Security-relevant
operations occurring within these components may not be directly visible
through high-level Python source analysis.
\cite{ning2024binarysca,wang2024binarysca}

This research builds upon established binary-analysis methodologies by
integrating native binary inspection and reverse engineering with
dependency intelligence, vulnerability information, serialized model
inspection, runtime analysis, and cross-layer security correlation.

\section{Static Analysis Limitations in ML/DL Libraries}\label{82-static-analysis-limitations-in-mldl-libraries}

Static analysis tools are widely used to identify software defects and
potential vulnerabilities without executing the target software. Prior
empirical research, however, has identified significant limitations when
these techniques are applied to complex ML/DL frameworks and libraries.

Studies examining real-world defects and vulnerabilities in frameworks
such as TensorFlow, PyTorch, and MXNet report that existing static
analysis techniques detect only a limited subset of known issues.
\cite{harzevili2023static,filus2023software}
Contributing factors include:

\begin{itemize}
\tightlist
\item
  high framework complexity and data-dependent behavior
\item
  substantial reliance on native C/C++ implementations
\item
  runtime-dependent execution behavior
\item
  dynamically constructed computational operations
\item
  hardware- and backend-specific execution paths
\item
  and limited visibility into interactions across software layers.
\end{itemize}

These findings motivate security-analysis approaches that complement
conventional static analysis with dependency analysis, native binary
inspection, serialized model analysis, runtime observation, and
cross-layer correlation.

\section{Characteristics of Vulnerabilities in ML/DL Libraries}\label{83-characteristics-of-vulnerabilities-in-mldl-libraries}

Understanding the characteristics of vulnerabilities within ML/DL
frameworks is important for developing effective security-analysis
methodologies. Prior empirical studies have examined vulnerabilities
across widely used frameworks and numerical libraries, including
TensorFlow, PyTorch, NumPy, SciPy, and related components.
\cite{harzevili2023characterizing,lai2025security,zhou2025unraveling}

Reported vulnerability characteristics include:

\begin{itemize}
\tightlist
\item
  improper input validation
\item
  memory-management errors
\item
  memory corruption
\item
  integer-related weaknesses
\item
  crashes and undefined behavior
\item
  incorrect computational results
\item
  vulnerabilities distributed across multiple framework components
\item
  and defects affecting both API-level and implementation-level code.
\end{itemize}

These findings indicate that vulnerabilities within ML/DL software
ecosystems cannot always be adequately characterized solely from
high-level application code. Instead, effective security analysis may
require examination of framework internals, dependency relationships,
native implementations, and runtime behavior.
\cite{harzevili2023characterizing,lai2025security,zhou2025unraveling}

This observation supports the development of security methodologies
capable of \textbf{classifying and correlating vulnerabilities across
multiple software layers}.

\section{Native Code as a Significant Attack Surface}\label{84-native-code-as-a-significant-attack-surface}

ML/DL frameworks rely extensively on optimized native implementations to
provide efficient tensor computation, numerical processing, model
execution, and hardware acceleration. Consequently, security
vulnerabilities may occur within native C/C++ components that execute
beneath high-level Python APIs.

Prior research has identified native-code vulnerabilities involving:

\begin{itemize}
\tightlist
\item
  memory corruption
\item
  buffer overflows
\item
  integer overflows
\item
  NULL pointer dereferences
\item
  use-after-free conditions
\item
  and improper input validation.
\end{itemize}
\cite{xiao2018security,filus2023software,ning2024binarysca}

These findings demonstrate that native components constitute a
significant attack surface within Python-based ML/DL software
ecosystems. Security analysis limited to high-level Python source code
may therefore provide incomplete visibility into vulnerabilities
residing within compiled framework components and native dependencies.

\section{Testing and Dynamic Analysis}\label{85-testing-and-dynamic-analysis}

Testing approaches such as differential testing have been developed to
identify defects in ML/DL libraries. EAGLE constructs equivalent
computation graphs and compares their outputs within a single library,
reporting detected bugs in TensorFlow and PyTorch.
\cite{wang2022eagle}
ML/DL library testing also presents challenges involving:

\begin{itemize}
\tightlist
\item
  generation of semantically valid input combinations
\item
  large and complex API surfaces
\item
  incomplete exploration of internal execution paths
\item
  computational cost
\item
  framework-specific execution constraints
\item
  and limited integration with dependency, model, and binary-level
  security information.
\end{itemize}

These limitations suggest that library testing provides useful but
bounded evidence and should be considered complementary to other forms
of structural, runtime, and cross-layer analysis.
\cite{wang2022eagle}

\section{Multilayer Architecture of ML/DL Software Ecosystems}\label{86-multilayer-architecture-of-mldl-software-ecosystems}

Modern ML/DL software ecosystems operate across multiple interconnected
abstraction, software, and execution layers. Consistent with the
conceptual layer model defined in Section 6.10, this dissertation
represents the ML/DL software ecosystem using the following layers.
The cited studies inform the dependency and supply-chain portions of this
model; the seven-layer integration remains the subject of this
dissertation's evaluation.\cite{zhou2025unraveling,tan2026securing,baird2026attackchains}

\textbf{Application Layer} ↓ \textbf{Dependency / Software Supply-Chain
Layer} ↓ \textbf{ML/DL Framework Layer} ↓ \textbf{Serialization /
Computational Graph Layer} ↓ \textbf{Native Binary Layer} ↓
\textbf{Runtime / Hardware Acceleration Layer} ↓ \textbf{Operating
System / Infrastructure Layer}

Recent TensorFlow supply-chain and SBOM-graph studies describe
vulnerabilities and dependency relationships spanning software
components.\cite{zhou2025unraveling,tan2026securing,baird2026attackchains}
The multilayer architecture introduces further security-analysis
challenges:

\begin{itemize}
\tightlist
\item
  dependencies and trust relationships span multiple abstraction levels
\item
  high-level application APIs may obscure underlying framework and
  native execution
\item
  third-party and transitive dependencies may introduce vulnerabilities
  or compromised components
\item
  serialized model artifacts and computational graphs may influence
  runtime execution
\item
  native C/C++ components may introduce memory-safety and
  implementation-level vulnerabilities
\item
  execution behavior may vary across runtime, hardware, and
  infrastructure environments
\item
  and security findings originating within one layer may require
  evidence from other layers for accurate interpretation and risk
  assessment.
\end{itemize}

Existing security-analysis techniques frequently specialize in
individual portions of this software stack. Source analysis, software
composition analysis, model inspection, binary analysis, runtime
monitoring, and infrastructure-security tools therefore provide
different but potentially complementary forms of security evidence.

This fragmentation motivates a cross-layer security-analysis approach
capable of correlating evidence across the complete ML/DL software
ecosystem rather than evaluating each layer independently.

\section{Fault Injection and Runtime Behavior Analysis}\label{87-fault-injection-and-runtime-behavior-analysis}

Fault injection and runtime-analysis techniques have been used to
investigate the resilience and execution behavior of ML/DL systems under
abnormal or intentionally modified conditions. Such studies demonstrate
that runtime behavior may differ from expectations derived solely from
static program structure.

Reported observations include:

\begin{itemize}
\tightlist
\item
  runtime behavior that differs from static expectations
\item
  lower-layer faults affecting application-level outputs
\item
  silent or difficult-to-observe computational failures
\item
  hardware- or backend-dependent behavior
\item
  and execution paths that are difficult to characterize using source
  analysis alone.
\end{itemize}

These observations motivate runtime tracing, behavioral analysis, and
cross-layer correlation as complementary techniques for understanding
security-relevant behavior within ML/DL software ecosystems.
\cite{jia2021symptoms,tambon2023silent}

\section{Security Implications in Real-World Applications}\label{88-security-implications-in-real-world-applications}

ML/DL systems are increasingly deployed in domains such as:

\begin{itemize}
\tightlist
\item
  autonomous and cyber-physical systems
\item
  healthcare
\item
  financial services
\item
  energy
\item
  cybersecurity
\item
  and other operational environments.
\end{itemize}

As ML/DL technologies become embedded within software systems that
influence consequential decisions and operations, vulnerabilities within
underlying frameworks, dependencies, models, and runtime environments
may have broader security implications.

The widespread adoption of ML/DL therefore increases the importance of
security-analysis methodologies capable of examining not only
model-level attacks but also the underlying software ecosystem
supporting model development and execution.
\cite{lai2025security,tan2026securing,meiklejohn2025modelsupplychain}

\section{Gap Analysis and Research Positioning}\label{89-gap-analysis-and-research-positioning}

Based on the reviewed literature, several research gaps motivate this
dissertation and the development of the \textbf{Unified Multilayer
Security Analysis Framework (UMSAF)}.

\subsection{Limitations of Isolated Security Analysis}\label{1-limitations-of-isolated-security-analysis}

Existing security-analysis techniques provide valuable capabilities but
generally address specific portions of the ML/DL software ecosystem:

\begin{itemize}
\tightlist
\item
  source and static analysis examine application code or intermediate
  representations
\item
  software composition analysis identifies dependencies and known
  vulnerability information
\item
  serialized model analysis examines model artifacts and computational
  structures
\item
  binary analysis and reverse engineering inspect compiled native
  components
\item
  dynamic and runtime analysis observe execution behavior
\item
  infrastructure-security tools evaluate operating system, container,
  cloud, and deployment environments.
\end{itemize}

Individually, these approaches may provide incomplete visibility into
security relationships spanning multiple software layers.

The reviewed collection also contains limited peer-reviewed evaluation
of SBOM-generation completeness and vulnerability-correlation accuracy
specifically for deep-learning software ecosystems. CycloneDX and SPDX
define SBOM formats, but format conformance alone is not evidence of
inventory completeness or security-analysis effectiveness.
\cite{cyclonedxSpecification,spdxSpecification}

\subsection{Limited Cross-Layer Correlation}\label{2-limited-cross-layer-correlation}

Recent studies analyze TensorFlow vulnerability propagation through
software supply chains and multi-vulnerability chains using SBOM graphs.
These are important adjacent contributions, but do not by themselves
establish evidence correlation across all seven layers in the UMSAF
model.\cite{zhou2025unraveling,baird2026attackchains,tan2026securing}
This dissertation will evaluate that broader integration rather than
assume its benefit. The UMSAF ecosystem is represented as:

\textbf{Application Layer} ↓ \textbf{Dependency / Software Supply-Chain
Layer} ↓ \textbf{ML/DL Framework Layer} ↓ \textbf{Serialization /
Computational Graph Layer} ↓ \textbf{Native Binary Layer} ↓
\textbf{Runtime / Hardware Acceleration Layer} ↓ \textbf{Operating
System / Infrastructure Layer}

A significant research opportunity therefore exists in systematically
correlating security findings across these abstraction, dependency,
execution, and trust boundaries.

\subsection{Limited Native Binary Visibility}\label{3-limited-native-binary-visibility}

Python-based ML/DL applications may depend on substantial compiled
native functionality that is not directly represented in high-level
application source code. Security analysis operating primarily at the
application or package-metadata level may therefore provide limited
visibility into:

\begin{itemize}
\tightlist
\item
  embedded native components
\item
  native dependency relationships
\item
  low-level implementation behavior
\item
  memory-safety weaknesses
\item
  API-to-native execution relationships
\item
  and vulnerabilities within compiled framework components.
\end{itemize}

\subsection{Limited Integration of Serialized Model and Runtime Security Analysis}\label{4-limited-integration-of-serialized-model-and-runtime-security-analysis}

Serialized model artifacts, computational graphs, framework execution
engines, and runtime environments may introduce security-relevant
behaviors that cannot be fully characterized through application source,
dependency metadata, or native binary analysis alone.

A research gap therefore exists in integrating:

\begin{itemize}
\tightlist
\item
  serialized model inspection
\item
  computational graph analysis
\item
  native component analysis
\item
  runtime behavior analysis
\item
  and application-level context
\end{itemize}

within a unified security-analysis process.

\subsection{Fragmentation Across Security-Analysis Techniques}\label{5-fragmentation-across-security-analysis-techniques}

Relevant security-analysis capabilities---including:

\begin{itemize}
\tightlist
\item
  software composition and dependency analysis
\item
  SBOM generation and analysis
\item
  vulnerability intelligence
\item
  serialized model inspection
\item
  binary analysis
\item
  reverse engineering
\item
  runtime behavior analysis
\item
  cross-layer interaction modeling
\item
  and risk classification and prioritization
\end{itemize}

are frequently applied independently rather than as components of an
integrated cross-layer security-assurance methodology.

These gaps collectively motivate the development of a unified framework
capable of correlating security evidence across multiple layers of
modern ML/DL software ecosystems.

\section{Positioning of This Research}\label{810-positioning-of-this-research}

This dissertation addresses the identified research gaps by proposing
the \textbf{Unified Multilayer Security Analysis Framework (UMSAF)}, a
cross-layer security-analysis methodology designed to correlate security
evidence across the interconnected layers of ML/DL software ecosystems.
It builds on recent supply-chain propagation and SBOM-graph studies while
evaluating a broader evidence model spanning the seven specified layers.
\cite{zhou2025unraveling,baird2026attackchains}

Consistent with the threat model and conceptual architecture defined in
Section 6, UMSAF examines security across the following layers:

\textbf{Application Layer} ↓ \textbf{Dependency / Software Supply-Chain
Layer} ↓ \textbf{ML/DL Framework Layer} ↓ \textbf{Serialization /
Computational Graph Layer} ↓ \textbf{Native Binary Layer} ↓
\textbf{Runtime / Hardware Acceleration Layer} ↓ \textbf{Operating
System / Infrastructure Layer}

To analyze security across these layers, UMSAF integrates:

\begin{itemize}
\tightlist
\item
  software composition and dependency analysis
\item
  Software Bill of Materials (SBOM) generation and analysis
\item
  vulnerability intelligence and correlation
\item
  package and artifact integrity analysis
\item
  native binary extraction and analysis
\item
  binary reverse engineering
\item
  serialized model and computational graph inspection
\item
  runtime behavior and execution analysis
\item
  cross-layer interaction and vulnerability propagation modeling
\item
  and risk prioritization, using rule-based scoring as the required
  approach and machine-learning-based prioritization only if data
  adequacy supports it.
\end{itemize}

UMSAF is distinguished by its emphasis on \textbf{correlating security
evidence across software layers rather than analyzing individual
components independently}. The framework is designed to establish
relationships among application code, dependencies, framework
components, serialized model artifacts, native binaries, runtime
behavior, and underlying execution environments.

The proposed framework therefore extends beyond conventional source
analysis, dependency scanning, binary inspection, or runtime monitoring
by integrating these forms of security evidence into a unified
cross-layer analysis process.

The proposed research will empirically evaluate whether this integrated
approach provides more comprehensive security visibility, vulnerability
characterization, cross-layer traceability, and risk identification than
isolated security-analysis approaches.

Through this integration, UMSAF aims to provide a scalable and
extensible foundation for holistic security assurance within modern
ML/DL software ecosystems.

\section{Evolution of ML/DL Software Security}\label{811-evolution-of-mldl-software-security}

Early ML/DL security research focused substantially on adversarial
attacks against trained models. Initial studies demonstrated that
carefully crafted input perturbations could manipulate neural-network
predictions, establishing adversarial robustness as an important area of
ML/DL security research.

Subsequent research expanded to include:

\begin{itemize}
\tightlist
\item
  data poisoning
\item
  model extraction
\item
  model inversion
\item
  privacy leakage
\item
  malicious model manipulation
\item
  and related attacks affecting the ML/DL lifecycle.
\end{itemize}

As ML/DL frameworks became increasingly integrated into production
software environments, security research expanded beyond model
robustness toward vulnerabilities within the underlying software
infrastructure. Studies have identified security concerns involving
frameworks and libraries such as TensorFlow, PyTorch, NumPy, and related
components, including memory-safety vulnerabilities, runtime failures,
serialization risks, and software supply-chain weaknesses.

More recent work has further broadened ML/DL security toward lifecycle
and software-ecosystem assurance involving:

\begin{itemize}
\tightlist
\item
  dependency management
\item
  software supply-chain security
\item
  model distribution
\item
  serialized model security
\item
  native binary components
\item
  runtime behavior
\item
  and cross-layer security analysis.
\end{itemize}

This progression demonstrates the increasing importance of integrating
traditional software security, software engineering, supply-chain
analysis, and runtime assurance with established ML/DL security
research.\cite{xiao2018security,meiklejohn2025modelsupplychain,lai2025security,tan2026securing}

\section{Security Vulnerabilities in ML/DL Frameworks}\label{812-security-vulnerabilities-in-mldl-frameworks}

Frameworks and libraries such as TensorFlow, PyTorch, NumPy, and related
computational components contain substantial native codebases
implemented in C/C++ and other low-level languages for performance
optimization, numerical computation, and hardware acceleration.

Prior research has identified vulnerabilities involving:

\begin{itemize}
\tightlist
\item
  buffer overflows
\item
  integer overflows
\item
  use-after-free conditions
\item
  NULL pointer dereferences
\item
  memory corruption
\item
  and improper input validation.
\end{itemize}

Many such vulnerabilities occur within native components responsible for
tensor operations, serialization, graph execution, numerical
computation, and hardware acceleration. Additional risks may arise
through transitive dependencies and embedded third-party components.

Because these components execute beneath high-level Python abstractions,
application developers and source-level analysis tools may have limited
visibility into their implementation and associated vulnerabilities.
These findings motivate security-analysis approaches capable of
examining both high-level ML/DL applications and their underlying native
software components.\cite{xiao2018security,filus2023software,lai2025security,zhou2025unraveling}

\section{Bugs and Silent Failures in Deep Learning Systems}\label{813-bugs-and-silent-failures-in-deep-learning-systems}

Deep learning systems may exhibit failure characteristics that differ
from conventional deterministic software. In addition to explicit
crashes and runtime exceptions, DL frameworks may experience silent
computational failures in which incorrect or unexpected results occur
without immediately observable failure indicators.

Prior research has reported issues involving:

\begin{itemize}
\tightlist
\item
  incorrect tensor computations
\item
  silent training corruption
\item
  numerical instability
\item
  inconsistent inference behavior
\item
  and invalid gradient calculations.
\end{itemize}
\cite{jia2021symptoms,tambon2023silent}

Within the reviewed collection, recent peer-reviewed studies focused
specifically on security-oriented runtime tracing of deep-learning
framework execution were limited. This remains a literature gap for the
proposed work, not evidence that runtime monitoring is generally absent
from the field.

Such failures can be difficult to identify because ML/DL outputs may be
probabilistic, data-dependent, and affected by runtime or hardware
characteristics.

Runtime behavior may also vary across:

\begin{itemize}
\tightlist
\item
  hardware configurations
\item
  compiler optimizations
\item
  execution backends
\item
  framework versions
\item
  and distributed processing environments.
\end{itemize}

These observations motivate the use of runtime analysis, behavioral
tracing, anomaly detection, and cross-layer monitoring as complementary
components of ML/DL software assurance.
\cite{jia2021symptoms,tambon2023silent,wang2022eagle}

\section{Reverse Engineering of ML/DL Systems}\label{814-reverse-engineering-of-mldl-systems}

Reverse engineering methodologies provide mechanisms for examining
compiled software structures and execution behavior that may not be
visible at the source-code level. Relevant techniques include:

\begin{itemize}
\tightlist
\item
  disassembly
\item
  decompilation
\item
  control-flow reconstruction
\item
  symbol analysis
\item
  static binary inspection
\item
  and dynamic tracing.
\end{itemize}

Binary software-composition studies demonstrate techniques for
identifying components in compiled programs. Depending on the target
artifact, binary analysis may reveal information concerning:

\begin{itemize}
\tightlist
\item
  native operator implementations
\item
  computational structures
\item
  framework internals
\item
  API-to-native relationships
\item
  execution paths
\item
  and other implementation characteristics.
\end{itemize}

Because modern ML/DL frameworks rely extensively on compiled native
libraries, binary analysis and reverse engineering provide a potentially
valuable complement to dependency, source-level, model, and runtime
analysis.

These findings support incorporating reverse engineering into a broader
multilayer security-analysis methodology such as UMSAF, while treating
binary results as one evidence source rather than proof of runtime
behavior or exploitability.\cite{ning2024binarysca,wang2024binarysca}

\section{ML/DL Threat Assessment and Lifecycle Security}\label{815-mldl-threat-assessment-and-lifecycle-security}

ML/DL systems introduce security concerns across multiple stages of the
software and model lifecycle, including:

\begin{itemize}
\tightlist
\item
  data acquisition and preparation
\item
  model training
\item
  model serialization
\item
  software packaging and distribution
\item
  deployment
\item
  inference
\item
  and runtime execution.
\end{itemize}

Security risks may arise across the same interconnected software layers
defined by the UMSAF conceptual architecture:

\begin{itemize}
\tightlist
\item
  Application Layer
\item
  Dependency / Software Supply-Chain Layer
\item
  ML/DL Framework Layer
\item
  Serialization / Computational Graph Layer
\item
  Native Binary Layer
\item
  Runtime / Hardware Acceleration Layer
\item
  Operating System / Infrastructure Layer.
\end{itemize}

Training-phase attacks may affect model integrity before deployment,
while compromised dependencies, malicious software artifacts, or
tampered serialized models may introduce security risks during packaging
and deployment. Runtime environments introduce additional concerns
involving native execution, hardware acceleration, distributed
processing, and infrastructure interaction.
\cite{kellas2025pickleball,digregorio2026insecurity,zhu2025mymodel,gan2025sentry}

Software supply-chain security is particularly relevant because ML/DL
applications may depend on numerous third-party packages, native
libraries, pretrained models, and external repositories. These
dependencies increase the number of trust relationships that must be
considered when evaluating the security of the complete software
ecosystem.\cite{tan2026securing,zhou2025unraveling,meiklejohn2025modelsupplychain}

These observations reinforce the need for lifecycle-oriented and
cross-layer security methodologies capable of analyzing ML/DL software
ecosystems holistically.

\section{Testing and Validation of Deep Learning Libraries}\label{816-testing-and-validation-of-deep-learning-libraries}

Testing DL systems presents challenges associated with:

\begin{itemize}
\tightlist
\item
  nondeterministic execution
\item
  numerical instability
\item
  hardware acceleration
\item
  complex framework APIs
\item
  and runtime-dependent behavior.
\end{itemize}

Traditional software-testing approaches remain valuable but may not
identify all failure modes associated with ML/DL frameworks. Researchers
have therefore investigated specialized testing approaches, including:

\begin{itemize}
\tightlist
\item
  fuzz testing
\item
  mutation testing
\item
  differential testing
\item
  metamorphic testing
\item
  runtime validation
\item
  and behavioral monitoring.
\end{itemize}

Mutation testing evaluates whether test mechanisms detect intentionally
introduced changes or faults. Differential testing compares behavior
across implementations, versions, or execution environments, while
metamorphic testing evaluates expected relationships among transformed
inputs and outputs.

These techniques provide useful mechanisms for evaluating ML/DL
framework behavior and may be incorporated as supporting experimental
techniques where appropriate. However, this dissertation does not
position mutation testing or any individual testing technique as the
primary research contribution. Instead, such techniques may support the
\textbf{empirical evaluation and validation of UMSAF}.

Collectively, the reviewed literature indicates that effective ML/DL
software security assurance may require the integration of:

\begin{itemize}
\tightlist
\item
  software composition and dependency analysis
\item
  static and binary analysis
\item
  serialized model inspection
\item
  dynamic and runtime analysis
\item
  behavioral monitoring
\item
  vulnerability intelligence
\item
  and cross-layer security correlation.
\end{itemize}

The absence of a unified methodology integrating these forms of security
evidence motivates the development and empirical evaluation of the
\textbf{Unified Multilayer Security Analysis Framework (UMSAF)}.

\chapter{Preliminary Study}\label{9-preliminary-study}

\section{Overview}\label{91-overview}

\textbf{Evidence boundary:} The recorded April 20, 2026 run contains
dependency analysis, pip-audit/OSV scanning, CycloneDX generation,
RECORD checks/PyPI expected-digest lookups, static findings and a
derived containment policy, and PE extraction/fingerprinting. PyPI
\texttt{verified} indicates expected-digest metadata retrieval, not an
archive-digest comparison or authenticated provenance. The run does not
contain implemented model inspection, runtime traces, Ghidra
reverse-engineering results, API/native execution mappings, integrated
propagation/risk reports, or final baseline comparisons. Sections
9.6-9.9 describe intended extensions and conceptual paths unless a
specific saved artifact is cited. A saved model is not model-inspection
evidence; a DLL import is not proof of runtime execution; a generated
policy is not enforced containment.

To evaluate the feasibility of the proposed \textbf{Unified Multilayer
Security Analysis Framework (UMSAF)}, a preliminary study was conducted
using a \textbf{deep learning--based image classification application
implemented as a Convolutional Neural Network (CNN) with
TensorFlow/Keras}. Although deep learning (DL) is a subset of machine
learning (ML), the experimental workload is more precisely classified as
deep learning because the model uses multiple trainable neural-network
layers, including convolutional, pooling, and fully connected layers, to
learn hierarchical image features directly from pixel data through
training and backpropagation.

The dissertation therefore retains \textbf{machine learning and deep
learning software security} as the broader research domain while
identifying the specific preliminary experimental workload as a
\textbf{deep learning CNN workload}.

The purpose of the preliminary study is to establish the feasibility of
examining security characteristics across the interconnected layers of
an ML/DL software ecosystem. Consistent with the threat model and
conceptual architecture defined in Section 6, the study considers the
following layers:

\textbf{Application Layer} ↓ \textbf{Dependency / Software Supply-Chain
Layer} ↓ \textbf{ML/DL Framework Layer} ↓ \textbf{Serialization /
Computational Graph Layer} ↓ \textbf{Native Binary Layer} ↓
\textbf{Runtime / Hardware Acceleration Layer} ↓ \textbf{Operating
System / Infrastructure Layer}

The preliminary study focuses particularly on relationships among Python
application code, third-party dependencies, TensorFlow/Keras framework
functionality, serialized model artifacts, native C/C++ components, and
runtime execution. These relationships provide a controlled environment
for investigating the feasibility of dependency analysis, Software Bill
of Materials (SBOM) generation, vulnerability analysis, native binary
extraction, reverse engineering, serialized-model inspection, runtime
analysis, and cross-layer security correlation.

The study focuses on widely used ML/DL technologies, particularly
TensorFlow, Keras, and NumPy, which provide a representative software
ecosystem for examining interactions between high-level Python
abstractions and underlying native computational components.

\section{Experimental Setup}\label{92-experimental-setup}

A controlled experimental environment was developed using a simplified
\textbf{deep learning Convolutional Neural Network (CNN)} implemented in
Python with TensorFlow/Keras. The application performs supervised
multiclass image classification using the CIFAR-10 dataset.

The CNN architecture contains two convolutional feature-extraction
blocks followed by fully connected classification layers. This
architecture makes the application a deep learning workload rather than
a conventional non-neural machine learning model.

The purpose of the experimental application is not to optimize
image-classification performance. Instead, it provides a controlled and
reproducible deep-learning workload through which multiple components
and execution layers of the selected software ecosystem can be examined.

The experimental workload supports investigation of:

\begin{itemize}
\tightlist
\item
  application-level Python code
\item
  direct and transitive software dependencies
\item
  software supply-chain relationships
\item
  TensorFlow/Keras framework operations
\item
  native C/C++ components
\item
  serialized model artifacts and computational structures
\item
  runtime execution behavior
\item
  and interactions among multiple software layers.
\end{itemize}

The application uses the \textbf{CIFAR-10 dataset}, which contains
60,000 32 × 32 color images distributed across 10 object categories. The
dataset provides a manageable and widely used image-classification
workload suitable for controlled deep learning experimentation.

The preliminary CNN application performs the following operations:

\begin{itemize}
\tightlist
\item
  dataset loading
\item
  data normalization
\item
  convolutional neural network construction
\item
  model compilation
\item
  supervised training
\item
  model evaluation
\item
  prediction and confidence estimation
\item
  and model serialization.
\end{itemize}

The following experimental application was used as the preliminary
controlled workload.

\textbf{Listing 1 --- Preliminary Deep Learning CNN Experimental
Application}

\begin{Shaded}
\begin{Highlighting}[]
\ImportTok{import}\NormalTok{ numpy }\ImportTok{as}\NormalTok{ np}
\ImportTok{import}\NormalTok{ tensorflow }\ImportTok{as}\NormalTok{ tf}

\NormalTok{keras }\OperatorTok{=}\NormalTok{ tf.keras}
\NormalTok{layers }\OperatorTok{=}\NormalTok{ tf.keras.layers}

\KeywordTok{def}\NormalTok{ main():}
    \CommentTok{\# 1) Load a tiny built{-}in dataset}
    \CommentTok{\#    (CIFAR{-}10: 32x32 color images, 10 classes)}
\NormalTok{    (x\_train, y\_train), (x\_test, y\_test) }\OperatorTok{=}\NormalTok{ keras.datasets.cifar10.load\_data()}

\NormalTok{    class\_names }\OperatorTok{=}\NormalTok{ [}
        \StringTok{"airplane"}\NormalTok{, }\StringTok{"automobile"}\NormalTok{, }\StringTok{"bird"}\NormalTok{, }\StringTok{"cat"}\NormalTok{, }\StringTok{"deer"}\NormalTok{,}
        \StringTok{"dog"}\NormalTok{, }\StringTok{"frog"}\NormalTok{, }\StringTok{"horse"}\NormalTok{, }\StringTok{"ship"}\NormalTok{, }\StringTok{"truck"}
\NormalTok{    ]}

    \CommentTok{\# 2) Normalize pixel values to [0, 1]}
\NormalTok{    x\_train }\OperatorTok{=}\NormalTok{ x\_train.astype(}\StringTok{"float32"}\NormalTok{) }\OperatorTok{/} \FloatTok{255.0}
\NormalTok{    x\_test }\OperatorTok{=}\NormalTok{ x\_test.astype(}\StringTok{"float32"}\NormalTok{) }\OperatorTok{/} \FloatTok{255.0}

    \CommentTok{\# 3) Build a small CNN}
\NormalTok{    model }\OperatorTok{=}\NormalTok{ keras.Sequential([}
\NormalTok{        layers.Input(shape}\OperatorTok{=}\NormalTok{(}\DecValTok{32}\NormalTok{, }\DecValTok{32}\NormalTok{, }\DecValTok{3}\NormalTok{)),}

\NormalTok{        layers.Conv2D(}\DecValTok{32}\NormalTok{, (}\DecValTok{3}\NormalTok{, }\DecValTok{3}\NormalTok{), activation}\OperatorTok{=}\StringTok{"relu"}\NormalTok{),}
\NormalTok{        layers.MaxPooling2D(),}

\NormalTok{        layers.Conv2D(}\DecValTok{64}\NormalTok{, (}\DecValTok{3}\NormalTok{, }\DecValTok{3}\NormalTok{), activation}\OperatorTok{=}\StringTok{"relu"}\NormalTok{),}
\NormalTok{        layers.MaxPooling2D(),}

\NormalTok{        layers.Flatten(),}
\NormalTok{        layers.Dense(}\DecValTok{64}\NormalTok{, activation}\OperatorTok{=}\StringTok{"relu"}\NormalTok{),}
\NormalTok{        layers.Dense(}\DecValTok{10}\NormalTok{, activation}\OperatorTok{=}\StringTok{"softmax"}\NormalTok{),}
\NormalTok{    ])}

\NormalTok{    model.}\BuiltInTok{compile}\NormalTok{(}
\NormalTok{        optimizer}\OperatorTok{=}\StringTok{"adam"}\NormalTok{,}
\NormalTok{        loss}\OperatorTok{=}\StringTok{"sparse\_categorical\_crossentropy"}\NormalTok{,}
\NormalTok{        metrics}\OperatorTok{=}\NormalTok{[}\StringTok{"accuracy"}\NormalTok{]}
\NormalTok{    )}

\NormalTok{    model.summary()}

    \CommentTok{\# 4) Train}
\NormalTok{    history }\OperatorTok{=}\NormalTok{ model.fit(}
\NormalTok{        x\_train, y\_train,}
\NormalTok{        validation\_split}\OperatorTok{=}\FloatTok{0.1}\NormalTok{,}
\NormalTok{        epochs}\OperatorTok{=}\DecValTok{5}\NormalTok{,}
\NormalTok{        batch\_size}\OperatorTok{=}\DecValTok{64}\NormalTok{,}
\NormalTok{        verbose}\OperatorTok{=}\DecValTok{2}
\NormalTok{    )}

    \CommentTok{\# 5) Evaluate}
\NormalTok{    test\_loss, test\_acc }\OperatorTok{=}\NormalTok{ model.evaluate(x\_test, y\_test, verbose}\OperatorTok{=}\DecValTok{0}\NormalTok{)}
    \BuiltInTok{print}\NormalTok{(}\SpecialStringTok{f"}\CharTok{\textbackslash{}n}\SpecialStringTok{Test accuracy: }\SpecialCharTok{\{}\NormalTok{test\_acc}\SpecialCharTok{:.4f\}}\SpecialStringTok{"}\NormalTok{)}

    \CommentTok{\# 6) Predict a few images}
\NormalTok{    idxs }\OperatorTok{=}\NormalTok{ np.random.choice(}\BuiltInTok{len}\NormalTok{(x\_test), size}\OperatorTok{=}\DecValTok{5}\NormalTok{, replace}\OperatorTok{=}\VariableTok{False}\NormalTok{)}
\NormalTok{    preds }\OperatorTok{=}\NormalTok{ model.predict(x\_test[idxs], verbose}\OperatorTok{=}\DecValTok{0}\NormalTok{)}
\NormalTok{    pred\_labels }\OperatorTok{=}\NormalTok{ np.argmax(preds, axis}\OperatorTok{=}\DecValTok{1}\NormalTok{)}

    \BuiltInTok{print}\NormalTok{(}\StringTok{"}\CharTok{\textbackslash{}n}\StringTok{Sample predictions:"}\NormalTok{)}
    \ControlFlowTok{for}\NormalTok{ i, idx }\KeywordTok{in} \BuiltInTok{enumerate}\NormalTok{(idxs):}
\NormalTok{        true\_label }\OperatorTok{=} \BuiltInTok{int}\NormalTok{(y\_test[idx][}\DecValTok{0}\NormalTok{])}
\NormalTok{        pred\_label }\OperatorTok{=} \BuiltInTok{int}\NormalTok{(pred\_labels[i])}
\NormalTok{        confidence }\OperatorTok{=} \BuiltInTok{float}\NormalTok{(np.}\BuiltInTok{max}\NormalTok{(preds[i]))}

        \BuiltInTok{print}\NormalTok{(}
            \SpecialStringTok{f"Image \#}\SpecialCharTok{\{}\NormalTok{idx}\SpecialCharTok{:5d\}}\SpecialStringTok{ | "}
            \SpecialStringTok{f"True: }\SpecialCharTok{\{}\NormalTok{class\_names[true\_label]}\SpecialCharTok{:10s\}}\SpecialStringTok{ | "}
            \SpecialStringTok{f"Pred: }\SpecialCharTok{\{}\NormalTok{class\_names[pred\_label]}\SpecialCharTok{:10s\}}\SpecialStringTok{ | "}
            \SpecialStringTok{f"Conf: }\SpecialCharTok{\{}\NormalTok{confidence}\SpecialCharTok{:.2f\}}\SpecialStringTok{"}
\NormalTok{        )}

    \CommentTok{\# 7) Save the model}
\NormalTok{    model.save(}\StringTok{"basic\_cifar10\_cnn.keras"}\NormalTok{)}
    \BuiltInTok{print}\NormalTok{(}\StringTok{"}\CharTok{\textbackslash{}n}\StringTok{Saved model to basic\_cifar10\_cnn.keras"}\NormalTok{)}

\ControlFlowTok{if} \VariableTok{\_\_name\_\_} \OperatorTok{==} \StringTok{"\_\_main\_\_"}\NormalTok{:}
\NormalTok{    main()}
\end{Highlighting}
\end{Shaded}

The deep learning CNN application serves as the entry point for the
multilayer analysis. Execution of the application involves interactions
among Python code, TensorFlow/Keras APIs, dependency components,
serialized model artifacts, native framework components, runtime
execution engines, and supporting system resources.

The intentionally simplified workload therefore provides a practical
foundation for investigating security visibility across multiple ML/DL
software layers without introducing unnecessary application complexity.

\subsection{Experimental Environment}\label{921-experimental-environment}

The preliminary experiments utilize a controlled deep learning analysis
environment designed to support dependency analysis, vulnerability
scanning, binary extraction, reverse engineering, serialized-model
inspection, runtime analysis, and cross-layer security investigation.

The primary experimental environment consists of a \textbf{Windows
x86\_64 system using CPython 3.12.10}. Windows is selected as the
primary platform because it provides the controlled and reproducible
environment currently used to implement and execute the UMSAF prototype.
The platform choice defines the operating-system-specific scope of the
empirical results; the broader UMSAF methodology remains applicable to
other operating systems through platform-specific analysis adapters.

The primary ML/DL software stack includes:

\begin{itemize}
\tightlist
\item
  Python 3.12
\item
  TensorFlow 2.21.0
\item
  Keras 3.x
\item
  NumPy 2.x
\item
  and supporting scientific-computing and framework dependencies.
\end{itemize}

Additional experimental configurations may include PyTorch and ONNX
environments where comparative analysis or framework interoperability
testing is necessary. These additional environments are not required for
the preliminary CNN workload and may be incorporated selectively during
the broader empirical evaluation of UMSAF.

Reverse engineering and binary-analysis capabilities include:

\begin{itemize}
\tightlist
\item
  Ghidra 11.x
\item
  objdump
\item
  \texttt{pefile}
\item
  strings
\item
  Windows dependency-inspection utilities
\item
  and related Windows binary-analysis tools.
\end{itemize}

The primary native-binary analysis targets Windows PE-format
\texttt{.pyd} and \texttt{.dll} components. Linux ELF analysis is
outside the primary experimental configuration and is not used as
evidence for the results reported in this dissertation. A future Linux
configuration may be evaluated as comparative work, but it is not
required for the present UMSAF evaluation.

Software composition, dependency, SBOM, and vulnerability-analysis
capabilities include:

\begin{itemize}
\tightlist
\item
  Python package and dependency inspection
\item
  CycloneDX SBOM tooling
\item
  pip-audit
\item
  OSV vulnerability intelligence
\item
  CVE/NVD vulnerability information
\item
  and related software composition analysis techniques.
\end{itemize}

Isolated Python virtual environments and containerized execution
environments may be used to improve experimental reproducibility and
reduce interference among dependency and framework configurations.

Hardware acceleration may also be incorporated into selected experiments
using NVIDIA GPU and CUDA runtime environments when investigation of
hardware-accelerated execution paths is required. Such analysis may
examine:

\begin{itemize}
\tightlist
\item
  GPU-related execution paths
\item
  hardware-accelerated tensor operations
\item
  runtime dependency loading
\item
  framework-to-runtime interactions
\item
  and interactions with hardware acceleration components.
\end{itemize}

This experimental environment provides a reproducible and extensible
foundation for developing and evaluating UMSAF across controlled ML/DL
software configurations.

\section{Dependency, Software Supply-Chain, and SBOM Analysis}\label{93-dependency-software-supply-chain-and-sbom-analysis}

The first phase of the preliminary study examines the software
dependency structure associated with the CNN application and its
supporting ML/DL framework environment.

Although the application itself contains relatively little Python code,
its execution depends on TensorFlow, Keras, NumPy, and numerous direct
and transitive dependencies. Some of these dependencies contain or
interact with compiled native components that are not immediately
apparent from inspection of the application source code.

The dependency-analysis process includes:

\begin{enumerate}
\def\labelenumi{\arabic{enumi}.}
\tightlist
\item
  identifying installed Python packages and versions
\item
  enumerating direct and transitive dependencies
\item
  identifying package-distribution artifacts
\item
  identifying packages containing native binary components
\item
  generating an SBOM
\item
  and correlating dependency information with available vulnerability
  intelligence.
\end{enumerate}

The SBOM captures information such as:

\begin{itemize}
\tightlist
\item
  package names
\item
  package versions
\item
  dependency relationships
\item
  package-distribution artifacts
\item
  and associated software-component metadata.
\end{itemize}

The preliminary run generated CycloneDX output; SPDX generation was
skipped. These format specifications define interoperable SBOM
structures, but do not by themselves establish inventory completeness
or vulnerability-detection accuracy.\cite{cyclonedxSpecification,spdxSpecification}

Preliminary analysis demonstrates that even a comparatively small CNN
application depends on a substantially larger supporting software
ecosystem than is visible from the application source code alone.

This observation supports the inclusion of a distinct \textbf{Dependency
/ Software Supply-Chain Layer} within the UMSAF architecture and
demonstrates the importance of examining ML/DL applications as software
ecosystems rather than isolated Python programs.

\section{Vulnerability Intelligence and Correlation}\label{94-vulnerability-intelligence-and-correlation}

Dependency information obtained during software composition analysis can
be correlated with established vulnerability-intelligence sources to
identify known vulnerabilities associated with components within the
experimental environment.

Relevant vulnerability sources include:

\begin{itemize}
\tightlist
\item
  Common Vulnerabilities and Exposures (CVE)
\item
  Open Source Vulnerabilities (OSV)
\item
  National Vulnerability Database (NVD)
\item
  and vulnerability information exposed through software composition
  analysis tools.
\end{itemize}

The purpose of this analysis is not merely to enumerate known
vulnerabilities. Instead, UMSAF seeks to associate vulnerability
information with the software components and layers in which affected
packages or binaries participate.

The recorded April 2026 pip-audit/OSV snapshot reported zero findings.
It does not demonstrate a detected transitive vulnerability or absence
of vulnerabilities. The planned evaluation will use verified
advisory-linked cases to determine how vulnerability information relates
to transitive components; NVD enrichment and other omitted checks
require explicit implementation/decision records.

The full empirical evaluation will extend this process by correlating
vulnerability intelligence with dependency relationships, binary
components, model artifacts, and runtime evidence where appropriate.

\section{Native Binary Extraction and Analysis}\label{95-native-binary-extraction-and-analysis}

A central component of the preliminary study involves identifying native
binary components associated with the Python-based ML/DL environment.

The binary-analysis process includes:

\begin{enumerate}
\def\labelenumi{\arabic{enumi}.}
\tightlist
\item
  obtaining relevant TensorFlow, NumPy, and supporting package artifacts
\item
  inspecting package contents
\item
  identifying compiled native components
\item
  identifying dynamically linked dependencies
\item
  extracting selected binaries for further inspection
\item
  and mapping native components to their corresponding packages and
  framework functions where feasible.
\end{enumerate}

Relevant binary formats may include:

\begin{itemize}
\tightlist
\item
  \texttt{.pyd} Python extension modules on Windows
\item
  \texttt{.dll} dynamic libraries on Windows
\item
  \texttt{.so} shared objects on Linux for future comparative
  configurations
\item
  and other compiled framework or dependency components.
\end{itemize}

Preliminary inspection identifies substantial compiled functionality
associated with the ML/DL environment. TensorFlow contains native
framework components, including internal wrapper functionality, while
NumPy contains compiled extensions supporting numerical operations.

Native components may also interact with external numerical, runtime,
system, or hardware-specific libraries.

These observations demonstrate that a Python-based deep learning
application cannot be characterized solely through analysis of its
Python source code. Significant computational functionality is delegated
through framework interfaces to compiled native components.

This finding supports the inclusion of a dedicated \textbf{Native Binary
Layer} within UMSAF and motivates deeper binary inspection and reverse
engineering.

\section{Reverse Engineering Analysis}\label{96-reverse-engineering-analysis}

Selected native binary components can be subjected to reverse
engineering to investigate implementation characteristics that are not
visible through high-level Python application analysis.

The planned targeted reverse engineering process will examine
characteristics such as:

\begin{itemize}
\tightlist
\item
  binary structure
\item
  imported and exported functions
\item
  linked libraries
\item
  strings
\item
  symbols
\item
  function structures
\item
  control-flow characteristics
\item
  native interfaces
\item
  and low-level computational functionality.
\end{itemize}

Ghidra provides a primary reverse engineering environment for examining
selected native components, supplemented by \texttt{pefile}, strings,
Windows dependency-inspection utilities, and other platform-appropriate
binary-analysis tools. The primary analysis concerns PE-format binaries;
ELF-oriented tools such as \texttt{readelf} may be used only in future
comparative Linux experiments.

The objective is not to reverse engineer the complete TensorFlow
codebase. Such an undertaking would be unnecessarily broad for the
purpose of this dissertation. Instead, reverse engineering will be used
selectively to investigate native components and execution paths
relevant to dependency, vulnerability, model, or runtime evidence. The
existing PE fingerprints/imports/exports are not Ghidra
function/control-flow evidence.

Preliminary inspection demonstrates the feasibility of extracting and
examining native components underlying the high-level ML/DL framework
environment.

The broader UMSAF evaluation will use reverse engineering to investigate
security-relevant native components, relationships among framework APIs
and native implementations, and selected vulnerable or potentially
security-relevant execution paths.

\section{Serialized Model and Computational Graph Analysis}\label{97-serialized-model-and-computational-graph-analysis}

The experimental CNN explicitly produces a serialized model artifact
through the following operation:

\texttt{model.save("basic\_cifar10\_cnn.keras")}

This artifact introduces an additional analysis surface beyond the
application source code, dependencies, and native binaries.

Serialized model and computational structure analysis considers
characteristics such as:

\begin{itemize}
\tightlist
\item
  model architecture
\item
  layer definitions
\item
  model metadata
\item
  serialization format
\item
  model configuration
\item
  computational operations
\item
  framework relationships
\item
  and information influencing model loading and execution.
\end{itemize}

The purpose of this analysis is to determine what security-relevant
information can be extracted from serialized ML/DL artifacts and how
such information can be correlated with the framework, dependency,
native binary, and runtime layers.

The preliminary CNN demonstrates that model serialization is part of the
normal application lifecycle and therefore represents an important
component of the software ecosystem rather than merely an output file.

The full UMSAF evaluation will extend this analysis to investigate
security-relevant characteristics of serialized artifacts, including
potentially unexpected structures, unsafe behaviors, manipulated
artifacts, and relationships between serialized model operations and
underlying framework execution.

This analysis corresponds to the \textbf{Serialization / Computational
Graph Layer} defined in the UMSAF conceptual architecture.

\section{Runtime Behavior Analysis}\label{98-runtime-behavior-analysis}

Static inspection of application code, dependencies, serialized
artifacts, and native binaries provides important structural information
but does not necessarily reveal how these components interact during
execution.

Runtime analysis therefore provides an additional source of security
evidence.

The experimental CNN creates multiple runtime events during:

\begin{itemize}
\tightlist
\item
  dataset loading
\item
  model construction
\item
  model training
\item
  model evaluation
\item
  inference
\item
  model serialization
\item
  native library invocation
\item
  memory allocation
\item
  and, where applicable, hardware-accelerated execution.
\end{itemize}

Runtime analysis may be used to observe:

\begin{itemize}
\tightlist
\item
  dynamically loaded libraries
\item
  process and thread behavior
\item
  framework-to-native execution transitions
\item
  file-system interactions
\item
  resource utilization
\item
  native component invocation
\item
  and hardware/runtime interactions.
\end{itemize}

The purpose of runtime analysis within UMSAF is to identify execution
characteristics that may not be observable through static analysis alone
and to correlate those observations with evidence obtained from other
software layers.

The preliminary experimental environment establishes the foundation for
such analysis. More comprehensive runtime tracing and behavioral
analysis will be conducted during the empirical evaluation of UMSAF.

This capability corresponds primarily to the \textbf{Runtime / Hardware
Acceleration Layer}, while also providing evidence about interactions
with the Native Binary and Operating System / Infrastructure Layers.

\section{Cross-Layer Interaction Analysis}\label{99-cross-layer-interaction-analysis}

The central purpose of UMSAF is not merely to execute multiple
security-analysis techniques independently but to correlate the
resulting evidence across software layers.

The CNN workload provides a controlled environment for investigating
these interactions.

Consistent with the UMSAF architecture defined in Section 6, the
complete conceptual layer model is:

\textbf{Application Layer} ↓ \textbf{Dependency / Software Supply-Chain
Layer} ↓ \textbf{ML/DL Framework Layer} ↓ \textbf{Serialization /
Computational Graph Layer} ↓ \textbf{Native Binary Layer} ↓
\textbf{Runtime / Hardware Acceleration Layer} ↓ \textbf{Operating
System / Infrastructure Layer}

Within this architecture, a specific observed or investigated execution
path may traverse a subset of these layers. For example:

\textbf{Python CNN Application} ↓ \textbf{TensorFlow/Keras API} ↓
\textbf{Framework / Wrapper Functionality} ↓ \textbf{Native C/C++
Component} ↓ \textbf{Runtime / System Library / Hardware}

Similarly, serialization introduces another interaction path:

\textbf{Python CNN Application} ↓ \textbf{TensorFlow/Keras Framework} ↓
\textbf{Serialized Model Artifact} ↓ \textbf{Model Loading / Framework
Execution} ↓ \textbf{Native and Runtime Components}

These paths should not be interpreted as alternative UMSAF
architectures. Instead, they represent specific interaction and
execution paths occurring within the broader seven-layer model.

The preliminary dependency and PE inventories provide inputs for future
mapping. The conceptual paths above are not validated API-to-native
execution mappings; Task 15 must produce source-linked evidence and
confidence labels before such mappings can be reported as observed.

The full UMSAF evaluation will extend this capability by correlating
evidence from dependency analysis, vulnerability intelligence, binary
inspection, serialized model analysis, reverse engineering, and runtime
observation to investigate how security-relevant characteristics and
vulnerabilities may interact or propagate across software layers.

\section{Preliminary Findings}\label{910-preliminary-findings}

The preliminary study provides several observations that inform the
continued development of UMSAF.

\subsection{Significant Software Ecosystem Complexity}\label{1-significant-software-ecosystem-complexity}

Even a comparatively small deep learning application depends on a
substantially larger ecosystem of frameworks, packages, native
components, and runtime resources than is apparent from its Python
source code.

This finding supports treating ML/DL applications as \textbf{software
ecosystems} rather than isolated application programs.

\subsection{Significant Reliance on Native Components}\label{2-significant-reliance-on-native-components}

The TensorFlow/Keras experimental workload delegates substantial
computational functionality from high-level Python interfaces to
underlying framework and native components.

This creates a security-analysis requirement extending beyond Python
source code and supports the inclusion of native binary extraction and
reverse engineering within UMSAF.

\subsection{Dependency and Software Supply-Chain Visibility Is Important}\label{3-dependency-and-software-supply-chain-visibility-is-important}

Direct and transitive dependencies substantially expand the software
components associated with the experimental application.

SBOM generation and dependency analysis provide mechanisms for making
these relationships more visible and for correlating them with
vulnerability intelligence.

\subsection{Serialized Models Constitute an Additional Analysis Surface}\label{4-serialized-models-constitute-an-additional-analysis-surface}

The experimental CNN produces a serialized \texttt{.keras} model
artifact as part of its normal workflow.

This demonstrates that ML/DL security analysis should consider model
artifacts and their computational structures in addition to source code,
packages, and binaries.

\subsection{Runtime Analysis Is a Planned Complementary Evidence Source}\label{5-runtime-analysis-is-a-planned-complementary-evidence-source}

Structural analysis alone cannot fully characterize runtime interactions
among Python APIs, framework components, native libraries, runtime
environments, and hardware resources.

Runtime observation is expected to provide complementary evidence; its
value must be tested with saved traces in the final evaluation.

\subsection{Cross-Layer Integration Remains To Be Demonstrated}\label{6-cross-layer-integration-remains-to-be-demonstrated}

The preliminary pipeline collected application, dependency,
vulnerability, SBOM, integrity, and binary-inventory evidence. A
serialized model exists, but model inspection, runtime observation, and
integration of these producers have not yet been demonstrated.

This establishes the technical foundation for developing mechanisms that
correlate those findings across the formal UMSAF layers.

\subsection{Existing Techniques Provide Complementary but Fragmented Visibility}\label{7-existing-techniques-provide-complementary-but-fragmented-visibility}

Dependency scanning, SBOM analysis, vulnerability intelligence, binary
inspection, reverse engineering, serialized-model analysis, and runtime
observation provide different forms of security evidence.

These techniques motivate the research premise that correlation may
improve visibility; the preliminary run does not yet establish that
improvement.

\section{Implications for UMSAF}\label{911-implications-for-umsaf}

The preliminary study establishes an initial feasibility basis for the
continued design, implementation, and empirical evaluation of UMSAF.

Specifically, the recorded pipeline provides evidence for:

\begin{itemize}
\tightlist
\item
  identifying direct and transitive dependencies
\item
  generating SBOM information
\item
  correlating dependencies with vulnerability intelligence
\item
  identifying and extracting native binary components
\item
  recording PE structure, imports, exports, and hashes (not completed
  targeted reverse engineering)
\item
  identifying serialized model artifacts as a distinct analysis surface
\item
  deriving a proposed containment policy (not runtime tracing or
  enforcement)
\item
  and collecting preliminary inputs for future cross-layer correlation.
\end{itemize}

These preliminary findings do not establish overall UMSAF effectiveness
or implementation/integration of all required capabilities. Tasks 11-28
must produce and validate the missing evidence before those claims can
be made.

The subsequent research will extend these preliminary capabilities into
an integrated framework and empirically evaluate whether cross-layer
correlation provides more comprehensive security visibility,
vulnerability characterization, traceability, and risk identification
than isolated security-analysis approaches.

Accordingly, the preliminary study provides the experimental foundation
for transitioning from independently applied security-analysis
techniques toward the proposed \textbf{Unified Multilayer Security
Analysis Framework (UMSAF)}.

\chapter{Proposed Framework}\label{10-proposed-framework}

\section{Overview}\label{101-overview}

This research proposes a \textbf{Unified Multilayer Security Analysis
Framework (UMSAF)} for cross-layer analysis of Python-based ML/DL
software ecosystems and their interconnected dependencies, serialized
model artifacts, native C/C++ components, runtime environments, hardware
acceleration, and supporting infrastructure. The primary design and
empirical evaluation are centered on the selected TensorFlow/Keras
deep-learning case; the study does not claim that this case validates
UMSAF across the broader ML/DL domain.

UMSAF integrates software composition and dependency analysis, Software
Bill of Materials (SBOM) generation, vulnerability intelligence, package
and artifact integrity analysis, native binary extraction, binary
reverse engineering, serialized model and computational graph
inspection, runtime behavior analysis, cross-layer interaction modeling,
vulnerability propagation analysis, and risk prioritization.
Interpretable rule-based prioritization is the required approach;
machine-learning-based classification or prioritization remains a
separate, conditional extension if defensible labeled data is adequate.

The primary goal of UMSAF is to overcome the limitations of isolated
security-analysis approaches by providing \textbf{end-to-end security
visibility across multiple interconnected layers of the software
ecosystem}.

Consistent with the threat model and conceptual architecture defined in
Section 6, UMSAF analyzes the following seven layers:

\textbf{Application Layer} ↓ \textbf{Dependency / Software Supply-Chain
Layer} ↓ \textbf{ML/DL Framework Layer} ↓ \textbf{Serialization /
Computational Graph Layer} ↓ \textbf{Native Binary Layer} ↓
\textbf{Runtime / Hardware Acceleration Layer} ↓ \textbf{Operating
System / Infrastructure Layer}

These layers represent the \textbf{security-analysis domain} of UMSAF.
They should be distinguished from the analytical components of the
framework, which collect, process, correlate, and evaluate security
evidence associated with those layers.

Unlike approaches that analyze source code, dependencies, serialized
models, native binaries, runtime behavior, or infrastructure
independently, UMSAF is designed as a \textbf{holistic cross-layer
security-assurance framework} capable of correlating findings across
multiple layers.

Through this correlation, UMSAF aims to provide a more comprehensive
understanding of how vulnerabilities and security-relevant behaviors
originate, propagate, interact, and manifest throughout ML/DL software
ecosystems.

The framework will be developed and empirically evaluated primarily
using a controlled deep-learning workload: a \textbf{deep
learning--based image classification application implemented as a
Convolutional Neural Network (CNN) with TensorFlow/Keras using the
CIFAR-10 dataset}. Additional workloads are optional and should be added
only when needed to evaluate a specific capability and approved as
feasible.

\section{Architectural Design}\label{102-architectural-design}

UMSAF is designed as a modular analytical architecture in which
specialized components generate security evidence that is subsequently
correlated across the ML/DL software ecosystem.

The major analytical components are:

\textbf{ML/DL Application and Software Ecosystem} ↓ \textbf{(1)
Dependency and Software Supply-Chain Intelligence Engine} ↓ \textbf{(2)
SBOM, Integrity, and Vulnerability Intelligence Analyzer} ↓ \textbf{(3)
Native Binary Extraction and Characterization Module} ↓ \textbf{(4)
Reverse Engineering and Binary Analysis Engine} ↓ \textbf{(5) Serialized
Model and Computational Graph Inspection Engine} ↓ \textbf{(6) Runtime
Behavior Analysis Module} ↓ \textbf{(7) Cross-Layer Interaction and
Vulnerability Propagation Mapper} ↓ \textbf{(8) Integrated Risk Scoring
and ML-Based Classification Engine} ↓ \textbf{Integrated Security Report
and Recommendations}

Although represented sequentially for conceptual clarity, these
components do not necessarily operate as a strictly linear pipeline.
Findings generated by one analytical component may inform or trigger
deeper analysis by another component.

For example, dependency analysis may identify a vulnerable package
containing native binaries. The affected binaries may subsequently be
extracted and examined through reverse engineering. Runtime analysis may
then determine whether the relevant component is dynamically loaded or
executed by the experimental workload. The Cross-Layer Interaction and
Vulnerability Propagation Mapper can correlate these findings to
establish relationships among the affected dependency, native component,
framework operation, and runtime behavior.

This modular design enables UMSAF to combine multiple forms of security
evidence while maintaining traceability between findings and their
corresponding software layers.

\section{Component Descriptions}\label{103-component-descriptions}

\subsection{Dependency and Software Supply-Chain Intelligence Engine}\label{1031-dependency-and-software-supply-chain-intelligence-engine}

The \textbf{Dependency and Software Supply-Chain Intelligence Engine}
identifies and characterizes the direct and transitive software
dependencies associated with an ML/DL application.

The component provides visibility into the \textbf{Dependency / Software
Supply-Chain Layer} and establishes relationships between high-level
application components and their supporting packages and artifacts.

\subsubsection{Functions}\label{functions}

\begin{itemize}
\tightlist
\item
  identify direct dependencies
\item
  resolve transitive dependency relationships
\item
  identify package names and versions
\item
  identify package-distribution artifacts
\item
  identify dependencies containing native components
\item
  construct dependency graphs
\item
  collect package metadata
\item
  identify dependency provenance where available
\item
  and support software supply-chain analysis.
\end{itemize}

\subsubsection{Outputs}\label{outputs}

\begin{itemize}
\tightlist
\item
  dependency inventory
\item
  dependency graph
\item
  package metadata
\item
  package-to-artifact mappings
\item
  and software supply-chain relationship data.
\end{itemize}

\subsection{SBOM, Integrity, and Vulnerability Intelligence Analyzer}\label{1032-sbom-integrity-and-vulnerability-intelligence-analyzer}

The \textbf{SBOM, Integrity, and Vulnerability Intelligence Analyzer}
generates standardized software inventories and correlates identified
software components with vulnerability and artifact-integrity
information.

\subsubsection{Functions}\label{functions-1}

\begin{itemize}
\tightlist
\item
  generate SBOMs using formats such as CycloneDX or SPDX
\item
  identify component versions
\item
  query vulnerability intelligence sources
\item
  correlate components with CVE, OSV, and NVD information
\item
  identify known vulnerable versions
\item
  calculate or verify package and artifact hashes where applicable
\item
  identify unexpected or modified artifacts
\item
  and associate vulnerability information with dependency and component
  relationships.
\end{itemize}

\subsubsection{Outputs}\label{outputs-1}

\begin{itemize}
\tightlist
\item
  SBOM
\item
  vulnerability inventory
\item
  component-to-vulnerability mappings
\item
  artifact-integrity information
\item
  and vulnerability correlation data.
\end{itemize}

\subsection{Native Binary Extraction and Characterization Module}\label{1033-native-binary-extraction-and-characterization-module}

The \textbf{Native Binary Extraction and Characterization Module}
identifies and extracts compiled native components associated with
Python packages, ML/DL frameworks, and supporting dependencies.

\subsubsection{Functions}\label{functions-2}

\begin{itemize}
\tightlist
\item
  inspect Python package artifacts
\item
  extract native binary components
\item
  identify \texttt{.so}, \texttt{.pyd}, \texttt{.dll}, and related
  compiled artifacts
\item
  identify compiled Python extensions
\item
  collect cryptographic hashes
\item
  identify binary architecture and format
\item
  collect binary metadata
\item
  inspect dynamically linked dependencies
\item
  and associate native components with their originating packages.
\end{itemize}

\subsubsection{Outputs}\label{outputs-2}

\begin{itemize}
\tightlist
\item
  native binary inventory
\item
  package-to-binary mappings
\item
  binary metadata
\item
  binary hashes
\item
  linked-library information
\item
  and selected artifacts for deeper analysis.
\end{itemize}

\subsection{Reverse Engineering and Binary Analysis Engine}\label{1034-reverse-engineering-and-binary-analysis-engine}

The \textbf{Reverse Engineering and Binary Analysis Engine} examines
selected native components to identify structural and security-relevant
characteristics that may not be visible through application-level or
dependency-level analysis.

Ghidra serves as a primary reverse engineering platform, supplemented
where appropriate by native binary-analysis utilities.

\subsubsection{Functions}\label{functions-3}

\begin{itemize}
\tightlist
\item
  disassemble selected binaries
\item
  decompile selected functions
\item
  inspect imported and exported symbols
\item
  identify linked libraries
\item
  recover function and control-flow information
\item
  inspect native interfaces
\item
  identify potentially security-relevant memory operations
\item
  investigate input-validation logic
\item
  identify potentially unsafe API usage
\item
  and analyze selected vulnerability-relevant execution paths.
\end{itemize}

\subsubsection{Outputs}\label{outputs-3}

\begin{itemize}
\tightlist
\item
  binary-analysis reports
\item
  function and symbol information
\item
  control-flow information
\item
  native dependency relationships
\item
  and security-relevant binary findings.
\end{itemize}

The objective is not to reverse engineer complete ML/DL frameworks.
Instead, analysis will be targeted toward native components identified
as relevant through dependency intelligence, vulnerability information,
serialized-model analysis, or runtime observations.

\subsection{Serialized Model and Computational Graph Inspection Engine}\label{1035-serialized-model-and-computational-graph-inspection-engine}

The \textbf{Serialized Model and Computational Graph Inspection Engine}
analyzes ML/DL model artifacts and their computational structures.

This component provides visibility into the \textbf{Serialization /
Computational Graph Layer} and establishes relationships among model
artifacts, framework functionality, native components, and runtime
execution.

\subsubsection{Functions}\label{functions-4}

\begin{itemize}
\tightlist
\item
  identify serialized model artifacts
\item
  inspect model structure and metadata
\item
  analyze model configuration
\item
  inspect computational graphs where applicable
\item
  enumerate operations and operators
\item
  identify unexpected or potentially security-relevant operations
\item
  detect suspicious or unusual execution relationships
\item
  analyze artifact integrity where applicable
\item
  and map model structures to framework and runtime components where
  feasible.
\end{itemize}

\subsubsection{Outputs}\label{outputs-4}

\begin{itemize}
\tightlist
\item
  serialized-model analysis report
\item
  model and graph structure information
\item
  operator inventory
\item
  artifact metadata
\item
  potentially security-relevant model findings
\item
  and model-to-framework execution mappings.
\end{itemize}

\subsection{Runtime Behavior Analysis Module}\label{1036-runtime-behavior-analysis-module}

The \textbf{Runtime Behavior Analysis Module} observes security-relevant
behavior during ML/DL application execution.

Static analysis provides information about software structure, whereas
runtime analysis provides evidence about which components and execution
paths are actually exercised.

\subsubsection{Functions}\label{functions-5}

\begin{itemize}
\tightlist
\item
  perform runtime tracing
\item
  observe dynamically loaded libraries
\item
  monitor relevant process and thread activity
\item
  observe framework-to-native execution transitions
\item
  monitor relevant file and resource interactions
\item
  identify runtime anomalies
\item
  compare execution behavior across controlled configurations where
  appropriate
\item
  investigate silent or unexpected computational behavior
\item
  and collect evidence concerning runtime and hardware-acceleration
  interactions.
\end{itemize}

\subsubsection{Outputs}\label{outputs-5}

\begin{itemize}
\tightlist
\item
  runtime behavior report
\item
  dynamically loaded component inventory
\item
  execution trace information
\item
  runtime anomaly information
\item
  silent-failure observations where applicable
\item
  and runtime interaction mappings.
\end{itemize}

\subsection{Cross-Layer Interaction and Vulnerability Propagation Mapper}\label{1037-cross-layer-interaction-and-vulnerability-propagation-mapper}

The \textbf{Cross-Layer Interaction and Vulnerability Propagation
Mapper} is a central integrative component of UMSAF.

Its purpose is to correlate evidence generated by the individual
analytical components and model relationships across the seven-layer
ML/DL software ecosystem.

Rather than limiting analysis to interactions between Python code and
native binaries, this component examines relationships spanning:

\textbf{Application Layer} ↓ \textbf{Dependency / Software Supply-Chain
Layer} ↓ \textbf{ML/DL Framework Layer} ↓ \textbf{Serialization /
Computational Graph Layer} ↓ \textbf{Native Binary Layer} ↓
\textbf{Runtime / Hardware Acceleration Layer} ↓ \textbf{Operating
System / Infrastructure Layer}

\subsubsection{Functions}\label{functions-6}

\begin{itemize}
\tightlist
\item
  correlate application dependencies with framework components
\item
  associate dependencies with native binaries
\item
  map application APIs to framework operations
\item
  map framework operations to native components where feasible
\item
  correlate serialized model structures with framework execution
\item
  correlate native components with runtime observations
\item
  associate vulnerabilities with affected components and software layers
\item
  construct cross-layer interaction graphs
\item
  identify potential vulnerability propagation paths
\item
  and preserve traceability among security findings across layers.
\end{itemize}

\subsubsection{Example Execution
Mapping}\label{example-execution-mapping}

\textbf{Python CNN Application} ↓ \textbf{TensorFlow/Keras API} ↓
\textbf{Framework / Wrapper Functionality} ↓ \textbf{Native C/C++
Component} ↓ \textbf{Runtime / System Library / Hardware}

A vulnerability-oriented mapping may additionally represent:

\textbf{Vulnerable Dependency} ↓ \textbf{Affected Native Component} ↓
\textbf{Framework Operation} ↓ \textbf{Runtime Invocation} ↓
\textbf{Potential Application or Infrastructure Impact}

\subsubsection{Outputs}\label{outputs-6}

\begin{itemize}
\tightlist
\item
  cross-layer interaction graph
\item
  component relationship mappings
\item
  execution-flow mappings
\item
  vulnerability propagation paths
\item
  layer-to-layer security relationships
\item
  and traceability information.
\end{itemize}

\subsection{Integrated Risk Scoring and ML-Based Classification Engine}\label{1038-integrated-risk-scoring-and-ml-based-classification-engine}

The \textbf{Integrated Risk Scoring and ML-Based Classification Engine}
aggregates security findings generated throughout UMSAF and supports
risk assessment and vulnerability prioritization.

The component will implement interpretable rule-based scoring and a
severity-only comparison. ML-based classification/prioritization will be
implemented and evaluated only if Task 20 establishes defensible labels
and leakage-resistant splits; otherwise the data-adequacy decision and
limitations must be reported. This conditional branch does not remove
the required H6 evaluation of integrated evidence versus severity-only
prioritization.

\subsubsection{Candidate Features}\label{candidate-features}

Potential features include:

\begin{itemize}
\tightlist
\item
  vulnerability severity
\item
  exploitability information
\item
  dependency depth
\item
  dependency centrality
\item
  component exposure
\item
  package provenance
\item
  binary characteristics
\item
  native-code complexity
\item
  API usage patterns
\item
  serialized-model characteristics
\item
  runtime behavior
\item
  cross-layer reachability
\item
  number of affected layers
\item
  and vulnerability propagation characteristics.
\end{itemize}

\subsubsection{Functions}\label{functions-7}

\begin{itemize}
\tightlist
\item
  aggregate findings from UMSAF analytical components
\item
  normalize security evidence
\item
  calculate rule-based risk scores
\item
  derive cross-layer security features
\item
  assess data adequacy and, if justified, train and evaluate ML-based
  risk-classification or prioritization models
\item
  identify high-risk components
\item
  prioritize vulnerabilities for further investigation
\item
  and provide supporting evidence for assigned risk levels.
\end{itemize}

\subsubsection{Outputs}\label{outputs-7}

\begin{itemize}
\tightlist
\item
  integrated risk scores
\item
  vulnerability prioritization
\item
  component risk rankings
\item
  ML-based classification results if justified, otherwise a documented
  data-adequacy decision
\item
  and supporting risk evidence.
\end{itemize}

\section{Data Flow and Analysis Process}\label{104-data-flow-and-analysis-process}

UMSAF operates through an integrated analysis process rather than a
purely linear vulnerability-scanning pipeline.

A representative analysis workflow is:

\begin{enumerate}
\def\labelenumi{\arabic{enumi}.}
\tightlist
\item
  ingest the Python-based ML/DL application and associated environment
\item
  identify direct and transitive dependencies
\item
  construct dependency relationships
\item
  generate an SBOM
\item
  collect package and artifact integrity information
\item
  correlate components with vulnerability intelligence
\item
  identify and extract native binary components
\item
  characterize selected binaries
\item
  perform targeted reverse engineering and binary analysis
\item
  identify and inspect serialized model artifacts and computational
  structures
\item
  execute the application within a controlled environment
\item
  collect runtime behavior and execution evidence
\item
  correlate evidence across the seven UMSAF layers
\item
  identify potential cross-layer vulnerability propagation paths
\item
  calculate integrated risk scores and vulnerability priorities
\item
  apply and evaluate ML-based classification or prioritization where
  appropriate
\item
  and generate an integrated security report.
\end{enumerate}

The process supports iterative analysis. Findings from later stages may
trigger additional investigation in earlier analytical components.

For example, runtime analysis may identify an unexpected native library,
which may then be returned to the Binary Extraction and Reverse
Engineering components for deeper inspection. Similarly,
serialized-model inspection may identify an operation requiring
investigation of the corresponding framework or native implementation.

\section{Key Research Innovations}\label{105-key-research-innovations}

The proposed UMSAF introduces several research contributions that
collectively distinguish it from isolated security-analysis approaches.

\subsection{Unified Multilayer Security Analysis}\label{1051-unified-multilayer-security-analysis}

UMSAF integrates multiple forms of security analysis, including:

\begin{itemize}
\tightlist
\item
  dependency and software composition analysis
\item
  SBOM analysis
\item
  vulnerability intelligence
\item
  artifact-integrity analysis
\item
  serialized-model analysis
\item
  native binary analysis
\item
  reverse engineering
\item
  runtime behavior analysis
\item
  and integrated risk assessment.
\end{itemize}

Rather than treating these techniques independently, UMSAF correlates
their findings across a common multilayer security model.

\subsection{Cross-Layer Interaction and Vulnerability Propagation Modeling}\label{1052-cross-layer-interaction-and-vulnerability-propagation-modeling}

UMSAF explicitly models security relationships across:

\textbf{Application → Dependency / Supply Chain → ML/DL Framework →
Serialization / Computational Graph → Native Binary → Runtime / Hardware
→ Operating System / Infrastructure}

This capability is intended to provide visibility into how security
findings associated with one layer relate to components and behaviors
occurring in other layers.

\subsection{Integration of Software Composition Analysis and Native Binary Analysis}\label{1053-integration-of-software-composition-analysis-and-native-binary-analysis}

Conventional software composition analysis primarily identifies software
components, versions, dependency relationships, and known vulnerability
metadata.

UMSAF extends this visibility by associating package-level findings with
the native binary components contained within or used by those
dependencies and by enabling selected binaries to undergo deeper
structural and reverse engineering analysis.

\subsection{Integration of Serialized Model Security}\label{1054-integration-of-serialized-model-security}

UMSAF treats serialized model artifacts and computational structures as
first-class components of the ML/DL software ecosystem.

This enables security evidence derived from model artifacts to be
correlated with application, framework, native binary, and runtime
information.

\subsection{Integration of Static and Runtime Security Evidence}\label{1055-integration-of-static-and-runtime-security-evidence}

UMSAF combines structural evidence derived from dependencies, SBOMs,
binaries, and serialized models with behavioral evidence obtained during
runtime execution.

This integration is intended to improve visibility into whether and how
identified components participate in actual ML/DL execution.

\subsection{Native-Level Security Visibility for ML/DL Systems}\label{1056-native-level-security-visibility-for-mldl-systems}

UMSAF provides targeted analysis of native C/C++ components underlying
Python-based ML/DL frameworks.

This capability addresses security characteristics and vulnerabilities
that may not be directly visible through high-level Python source or
dependency metadata alone.

\subsection{ML-Based Risk Classification and Vulnerability Prioritization}\label{1057-ml-based-risk-classification-and-vulnerability-prioritization}

UMSAF will assess whether suitable labeled data supports machine
learning--based experiments on multidimensional security evidence. If it
does not, the research will document that decision and retain the
required rule-based/severity-only comparison; ML is not an unconditional
implemented capability.

The proposed approach uses features derived from dependency
relationships, vulnerability intelligence, binary characteristics,
serialized-model structures, API interactions, runtime observations, and
cross-layer propagation information.

The effectiveness of these mechanisms will be determined through
empirical evaluation rather than assumed as an inherent advantage of
machine learning.

\section{System Implementation Plan}\label{106-system-implementation-plan}

UMSAF will be implemented as a modular research prototype primarily
using Python for orchestration, analysis integration, evidence
normalization, and reporting.

\subsection{Core Technologies}\label{core-technologies}

\begin{itemize}
\tightlist
\item
  Python
\item
  TensorFlow/Keras
\item
  NumPy
\item
  scikit-learn, where appropriate for ML-based classification
  experiments
\end{itemize}

\subsection{Dependency and Software Composition
Analysis}\label{dependency-and-software-composition-analysis}

Potential tools include:

\begin{itemize}
\tightlist
\item
  pip
\item
  pipdeptree
\item
  Python package metadata utilities
\item
  CycloneDX tooling
\item
  SPDX-compatible tooling where appropriate.
\end{itemize}

\subsection{Vulnerability
Intelligence}\label{vulnerability-intelligence}

Potential sources and tools include:

\begin{itemize}
\tightlist
\item
  pip-audit
\item
  OSV
\item
  CVE information
\item
  NVD vulnerability information.
\end{itemize}

\subsection{Binary Analysis and Reverse
Engineering}\label{binary-analysis-and-reverse-engineering}

Potential tools include:

\begin{itemize}
\tightlist
\item
  Ghidra
\item
  \texttt{pefile}
\item
  strings
\item
  Windows dependency-inspection utilities
\item
  and related platform-specific binary-analysis tools.
\end{itemize}

The implementation is platform-specific at the binary-analysis boundary.
The primary Windows implementation fingerprints and characterizes
PE-format \texttt{.pyd} and \texttt{.dll} components, while the UMSAF
evidence schema and cross-layer analysis are designed to support
additional binary formats in future configurations.

\subsection{Serialized Model Analysis}\label{serialized-model-analysis}

Analysis capabilities will include:

\begin{itemize}
\tightlist
\item
  TensorFlow/Keras model inspection
\item
  serialized artifact metadata extraction
\item
  model-structure inspection
\item
  and computational graph or operator analysis where supported by the
  relevant model format.
\end{itemize}

\subsection{Runtime Analysis}\label{runtime-analysis}

Runtime-analysis mechanisms may include:

\begin{itemize}
\tightlist
\item
  process and library tracing
\item
  dynamic dependency observation
\item
  framework execution monitoring
\item
  controlled runtime instrumentation
\item
  and system-level tracing utilities where appropriate.
\end{itemize}

\subsection{Evidence Integration}\label{evidence-integration}

UMSAF will normalize findings from the analytical modules into common
component and relationship representations that support:

\begin{itemize}
\tightlist
\item
  cross-layer correlation
\item
  security finding traceability
\item
  vulnerability propagation analysis
\item
  risk scoring
\item
  ML-based classification and prioritization
\item
  and integrated reporting.
\end{itemize}

The modular architecture is intended to allow individual analytical
tools to be replaced or extended without requiring fundamental changes
to the overall UMSAF conceptual model.

\section{Expected Outputs}\label{107-expected-outputs}

UMSAF is expected to generate a collection of intermediate and
integrated security-analysis artifacts, including:

\begin{enumerate}
\def\labelenumi{\arabic{enumi}.}
\item
  \textbf{Dependency Inventory and Dependency Graph} Identifies direct
  and transitive software relationships.
\item
  \textbf{Software Bill of Materials (SBOM)} Documents software
  components, versions, and dependency metadata.
\item
  \textbf{Vulnerability Intelligence Report} Correlates identified
  software components with known vulnerability information.
\item
  \textbf{Package and Artifact Integrity Information} Records hashes and
  other relevant integrity characteristics for analyzed software
  artifacts.
\item
  \textbf{Native Binary Inventory} Identifies compiled components
  associated with ML/DL packages and frameworks.
\item
  \textbf{Reverse Engineering and Binary Analysis Report} Documents
  security-relevant structural characteristics of selected native
  components.
\item
  \textbf{Serialized Model and Computational Graph Analysis Report}
  Documents model structures, metadata, operations, and potentially
  security-relevant characteristics.
\item
  \textbf{Runtime Behavior Report} Records relevant execution behavior,
  dynamically loaded components, runtime interactions, and identified
  anomalies.
\item
  \textbf{Cross-Layer Interaction Graph} Represents relationships among
  application, dependency, framework, model, binary, runtime, and
  infrastructure components.
\item
  \textbf{Vulnerability Propagation and Traceability Report} Identifies
  potential relationships between vulnerabilities, affected components,
  software layers, and execution paths.
\item
  \textbf{Integrated Risk Assessment and Vulnerability Prioritization
  Report} Provides component-level and vulnerability-level risk
  assessments based on evidence collected across UMSAF.
\item
  \textbf{ML-Based Risk Classification Results} Documents ML evaluation
  when data adequacy justifies it, otherwise records the decision,
  supporting evidence, and limitations. H6 still requires evaluation of
  integrated risk prioritization.
\item
  \textbf{Integrated UMSAF Security Report} Consolidates findings from
  the individual analytical modules into a unified cross-layer security
  assessment with supporting evidence and recommendations.
\end{enumerate}

Collectively, these outputs are intended to provide end-to-end
traceability from high-level ML/DL application components through
dependencies, frameworks, serialized artifacts, native binaries, runtime
execution, and supporting infrastructure.

Below is the \textbf{complete revised Section 11}, aligned with the
finalized Sections 6, 8, 9, and 10.

I have preserved the original methodological intent while revising the
terminology, hypotheses, metrics, experimental procedure, and
cross-layer propagation analysis so they are consistent with the current
UMSAF architecture and do not assume the research outcome in advance.
The original Section 11 already established the \textbf{Design →
Implementation → Experimentation → Evaluation} structure, which remains
appropriate.

\chapter{Research Methodology and Hypotheses}\label{11-research-methodology-and-hypotheses}

\section{Overview}\label{111-overview}

This research adopts a \textbf{design science and empirical evaluation
methodology} to design, implement, and evaluate the proposed
\textbf{Unified Multilayer Security Analysis Framework (UMSAF)}, with
the primary empirical evaluation centered on one controlled
TensorFlow/Keras deep-learning workload.

The study combines framework design, prototype implementation,
controlled experimentation, comparative analysis, and empirical
evaluation using that deep-learning software ecosystem, controlled
security scenarios, selected vulnerable framework or library versions,
known vulnerability information, and security-relevant software
artifacts. Additional workloads may be considered only when needed for a
specific capability and feasible within the approved scope.

The preliminary experimental workload consists of a \textbf{deep
learning--based image classification application implemented as a
Convolutional Neural Network (CNN) with TensorFlow/Keras using the
CIFAR-10 dataset}. This workload provides a controlled and reproducible
environment for examining security evidence across multiple
interconnected layers of an ML/DL software ecosystem.

Consistent with the conceptual architecture defined in Section 6, the
methodology evaluates security across the following seven layers:

\textbf{Application Layer} ↓ \textbf{Dependency / Software Supply-Chain
Layer} ↓ \textbf{ML/DL Framework Layer} ↓ \textbf{Serialization /
Computational Graph Layer} ↓ \textbf{Native Binary Layer} ↓
\textbf{Runtime / Hardware Acceleration Layer} ↓ \textbf{Operating
System / Infrastructure Layer}

The methodology is structured into four major phases:

\textbf{Design → Implementation → Experimentation → Evaluation}

The \textbf{Design} phase defines the UMSAF conceptual architecture,
threat model, analytical components, and cross-layer relationships.

The \textbf{Implementation} phase develops the UMSAF prototype and
integrates the analytical capabilities required to collect and correlate
security evidence.

The \textbf{Experimentation} phase applies UMSAF to the primary
controlled deep-learning software environment and its security
scenarios.

The \textbf{Evaluation} phase assesses the effectiveness, coverage,
traceability, scalability, and analytical value of the integrated
framework relative to isolated security-analysis approaches.

\begin{center}\rule{0.5\linewidth}{0.5pt}\end{center}

\section{Research Design}\label{112-research-design}

\subsection{Design Science Research Approach}\label{1121-design-science-research-approach}

This research follows a \textbf{Design Science Research (DSR)} paradigm
in which the primary research artifact is the proposed UMSAF framework.

The research process includes:

\begin{itemize}
\tightlist
\item
  problem identification and motivation
\item
  definition of research objectives
\item
  UMSAF architecture and artifact design
\item
  prototype implementation
\item
  demonstration using the primary controlled deep-learning workload and
  its security scenarios
\item
  empirical evaluation
\item
  comparative analysis
\item
  and refinement based on experimental findings.
\end{itemize}

The identified research problem concerns the fragmentation of existing
security-analysis techniques across different portions of ML/DL software
ecosystems.

Current approaches may independently examine:

\begin{itemize}
\tightlist
\item
  source code
\item
  dependencies
\item
  known vulnerabilities
\item
  serialized models
\item
  native binaries
\item
  runtime behavior
\item
  or infrastructure components.
\end{itemize}

UMSAF is designed to investigate whether integrating and correlating
evidence from these different forms of analysis provides greater
security visibility than applying them independently.

\subsection{Research Artifact}\label{1122-research-artifact}

The principal research artifact is the \textbf{Unified Multilayer
Security Analysis Framework (UMSAF)}.

UMSAF consists of analytical components supporting:

\begin{itemize}
\tightlist
\item
  dependency and software supply-chain intelligence
\item
  SBOM generation and analysis
\item
  vulnerability intelligence
\item
  package and artifact integrity analysis
\item
  native binary extraction
\item
  binary reverse engineering
\item
  serialized model and computational graph inspection
\item
  runtime behavior analysis
\item
  cross-layer interaction modeling
\item
  vulnerability propagation analysis
\item
  integrated, interpretable rule-based risk assessment
\item
  and, if data adequacy supports it, machine-learning-based risk
  classification or vulnerability prioritization as a separately
  evaluated extension.
\end{itemize}

The artifact will be evaluated not only by whether its individual
components function correctly, but also by whether the integration and
correlation of their outputs provides measurable analytical value.

\begin{center}\rule{0.5\linewidth}{0.5pt}\end{center}

\section{Experimental Subjects and Data Sources}\label{113-experimental-subjects-and-data-sources}

\subsection{Experimental Subjects}\label{1131-experimental-subjects}

The empirical evaluation will use the primary controlled
TensorFlow/Keras deep-learning software environment and security
scenarios selected to exercise UMSAF. Additional workloads are not
assumed and require a specific evaluation need.

The experimental subjects may include the following categories.

\subsubsection{Category A --- Controlled Experimental
Workloads}\label{category-a--controlled-experimental-workloads}

The primary controlled workload is:

\begin{itemize}
\tightlist
\item
  a deep learning--based image classification application implemented as
  a CNN using TensorFlow/Keras and CIFAR-10.
\end{itemize}

Additional controlled workloads may be introduced where necessary to
evaluate specific UMSAF capabilities.

Controlled test cases may also be constructed to examine:

\begin{itemize}
\tightlist
\item
  vulnerable dependencies
\item
  manipulated package configurations
\item
  selected native binary conditions
\item
  modified serialized artifacts
\item
  runtime anomalies
\item
  and controlled cross-layer security scenarios.
\end{itemize}

\subsubsection{Category B --- Selected ML/DL Framework and Dependency
Configurations}\label{category-b--selected-mldl-framework-and-dependency-configurations}

The research may evaluate selected configurations involving:

\begin{itemize}
\tightlist
\item
  TensorFlow
\item
  Keras
\item
  NumPy
\item
  related direct and transitive dependencies
\item
  and, where appropriate, additional ML/DL frameworks or
  interoperability environments.
\end{itemize}

These configurations will be selected based on their relevance to the
research questions and experimental objectives rather than to provide
exhaustive coverage of all ML/DL ecosystems.

\subsubsection{Category C --- Known Vulnerable Framework or Library
Versions}\label{category-c--known-vulnerable-framework-or-library-versions}

Selected historical versions of ML/DL frameworks or dependencies
associated with known vulnerabilities may be used to establish
controlled evaluation cases.

These subjects may include:

\begin{itemize}
\tightlist
\item
  historical TensorFlow versions
\item
  historical NumPy versions
\item
  related supporting libraries
\item
  and components associated with documented CVEs or publicly available
  vulnerability information.
\end{itemize}

The purpose of these subjects is to provide known reference conditions
against which UMSAF analysis results can be evaluated.

\subsection{Data Sources}\label{1132-data-sources}

The research may use information derived from:

\begin{itemize}
\tightlist
\item
  CVE records
\item
  OSV vulnerability data
\item
  NVD vulnerability information
\item
  package metadata
\item
  dependency information
\item
  SBOM records
\item
  framework release information
\item
  publicly documented vulnerability cases
\item
  software repositories
\item
  binary metadata
\item
  serialized model artifacts
\item
  runtime execution traces
\item
  and published vulnerability datasets where their provenance and
  applicability can be verified.
\end{itemize}

Specific quantitative claims obtained from prior vulnerability datasets
will be independently verified against the original source before being
incorporated into the final empirical analysis.

\begin{center}\rule{0.5\linewidth}{0.5pt}\end{center}

\section{Evaluation Objectives}\label{114-evaluation-objectives}

The empirical evaluation is designed to assess several dimensions of
UMSAF.

\subsection{Security Visibility}\label{1141-security-visibility}

Determine whether UMSAF provides security evidence across more portions
of the ML/DL software ecosystem than isolated security-analysis
techniques.

\subsection{Vulnerability Identification and Characterization}\label{1142-vulnerability-identification-and-characterization}

Determine whether integrating dependency, vulnerability, binary, model,
and runtime evidence improves the identification or characterization of
known and controlled security conditions.

\subsection{Cross-Layer Traceability}\label{1143-cross-layer-traceability}

Determine whether UMSAF can establish meaningful relationships among
security findings originating from different software layers.

\subsection{Vulnerability Propagation Analysis}\label{1144-vulnerability-propagation-analysis}

Determine whether security-relevant relationships or potential
propagation paths can be identified across dependencies, frameworks,
serialized artifacts, native binaries, runtime components, and
infrastructure.

\subsection{Risk Prioritization}\label{1145-risk-prioritization}

Evaluate whether integrated cross-layer evidence improves the ability to
prioritize vulnerabilities or security-relevant components relative to
approaches using narrower evidence sources.

\subsection{Scalability}\label{1146-scalability}

Assess the computational and analytical cost of applying UMSAF to
increasingly complex ML/DL environments.

\subsection{Reproducibility}\label{1147-reproducibility}

Assess whether the analytical process can be consistently reproduced
across controlled environments and repeated experimental runs.

\begin{center}\rule{0.5\linewidth}{0.5pt}\end{center}

\section{Evaluation Metrics}\label{115-evaluation-metrics}

The evaluation will use multiple complementary metrics rather than
relying on a single measure of framework effectiveness.

\subsection{Vulnerability Identification Metrics}\label{1151-vulnerability-identification-metrics}

Where suitable ground truth exists, the following metrics may be used:

\begin{itemize}
\tightlist
\item
  \textbf{True Positives (TP):} known or controlled vulnerabilities
  correctly identified
\item
  \textbf{False Positives (FP):} findings incorrectly classified as
  vulnerabilities or security-relevant conditions
\item
  \textbf{False Negatives (FN):} known or controlled vulnerabilities not
  identified
\item
  \textbf{True Negatives (TN):} non-vulnerable conditions correctly
  identified where an appropriate negative dataset exists.
\end{itemize}

From these values, the following metrics may be calculated:

\textbf{Precision}

{[} Precision = \textbackslash frac\{TP\}\{TP + FP\} {]}

\textbf{Recall}

{[} Recall = \textbackslash frac\{TP\}\{TP + FN\} {]}

\textbf{F1 Score}

{[} F1 = 2 \textbackslash times \textbackslash frac\{Precision
\textbackslash times Recall\}\{Precision + Recall\} {]}

These metrics will be used only where the available ground truth permits
defensible classification of findings.

\subsection{Dependency and Software Composition Coverage Metrics}\label{1152-dependency-and-software-composition-coverage-metrics}

Potential measures include:

\begin{itemize}
\tightlist
\item
  percentage of installed dependencies identified
\item
  percentage of direct dependencies resolved
\item
  percentage of transitive dependencies resolved
\item
  dependency-tree depth
\item
  number of identified package relationships
\item
  number of packages containing native components
\item
  percentage of identified components represented within the SBOM
\item
  and percentage of SBOM components successfully correlated with
  vulnerability intelligence.
\end{itemize}

\subsection{Native Binary Analysis Metrics}\label{1153-native-binary-analysis-metrics}

Potential measures include:

\begin{itemize}
\tightlist
\item
  number of native binaries identified
\item
  number or percentage of relevant binaries successfully extracted
\item
  number of binaries mapped to originating packages
\item
  number of dynamically linked relationships identified
\item
  number of selected binaries successfully analyzed
\item
  and number of security-relevant binary characteristics identified.
\end{itemize}

\subsection{Serialized Model and Computational Graph Metrics}\label{1154-serialized-model-and-computational-graph-metrics}

Potential measures include:

\begin{itemize}
\tightlist
\item
  number of serialized artifacts identified
\item
  number of model structures successfully inspected
\item
  number of operators or computational components identified
\item
  percentage of model elements mapped to framework functionality where
  feasible
\item
  and number of potentially security-relevant model characteristics
  identified.
\end{itemize}

\subsection{Runtime Analysis Metrics}\label{1155-runtime-analysis-metrics}

Potential measures include:

\begin{itemize}
\tightlist
\item
  number of dynamically loaded components observed
\item
  number of runtime execution relationships identified
\item
  number of native components observed during execution
\item
  number of runtime anomalies identified
\item
  repeatability of runtime observations
\item
  and number of runtime findings correlated with static evidence.
\end{itemize}

\subsection{Cross-Layer Correlation Metrics}\label{1156-cross-layer-correlation-metrics}

Because cross-layer correlation is central to UMSAF, the evaluation will
measure more than Python-to-native function relationships.

Potential measures include:

\begin{itemize}
\tightlist
\item
  number of mapped cross-layer relationships
\item
  percentage of analyzed components mapped to one or more UMSAF layers
\item
  percentage of security findings correlated across multiple layers
\item
  number of findings with traceability from dependency to affected
  component
\item
  number of findings linked to actual runtime execution
\item
  completeness of identified interaction paths
\item
  number of potential vulnerability propagation paths
\item
  and number of security findings supported by evidence from multiple
  analytical modules.
\end{itemize}

\subsection{Risk Prioritization Metrics}\label{1157-risk-prioritization-metrics}

Where a suitable labeled dataset or defensible evaluation scheme exists,
ML-based prioritization may be evaluated using:

\begin{itemize}
\tightlist
\item
  precision
\item
  recall
\item
  F1 score
\item
  ranking quality
\item
  comparison with vulnerability-severity baselines
\item
  and consistency with known vulnerability characteristics.
\end{itemize}

The evaluation will distinguish between improvements produced by simple
rule-based scoring and those attributable to machine learning--based
classification or prioritization.

\subsection{Performance and Scalability Metrics}\label{1158-performance-and-scalability-metrics}

Potential performance measures include:

\begin{itemize}
\tightlist
\item
  total analysis execution time
\item
  execution time per analytical module
\item
  memory utilization
\item
  storage requirements
\item
  number of analyzed dependencies
\item
  number and size of native binaries
\item
  number of serialized artifacts
\item
  and scaling behavior as software ecosystem complexity increases.
\end{itemize}

\begin{center}\rule{0.5\linewidth}{0.5pt}\end{center}

\section{Baseline and Comparative Analysis}\label{116-baseline-and-comparative-analysis}

UMSAF will be evaluated against selected isolated security-analysis
approaches.

Potential baseline categories include:

\begin{enumerate}
\def\labelenumi{\arabic{enumi}.}
\tightlist
\item
  \textbf{Software Composition Analysis / Dependency Scanning}
\item
  \textbf{Known-Vulnerability Scanning}
\item
  \textbf{Static or Source-Level Analysis}
\item
  \textbf{Native Binary Analysis}
\item
  \textbf{Serialized-Model Analysis}
\item
  \textbf{Runtime or Dynamic Analysis}
\item
  \textbf{Fuzzing or specialized testing techniques where applicable}
\end{enumerate}

The purpose of the comparison is not to demonstrate that any individual
technique is ineffective. Instead, the evaluation will investigate
whether UMSAF provides additional security visibility by
\textbf{integrating and correlating evidence produced by complementary
techniques}.

Comparative analysis may examine differences in:

\begin{itemize}
\tightlist
\item
  component coverage
\item
  vulnerability visibility
\item
  cross-layer traceability
\item
  runtime relevance
\item
  contextual information
\item
  risk prioritization
\item
  and analytical completeness.
\end{itemize}

No predetermined detection advantage will be assumed.

The empirical results will determine whether and to what extent the
integrated UMSAF approach provides measurable benefits over isolated
analysis.

\begin{center}\rule{0.5\linewidth}{0.5pt}\end{center}

\section{Experimental Procedure}\label{117-experimental-procedure}

The empirical evaluation will follow a structured and reproducible
experimental procedure.

\subsection{Step 1 --- Experimental Subject
Selection}\label{step-1--experimental-subject-selection}

Select the primary deep-learning workload, its software configurations,
and known vulnerable components relevant to the evaluation objectives.
Add another workload only if it is necessary for a specific capability
and feasible within the approved scope.

Record:

\begin{itemize}
\tightlist
\item
  framework versions
\item
  package versions
\item
  operating environment
\item
  hardware configuration
\item
  model artifacts
\item
  and experiment-specific parameters.
\end{itemize}

\subsection{Step 2 --- Environment
Construction}\label{step-2--environment-construction}

Create isolated and reproducible Windows experimental environments using
the recorded CPython 3.12.10 interpreter, hash-verified dependency
lockfile, virtual environments, and equivalent isolation mechanisms
where appropriate.

Establish:

\begin{itemize}
\tightlist
\item
  software dependencies
\item
  package versions
\item
  model artifacts
\item
  runtime configuration
\item
  and known experimental conditions.
\end{itemize}

\subsection{Step 3 --- Dependency and Supply-Chain
Analysis}\label{step-3--dependency-and-supply-chain-analysis}

Execute the UMSAF Dependency and Software Supply-Chain Intelligence
Engine to:

\begin{itemize}
\tightlist
\item
  enumerate direct dependencies
\item
  enumerate transitive dependencies
\item
  collect package metadata
\item
  identify package artifacts
\item
  identify packages containing native components
\item
  and construct dependency relationships.
\end{itemize}

\subsection{Step 4 --- SBOM and Vulnerability Intelligence
Analysis}\label{step-4--sbom-and-vulnerability-intelligence-analysis}

Generate an SBOM and correlate components with available vulnerability
intelligence.

Collect:

\begin{itemize}
\tightlist
\item
  component identifiers
\item
  versions
\item
  vulnerability records
\item
  artifact-integrity information
\item
  and component-to-vulnerability relationships.
\end{itemize}

\subsection{Step 5 --- Native Binary Extraction and
Characterization}\label{step-5--native-binary-extraction-and-characterization}

Identify and extract relevant native components.

Collect:

\begin{itemize}
\tightlist
\item
  binary names
\item
  originating packages
\item
  binary formats
\item
  architectures
\item
  hashes
\item
  imported libraries
\item
  exported symbols
\item
  and linked dependencies.
\end{itemize}

\subsection{Step 6 --- Reverse Engineering and Binary
Analysis}\label{step-6--reverse-engineering-and-binary-analysis}

Perform targeted analysis of selected native components using Ghidra and
supporting binary-analysis tools.

Investigate:

\begin{itemize}
\tightlist
\item
  binary structure
\item
  symbols
\item
  imports and exports
\item
  control-flow characteristics
\item
  selected native interfaces
\item
  and security-relevant implementation characteristics.
\end{itemize}

\subsection{Step 7 --- Serialized Model and Computational Structure
Analysis}\label{step-7--serialized-model-and-computational-structure-analysis}

Identify and inspect serialized ML/DL artifacts.

Analyze where applicable:

\begin{itemize}
\tightlist
\item
  model metadata
\item
  architecture
\item
  configuration
\item
  computational structures
\item
  operators
\item
  execution-related characteristics
\item
  and relationships with framework functionality.
\end{itemize}

\subsection{Step 8 --- Runtime Behavior
Analysis}\label{step-8--runtime-behavior-analysis}

Execute selected workloads within controlled environments and collect
runtime evidence.

Potential observations include:

\begin{itemize}
\tightlist
\item
  dynamically loaded libraries
\item
  framework-to-native transitions
\item
  process and thread behavior
\item
  file and resource interactions
\item
  runtime anomalies
\item
  and hardware-acceleration interactions.
\end{itemize}

\subsection{Step 9 --- Cross-Layer
Correlation}\label{step-9--cross-layer-correlation}

Normalize and correlate evidence obtained from the preceding analytical
modules.

Construct relationships across:

\textbf{Application Layer} ↓ \textbf{Dependency / Software Supply-Chain
Layer} ↓ \textbf{ML/DL Framework Layer} ↓ \textbf{Serialization /
Computational Graph Layer} ↓ \textbf{Native Binary Layer} ↓
\textbf{Runtime / Hardware Acceleration Layer} ↓ \textbf{Operating
System / Infrastructure Layer}

\subsection{Step 10 --- Vulnerability Propagation
Analysis}\label{step-10--vulnerability-propagation-analysis}

Identify and evaluate potential security-relevant paths connecting
findings across multiple layers.

Examples may include:

\textbf{Vulnerable Dependency} ↓ \textbf{Affected Native Component} ↓
\textbf{Framework Operation} ↓ \textbf{Runtime Execution}

or:

\textbf{Manipulated Serialized Artifact} ↓ \textbf{Framework /
Computational Operation} ↓ \textbf{Native or Runtime Invocation} ↓
\textbf{System Interaction}

\subsection{Step 11 --- Risk Assessment and
Prioritization}\label{step-11--risk-assessment-and-prioritization}

Aggregate relevant features and calculate integrated risk scores.

Where appropriate, train and evaluate machine learning--based
classification or prioritization mechanisms.

\subsection{Step 12 --- Baseline
Comparison}\label{step-12--baseline-comparison}

Compare UMSAF results with the outputs of selected isolated analysis
approaches.

\subsection{Step 13 --- Metric
Calculation}\label{step-13--metric-calculation}

Calculate applicable:

\begin{itemize}
\tightlist
\item
  precision
\item
  recall
\item
  F1 score
\item
  component coverage
\item
  binary coverage
\item
  cross-layer correlation metrics
\item
  traceability measures
\item
  runtime coverage
\item
  performance measures
\item
  and risk-prioritization metrics.
\end{itemize}

\subsection{Step 14 --- Results
Analysis}\label{step-14--results-analysis}

Analyze:

\begin{itemize}
\tightlist
\item
  where UMSAF provided additional evidence
\item
  where UMSAF provided no measurable advantage
\item
  limitations of individual components
\item
  limitations of cross-layer correlation
\item
  false positives and false negatives
\item
  scalability characteristics
\item
  and implications for the research hypotheses.
\end{itemize}

\begin{center}\rule{0.5\linewidth}{0.5pt}\end{center}

\section{Threats to Validity}\label{118-threats-to-validity}

\subsection{Internal Validity}\label{1181-internal-validity}

Potential threats include:

\begin{itemize}
\tightlist
\item
  incorrect vulnerability-to-component mapping
\item
  incomplete dependency identification
\item
  incomplete native binary extraction
\item
  incorrect cross-layer relationships
\item
  instrumentation effects during runtime analysis
\item
  inaccurate model-artifact interpretation
\item
  or misclassification of findings.
\end{itemize}

\subsubsection{Mitigation}\label{mitigation}

Mitigation strategies include:

\begin{itemize}
\tightlist
\item
  using verified vulnerability information
\item
  validating results using multiple analytical sources
\item
  manually reviewing selected findings
\item
  recording software and framework versions
\item
  repeating experiments
\item
  preserving experimental artifacts
\item
  and distinguishing observed evidence from inferred relationships.
\end{itemize}

\subsection{External Validity}\label{1182-external-validity}

The selected experimental workloads and frameworks may not represent all
ML/DL software ecosystems.

Results obtained from TensorFlow/Keras environments may not generalize
directly to:

\begin{itemize}
\tightlist
\item
  other ML frameworks
\item
  different deployment architectures
\item
  distributed environments
\item
  specialized hardware
\item
  embedded ML/DL systems
\item
  or future framework versions.
\end{itemize}

\subsubsection{Mitigation}\label{mitigation-1}

The study will:

\begin{itemize}
\tightlist
\item
  use representative rather than purely synthetic workloads
\item
  document experimental boundaries
\item
  incorporate selected additional configurations where feasible
\item
  and avoid claiming universal applicability beyond the evaluated
  environments.
\end{itemize}

\subsection{Construct Validity}\label{1183-construct-validity}

Security visibility and software assurance are multidimensional concepts
that cannot be fully characterized by a single metric.

For example, identifying more findings does not necessarily indicate
higher analytical quality.

\subsubsection{Mitigation}\label{mitigation-2}

The evaluation will use complementary measures involving:

\begin{itemize}
\tightlist
\item
  detection
\item
  coverage
\item
  traceability
\item
  runtime relevance
\item
  cross-layer correlation
\item
  performance
\item
  and risk prioritization.
\end{itemize}

\subsection{Conclusion Validity}\label{1184-conclusion-validity}

Experimental results may be affected by:

\begin{itemize}
\tightlist
\item
  small sample sizes
\item
  limited numbers of known vulnerabilities
\item
  incomplete ground truth
\item
  variation among framework versions
\item
  and nondeterministic runtime behavior.
\end{itemize}

\subsubsection{Mitigation}\label{mitigation-3}

Where feasible, the study will:

\begin{itemize}
\tightlist
\item
  repeat experiments
\item
  document experimental parameters
\item
  use known or controlled vulnerability cases
\item
  separate quantitative findings from qualitative observations
\item
  and avoid statistical claims unsupported by the available dataset.
\end{itemize}

\subsection{Reproducibility Threats}\label{1185-reproducibility-threats}

Software ecosystems change over time as:

\begin{itemize}
\tightlist
\item
  packages are updated
\item
  vulnerabilities are reclassified
\item
  dependency relationships change
\item
  tools evolve
\item
  and runtime environments change.
\end{itemize}

\subsubsection{Mitigation}\label{mitigation-4}

The study will record:

\begin{itemize}
\tightlist
\item
  package versions
\item
  framework versions
\item
  tool versions
\item
  OS configuration
\item
  SBOM artifacts
\item
  vulnerability-data timestamps
\item
  and relevant experimental configurations.
\end{itemize}

\begin{center}\rule{0.5\linewidth}{0.5pt}\end{center}

\section{Expected Analytical Outcomes}\label{119-expected-analytical-outcomes}

Because the effectiveness of UMSAF is an empirical research question,
the methodology does not assume that the framework will outperform all
comparison approaches.

Instead, the evaluation is expected to determine whether UMSAF:

\begin{itemize}
\tightlist
\item
  provides security visibility across multiple ML/DL software layers
\item
  exposes relationships among application code, dependencies,
  frameworks, model artifacts, native binaries, and runtime behavior
\item
  provides additional contextual information for known vulnerabilities
\item
  improves traceability across interconnected software components
\item
  identifies potential cross-layer vulnerability propagation paths
\item
  provides useful integrated risk information
\item
  and reveals limitations or tradeoffs associated with cross-layer
  analysis.
\end{itemize}

The research will also document cases in which UMSAF provides limited or
no additional benefit.

Such findings remain important because they define the practical
boundaries and limitations of the proposed approach.

\begin{center}\rule{0.5\linewidth}{0.5pt}\end{center}

\section{Success Criteria}\label{1110-success-criteria}

The research will be considered successful if it demonstrates that the
proposed artifact can be implemented and empirically evaluated against
the research objectives.

Primary success criteria include:

\begin{itemize}
\tightlist
\item
  UMSAF successfully integrates multiple analytical techniques within a
  common framework
\item
  dependencies and software supply-chain relationships can be
  systematically identified
\item
  SBOM and vulnerability information can be correlated with analyzed
  components
\item
  relevant native binaries can be identified and selectively analyzed
\item
  serialized model artifacts can be incorporated into the
  security-analysis process
\item
  runtime behavior can be observed and correlated with structural
  evidence
\item
  cross-layer relationships can be represented and traced
\item
  potential vulnerability propagation paths can be modeled using
  supporting evidence
\item
  integrated risk assessment can be produced from multiple evidence
  sources
\item
  and the framework can be evaluated against appropriate baseline
  approaches.
\end{itemize}

An additional criterion is that the research clearly identifies the
conditions under which UMSAF provides additional security visibility and
the conditions under which its benefits are limited.

Success therefore does not require every hypothesis to be supported.

A rejected or only partially supported hypothesis may still represent a
valid and significant research result.

\begin{center}\rule{0.5\linewidth}{0.5pt}\end{center}

\section{Research Hypotheses}\label{1111-research-hypotheses}

This research investigates whether security assurance for ML/DL software
ecosystems can be improved by integrating and correlating evidence
across multiple interconnected software layers rather than relying
exclusively on isolated analysis techniques.

The following hypotheses guide the empirical evaluation.

\subsection{H1 --- Integrated Multilayer Security Analysis Improves
Security
Visibility}\label{h1--integrated-multilayer-security-analysis-improves-security-visibility}

\textbf{Hypothesis:}

Integrated multilayer security analysis provides greater security
visibility and vulnerability characterization than isolated
security-analysis techniques when applied to the selected ML/DL
experimental environments.

This hypothesis evaluates whether combining dependency intelligence,
vulnerability information, binary analysis, serialized-model inspection,
runtime analysis, and cross-layer correlation provides security evidence
that is unavailable or insufficiently characterized when these
techniques are applied independently.

\subsection{H2 --- Binary Reverse Engineering Provides Additional
Native-Level Security
Visibility}\label{h2--binary-reverse-engineering-provides-additional-native-level-security-visibility}

\textbf{Hypothesis:}

Targeted binary reverse engineering provides additional
security-relevant information about native ML/DL components beyond that
available from high-level application analysis and dependency metadata
alone.

This hypothesis evaluates whether reverse engineering can provide useful
information concerning:

\begin{itemize}
\tightlist
\item
  native implementation structures
\item
  linked components
\item
  low-level interfaces
\item
  execution relationships
\item
  and vulnerability-relevant native characteristics.
\end{itemize}

The evaluation will focus on selected TensorFlow, NumPy, and related
native components rather than assuming that every high-level framework
component corresponds directly to a standalone native binary.

\subsection{H3 --- Cross-Layer Interaction Modeling Improves Security
Traceability}\label{h3--cross-layer-interaction-modeling-improves-security-traceability}

\textbf{Hypothesis:}

Cross-layer interaction modeling improves the traceability of security
findings by establishing relationships among components and behaviors
distributed across multiple layers of ML/DL software ecosystems.

Potential relationships include:

\begin{itemize}
\tightlist
\item
  application-to-dependency relationships
\item
  framework-to-native relationships
\item
  serialized-model-to-framework relationships
\item
  native-to-runtime relationships
\item
  and runtime-to-infrastructure relationships.
\end{itemize}

This hypothesis evaluates whether correlated cross-layer evidence
provides greater contextual traceability than isolated findings.

\subsection{H4 --- Serialized Model Inspection Provides Additional
Security-Relevant
Evidence}\label{h4--serialized-model-inspection-provides-additional-security-relevant-evidence}

\textbf{Hypothesis:}

Serialized-model and computational-structure inspection provides
security-relevant evidence that is not available through
application-source and dependency analysis alone.

The evaluation may consider:

\begin{itemize}
\tightlist
\item
  model structure
\item
  metadata
\item
  computational operations
\item
  unexpected or modified artifact characteristics
\item
  and relationships between serialized artifacts and framework
  execution.
\end{itemize}

This hypothesis does not assume that all serialized models contain
malicious functionality. Instead, it evaluates whether the
serialized-model layer contributes distinct security information to the
overall analysis.

\subsection{H5 --- Runtime Analysis Provides Additional Execution-Level
Security
Evidence}\label{h5--runtime-analysis-provides-additional-execution-level-security-evidence}

\textbf{Hypothesis:}

Runtime behavior analysis provides security-relevant execution evidence
that is not fully observable through static, dependency, binary, or
serialized-artifact analysis alone.

The evaluation may consider:

\begin{itemize}
\tightlist
\item
  dynamically loaded components
\item
  executed native libraries
\item
  framework-to-native transitions
\item
  anomalous runtime behavior
\item
  silent computational deviations where measurable
\item
  and hardware/runtime interactions.
\end{itemize}

This hypothesis evaluates whether runtime evidence improves
understanding of actual software behavior.

\subsection{H6 --- Integrated Cross-Layer Evidence Improves Risk
Prioritization}\label{h6--integrated-cross-layer-evidence-improves-risk-prioritization}

\textbf{Hypothesis:}

Risk prioritization based on integrated cross-layer evidence provides
more contextually informed vulnerability or component prioritization
than prioritization based solely on conventional vulnerability-severity
information.

This hypothesis evaluates whether features such as:

\begin{itemize}
\tightlist
\item
  dependency relationships
\item
  runtime reachability
\item
  native-component involvement
\item
  affected software layers
\item
  model interactions
\item
  and cross-layer propagation characteristics
\end{itemize}

contribute meaningful information to risk assessment.

Where machine learning--based classification or prioritization is used,
its performance will be compared with simpler rule-based or
severity-based approaches.

\subsection{H7 --- Holistic Analysis Provides Broader Security Assurance
for the Evaluated Deep-Learning
Ecosystem}\label{h7--holistic-analysis-provides-broader-security-assurance-for-the-evaluated-deep-learning-ecosystem}

\textbf{Hypothesis:}

Analyzing the evaluated deep-learning application and its supporting
software ecosystem as an interconnected system provides broader security
assurance than evaluating its individual components or layers
independently.

This hypothesis represents the overarching proposition of the
dissertation and integrates the findings associated with H1 through H6.

Collectively, these hypotheses guide the design, implementation,
experimentation, and evaluation of UMSAF.

The empirical results will determine whether each hypothesis is
supported, partially supported, or not supported.

\begin{center}\rule{0.5\linewidth}{0.5pt}\end{center}

\section{Cross-Layer Vulnerability Propagation Analysis}\label{1112-cross-layer-vulnerability-propagation-analysis}

A central research premise investigated in this dissertation is that
security risks within ML/DL software ecosystems \textbf{may propagate
across multiple interconnected software and runtime layers rather than
remaining isolated within the component in which they originate}.

UMSAF therefore incorporates cross-layer vulnerability propagation
analysis to investigate relationships across the formal seven-layer
architecture:

\textbf{Application Layer} ↓ \textbf{Dependency / Software Supply-Chain
Layer} ↓ \textbf{ML/DL Framework Layer} ↓ \textbf{Serialization /
Computational Graph Layer} ↓ \textbf{Native Binary Layer} ↓
\textbf{Runtime / Hardware Acceleration Layer} ↓ \textbf{Operating
System / Infrastructure Layer}

The objective is not to assume that every vulnerability propagates
across multiple layers.

Instead, the analysis seeks to determine:

\begin{itemize}
\tightlist
\item
  whether a finding is isolated to a single component
\item
  whether affected components participate in other software layers
\item
  whether the vulnerable component is reachable during application
  execution
\item
  whether the finding influences framework, native, model, runtime, or
  infrastructure behavior
\item
  and whether a defensible propagation path can be established using
  available evidence.
\end{itemize}

\subsection{Cross-Layer Evidence Sources}\label{11121-cross-layer-evidence-sources}

UMSAF correlates evidence derived from:

\begin{itemize}
\tightlist
\item
  application dependencies
\item
  package metadata
\item
  SBOM records
\item
  vulnerability intelligence
\item
  framework APIs
\item
  serialized model structures
\item
  computational operations
\item
  native binary metadata
\item
  binary linkage relationships
\item
  reverse engineering results
\item
  dynamically loaded libraries
\item
  runtime traces
\item
  hardware/runtime interactions
\item
  and operating system or infrastructure dependencies.
\end{itemize}

Each source provides a different form of evidence.

Cross-layer propagation analysis attempts to combine these evidence
sources into traceable relationships.

\subsection{Dependency-to-Runtime Propagation}\label{11122-dependency-to-runtime-propagation}

One potential propagation scenario involves a vulnerable or compromised
dependency.

For example:

\textbf{Vulnerable or Compromised Dependency} ↓ \textbf{Affected
Framework or Native Component} ↓ \textbf{Framework Invocation} ↓
\textbf{Runtime Execution} ↓ \textbf{Potential Application or System
Impact}

The analysis will investigate whether the affected dependency is:

\begin{itemize}
\tightlist
\item
  installed
\item
  reachable
\item
  imported
\item
  linked
\item
  loaded
\item
  and executed within the experimental workload.
\end{itemize}

This distinction is important because the mere presence of a vulnerable
component does not necessarily establish exploitability or runtime
impact.

\subsection{Serialized-Artifact Propagation}\label{11123-serialized-artifact-propagation}

Another potential propagation scenario involves a manipulated or
security-relevant serialized model artifact.

For example:

\textbf{Serialized Model Artifact} ↓ \textbf{Framework Loading or
Interpretation} ↓ \textbf{Computational Operation} ↓ \textbf{Native /
Runtime Invocation} ↓ \textbf{Potential System Interaction}

UMSAF will investigate whether model structures or operations can be
associated with:

\begin{itemize}
\tightlist
\item
  specific framework functionality
\item
  native components
\item
  runtime execution
\item
  or external resource interactions.
\end{itemize}

The precise propagation path will depend on the serialization format and
framework behavior observed in the selected experimental environment.

\subsection{Native Component Propagation}\label{11124-native-component-propagation}

A vulnerability identified within a native binary may be mapped through:

\textbf{Native Binary} ↓ \textbf{Framework Functionality} ↓
\textbf{Application API Usage} ↓ \textbf{Runtime Reachability}

This analysis can help distinguish between:

\begin{itemize}
\tightlist
\item
  vulnerable native components merely present in the environment
\item
  components reachable through the selected workload
\item
  and components actually observed during runtime execution.
\end{itemize}

\subsection{Runtime Correlation}\label{11125-runtime-correlation}

Runtime observations will be correlated with structural findings
obtained from other UMSAF components.

For example:

\textbf{SBOM Vulnerability Record} ↓ \textbf{Affected Package} ↓
\textbf{Native Binary} ↓ \textbf{Dynamically Loaded Component} ↓
\textbf{Observed Runtime Execution}

Such correlation can provide stronger evidence that a vulnerability is
relevant to the analyzed workload than package-presence information
alone.

\subsection{Propagation Confidence}\label{11126-propagation-confidence}

Because some relationships may be directly observed while others are
inferred, UMSAF should distinguish different levels of propagation
evidence.

Potential categories include:

\begin{itemize}
\tightlist
\item
  \textbf{Directly Observed:} supported by runtime or explicit
  structural evidence
\item
  \textbf{Strongly Correlated:} supported by multiple independent
  evidence sources
\item
  \textbf{Potential:} technically plausible but not directly observed
\item
  \textbf{Unconfirmed:} insufficient evidence to establish propagation.
\end{itemize}

This classification can reduce the risk of presenting inferred
cross-layer relationships as verified execution paths.

\subsection{Example Propagation Paths}\label{11127-example-propagation-paths}

A representative dependency-oriented path is:

\textbf{Compromised Python Package} ↓ \textbf{Malicious or Vulnerable
Native Component} ↓ \textbf{Framework / Runtime Invocation} ↓
\textbf{Security-Relevant Runtime Behavior} ↓ \textbf{Potential System
Impact}

A representative serialized-artifact path is:

\textbf{Manipulated Serialized Model Artifact} ↓ \textbf{Framework /
Computational Structure Processing} ↓ \textbf{Native or Runtime
Invocation} ↓ \textbf{Potential Unauthorized or Unexpected System
Interaction}

A known-vulnerability correlation path may be:

\textbf{CVE / OSV Record} ↓ \textbf{Affected Dependency Version} ↓
\textbf{Associated Native Binary} ↓ \textbf{Framework Execution Path} ↓
\textbf{Observed Runtime Component}

These examples represent analytical patterns to be investigated rather
than assumed outcomes.

\subsection{Relationship to Risk Assessment}\label{11128-relationship-to-risk-assessment}

Cross-layer propagation analysis contributes directly to integrated risk
assessment.

A vulnerability associated with a transitive dependency but not
reachable within the experimental workload may receive different
contextual treatment from a vulnerability that:

\begin{itemize}
\tightlist
\item
  affects a loaded component
\item
  participates in an observed execution path
\item
  crosses multiple software layers
\item
  or interacts with runtime or infrastructure resources.
\end{itemize}

Accordingly, propagation-related features may contribute to risk
classification and vulnerability prioritization.

These features may include:

\begin{itemize}
\tightlist
\item
  number of affected layers
\item
  runtime reachability
\item
  dependency depth
\item
  execution-path participation
\item
  native-component involvement
\item
  artifact relationships
\item
  and evidence strength.
\end{itemize}

\subsection{Research Significance}\label{11129-research-significance}

Cross-layer vulnerability propagation analysis operationalizes one of
the central concepts of UMSAF: security findings should be evaluated not
only as isolated vulnerabilities but also in the context of their
relationships with other components and execution layers.

By correlating evidence across dependencies, framework components,
serialized artifacts, native binaries, runtime behavior, and supporting
infrastructure, UMSAF seeks to provide a more complete and traceable
understanding of security conditions within ML/DL software ecosystems.

The empirical evaluation will determine whether this cross-layer
correlation provides measurable improvements in security visibility,
traceability, vulnerability characterization, and risk prioritization
compared with isolated security-analysis approaches.

\chapter{Expected Contributions}\label{12-expected-contributions}

Collectively, the expected contributions of this research aim to advance
\textbf{machine learning (ML) and deep learning (DL) software security}
from isolated vulnerability detection toward a more holistic model of
\textbf{multilayer software ecosystem assurance}. The dissertation is
expected to contribute to software security, ML/DL security, software
engineering, software supply-chain security, and security analysis by
developing and empirically evaluating methods for collecting,
integrating, and correlating security evidence across interconnected
software layers.

The expected contributions are organized into four primary categories:
\textbf{theoretical, methodological, empirical, and practical
contributions}. Additional supporting contributions are described in
Section 12.5.

\section{Theoretical Contributions}\label{121-theoretical-contributions}

\subsection{Multilayer Security Model for ML/DL Software Ecosystems}\label{1211-multilayer-security-model-for-mldl-software-ecosystems}

This research proposes a formalized multilayer security model that
characterizes the interconnected software and execution layers
comprising modern ML/DL software ecosystems.

Consistent with the threat model and UMSAF architecture established in
this dissertation, the model consists of seven layers:

\textbf{Application Layer} ↓ \textbf{Dependency / Software Supply-Chain
Layer} ↓ \textbf{ML/DL Framework Layer} ↓ \textbf{Serialization /
Computational Graph Layer} ↓ \textbf{Native Binary Layer} ↓
\textbf{Runtime / Hardware Acceleration Layer} ↓ \textbf{Operating
System / Infrastructure Layer}

The proposed model provides a conceptual foundation for examining how
vulnerabilities and security-relevant behaviors may originate, interact,
propagate, and manifest across different portions of an ML/DL software
ecosystem.

Rather than treating a Python-based ML/DL application as an isolated
program, the model represents the application as part of an
interconnected software ecosystem involving third-party dependencies,
ML/DL frameworks, serialized artifacts, native implementations, runtime
environments, hardware-acceleration mechanisms, and supporting
infrastructure.

The expected theoretical contribution is therefore a structured model
for reasoning about \textbf{cross-layer security relationships, trust
boundaries, attack surfaces, vulnerability propagation, and software
assurance within ML/DL software ecosystems}.

\subsection{Cross-Layer Vulnerability and Security-Risk Taxonomy}\label{1212-cross-layer-vulnerability-and-security-risk-taxonomy}

This dissertation is expected to develop a cross-layer taxonomy for
characterizing vulnerabilities and security-relevant conditions within
ML/DL software ecosystems.

The taxonomy will classify findings according to dimensions such as:

\begin{itemize}
\tightlist
\item
  affected software layer
\item
  component type
\item
  vulnerability or security-risk category
\item
  root cause
\item
  execution context
\item
  dependency relationship
\item
  runtime reachability
\item
  potential propagation behavior
\item
  and evidence source.
\end{itemize}

The taxonomy may characterize findings across:

\begin{itemize}
\tightlist
\item
  application code
\item
  dependencies and software supply chains
\item
  ML/DL frameworks
\item
  serialized model artifacts and computational structures
\item
  native C/C++ components
\item
  runtime and hardware-acceleration environments
\item
  and operating system or infrastructure components.
\end{itemize}

Relevant security categories may include:

\begin{itemize}
\tightlist
\item
  dependency vulnerabilities
\item
  software supply-chain weaknesses
\item
  package or artifact integrity concerns
\item
  memory-safety weaknesses
\item
  input-validation weaknesses
\item
  serialization and deserialization risks
\item
  suspicious or unexpected model operations
\item
  runtime anomalies
\item
  execution inconsistencies
\item
  vulnerable native components
\item
  and infrastructure-related security conditions.
\end{itemize}

The taxonomy is intended to provide a systematic mechanism for relating
individual security findings to their broader ML/DL software ecosystem
context.

\subsection{Integrated Multilayer Security-Assurance Paradigm}\label{1213-integrated-multilayer-security-assurance-paradigm}

This research proposes an integrated security-analysis paradigm that
treats ML/DL security as a \textbf{multilayer software ecosystem
assurance problem} rather than solely as a source-code, model-security,
dependency, binary-security, or vulnerability-detection problem.

The proposed paradigm integrates:

\begin{itemize}
\tightlist
\item
  software composition and dependency analysis
\item
  SBOM generation and analysis
\item
  vulnerability intelligence
\item
  software supply-chain analysis
\item
  package and artifact integrity analysis
\item
  native binary analysis
\item
  binary reverse engineering
\item
  serialized model and computational structure inspection
\item
  runtime behavior analysis
\item
  cross-layer interaction modeling
\item
  vulnerability propagation analysis
\item
  and integrated risk assessment.
\end{itemize}

The theoretical contribution of this paradigm is the explicit
recognition that different analytical techniques provide complementary
forms of evidence and that their security value may be increased by
correlating findings across a common multilayer architecture.

This perspective extends beyond isolated vulnerability detection by
emphasizing \textbf{relationships, reachability, propagation, execution
context, traceability, and integrated security evidence}.

\section{Methodological Contributions}\label{122-methodological-contributions}

\subsection{Unified Multilayer Security Analysis Framework (UMSAF)}\label{1221-unified-multilayer-security-analysis-framework-umsaf}

The principal methodological contribution of this dissertation is the
proposed \textbf{Unified Multilayer Security Analysis Framework
(UMSAF)}.

UMSAF is designed as a modular analytical framework integrating:

\begin{itemize}
\tightlist
\item
  dependency and software supply-chain intelligence
\item
  SBOM generation and analysis
\item
  vulnerability intelligence and correlation
\item
  package and artifact integrity analysis
\item
  native binary extraction and characterization
\item
  reverse engineering and binary analysis
\item
  serialized model and computational graph inspection
\item
  runtime behavior analysis
\item
  cross-layer interaction mapping
\item
  vulnerability propagation analysis
\item
  integrated risk assessment
\item
  and, if data adequacy supports it, machine-learning-based risk
  classification or vulnerability prioritization as an optional
  extension.
\end{itemize}

UMSAF is not intended to replace specialized security-analysis tools.
Instead, it provides an integration and correlation methodology through
which evidence generated by different analytical techniques can be
associated with components and relationships across the seven-layer
ML/DL software ecosystem model.

The expected methodological contribution is a reproducible approach for
moving from independently generated security findings toward
\textbf{integrated cross-layer security evidence, contextualization, and
traceability}.

\subsection{Cross-Layer Interaction Mapping Methodology}\label{1222-cross-layer-interaction-mapping-methodology}

This research proposes a methodology for identifying and correlating
relationships among components distributed across multiple ML/DL
software layers.

The methodology extends beyond Python-to-native call mapping and
considers relationships across the complete UMSAF architecture:

\textbf{Application Layer} ↓ \textbf{Dependency / Software Supply-Chain
Layer} ↓ \textbf{ML/DL Framework Layer} ↓ \textbf{Serialization /
Computational Graph Layer} ↓ \textbf{Native Binary Layer} ↓
\textbf{Runtime / Hardware Acceleration Layer} ↓ \textbf{Operating
System / Infrastructure Layer}

The methodology is expected to correlate evidence involving:

\begin{itemize}
\tightlist
\item
  application dependencies
\item
  package relationships
\item
  framework APIs
\item
  serialized model structures
\item
  computational operations
\item
  native binary components
\item
  dynamically linked libraries
\item
  runtime execution
\item
  hardware-acceleration interactions
\item
  and supporting infrastructure.
\end{itemize}

This contribution is intended to improve the contextual interpretation
and traceability of security findings by identifying how components
observed through one analytical technique relate to evidence obtained
from other layers.

\subsection{Cross-Layer Vulnerability Propagation Analysis Methodology}\label{1223-cross-layer-vulnerability-propagation-analysis-methodology}

The dissertation proposes a methodology for investigating
\textbf{whether and how} security risks may propagate across
interconnected ML/DL software layers.

Rather than assuming that every vulnerability produces cross-layer
impact, the methodology distinguishes among findings that are:

\begin{itemize}
\tightlist
\item
  isolated to a particular component
\item
  structurally connected to other components
\item
  reachable through application or framework execution
\item
  observed during runtime
\item
  or associated with a defensible cross-layer propagation path.
\end{itemize}

A dependency-oriented propagation pattern may be represented as:

\textbf{Vulnerable or Compromised Dependency} ↓ \textbf{Affected
Framework or Native Component} ↓ \textbf{Framework Invocation} ↓
\textbf{Runtime Execution} ↓ \textbf{Potential Application or System
Impact}

A serialized-artifact-oriented pattern may be represented as:

\textbf{Manipulated Serialized Model Artifact} ↓ \textbf{Framework /
Computational Structure Processing} ↓ \textbf{Native or Runtime
Invocation} ↓ \textbf{Potential Security-Relevant System Interaction}

The methodology further distinguishes among \textbf{directly observed,
strongly correlated, potential, and unconfirmed relationships},
consistent with the propagation-evidence approach defined in Section 11.

This distinction is intended to prevent technically plausible
relationships from being represented as empirically verified propagation
paths without sufficient supporting evidence.

\subsection{Cross-Layer Risk Classification and Vulnerability Prioritization}\label{1224-cross-layer-risk-classification-and-vulnerability-prioritization}

This dissertation is expected to develop and evaluate a
risk-classification and vulnerability-prioritization methodology that
incorporates evidence from multiple UMSAF layers.

Potential risk features include:

\begin{itemize}
\tightlist
\item
  vulnerability severity
\item
  exploitability information
\item
  dependency depth
\item
  dependency relationships
\item
  component exposure
\item
  package provenance
\item
  native-component involvement
\item
  binary characteristics
\item
  serialized-model characteristics
\item
  API usage patterns
\item
  runtime reachability
\item
  execution-path participation
\item
  number of affected layers
\item
  vulnerability-propagation characteristics
\item
  and strength of supporting evidence.
\end{itemize}

The methodology may combine interpretable rule-based scoring with
machine learning--based classification or prioritization.

A central objective is to investigate whether integrated cross-layer
evidence provides more contextually informative risk prioritization than
approaches based primarily on conventional vulnerability-severity
information.

The effectiveness of machine learning--based classification or
prioritization will be determined through empirical evaluation rather
than assumed.

\section{Empirical Contributions}\label{123-empirical-contributions}

\subsection{Reproducible Deep-Learning Ecosystem Security Evidence Collection}\label{1231-reproducible-deep-learning-ecosystem-security-evidence-collection}

The research is expected to produce a structured collection of
experimental security evidence derived primarily from the controlled
TensorFlow/Keras deep-learning environment and its security scenarios.

The collected evidence may include mappings among:

\begin{itemize}
\tightlist
\item
  software packages and versions
\item
  direct and transitive dependencies
\item
  SBOM components
\item
  CVE, OSV, and related vulnerability records
\item
  native binaries
\item
  package-to-binary relationships
\item
  serialized model artifacts
\item
  computational structures
\item
  runtime observations
\item
  vulnerability propagation relationships
\item
  and cross-layer security evidence.
\end{itemize}

Where known vulnerable framework or library versions are evaluated, the
research will document version-specific relationships between
vulnerability information and affected software components.

The primary contribution is the \textbf{structured, traceable, and
reproducible correlation of experimental security evidence} rather than
the creation of an exhaustive vulnerability dataset covering all ML/DL
frameworks and ecosystems.

\subsection{Comparative Evaluation of Security-Analysis Approaches}\label{1232-comparative-evaluation-of-security-analysis-approaches}

The dissertation will empirically compare UMSAF with selected isolated
security-analysis approaches.

Comparison categories may include:

\begin{itemize}
\tightlist
\item
  software composition and dependency analysis
\item
  known-vulnerability scanning
\item
  static or source-level analysis
\item
  native binary analysis
\item
  serialized-model analysis
\item
  runtime or dynamic analysis
\item
  and fuzzing or specialized testing techniques where applicable.
\end{itemize}

The comparative evaluation will examine dimensions such as:

\begin{itemize}
\tightlist
\item
  component coverage
\item
  vulnerability visibility
\item
  security-context availability
\item
  cross-layer traceability
\item
  runtime relevance
\item
  vulnerability characterization
\item
  vulnerability propagation information
\item
  risk prioritization
\item
  performance
\item
  and analytical completeness.
\end{itemize}

The objective is not to establish in advance that UMSAF will outperform
every individual technique.

Instead, the empirical contribution will determine \textbf{whether,
where, and to what extent integrated cross-layer analysis provides
additional security evidence, traceability, or contextual information
beyond isolated approaches}.

Cases in which UMSAF provides limited or no additional benefit will also
be documented because such findings help establish the practical
boundaries of the proposed methodology.

\subsection{Controlled Case Study of the TensorFlow/Keras Deep Learning Ecosystem}\label{1233-controlled-case-study-of-the-tensorflowkeras-deep-learning-ecosystem}

The primary controlled case study will examine a \textbf{deep
learning--based image classification application implemented as a
Convolutional Neural Network (CNN) with TensorFlow/Keras using the
CIFAR-10 dataset}.

The case study provides a reproducible environment for examining:

\begin{itemize}
\tightlist
\item
  Python application behavior
\item
  TensorFlow/Keras framework interactions
\item
  direct and transitive dependencies
\item
  NumPy and related supporting components
\item
  SBOM generation
\item
  vulnerability intelligence
\item
  native binary extraction
\item
  selected reverse engineering activities
\item
  serialized \texttt{.keras} model artifacts
\item
  computational structures
\item
  runtime behavior
\item
  and cross-layer security relationships.
\end{itemize}

Rather than treating TensorFlow, Keras, and NumPy as necessarily
independent case studies, the research examines them primarily as
interconnected components within the controlled deep learning software
ecosystem.

Selected historical framework or dependency versions associated with
known vulnerabilities may also be incorporated to provide controlled
security-evaluation cases where appropriate.

An additional workload may be incorporated only if required to evaluate
a specific UMSAF capability, approved within the project scope, and
feasible. Such a workload would not by itself establish broad external
validity.

\subsection{Empirical Evidence of Cross-Layer Security Relationships}\label{1234-empirical-evidence-of-cross-layer-security-relationships}

An additional empirical contribution is expected to be evidence
concerning whether and how security findings can be correlated across
multiple ML/DL software layers.

The research will investigate relationships such as:

\begin{itemize}
\tightlist
\item
  vulnerable dependency → affected component
\item
  dependency → native binary
\item
  framework API → native implementation
\item
  serialized artifact → framework operation
\item
  native component → runtime execution
\item
  vulnerability record → runtime-reachable component
\item
  runtime observation → supporting infrastructure
\item
  and cross-layer finding → contextualized risk assessment.
\end{itemize}

This contribution is important because the dissertation does not merely
propose that ML/DL software ecosystems contain multiple interconnected
layers. It empirically investigates whether security evidence from those
layers can be systematically collected, correlated, and evaluated.

\subsection{Empirical Evaluation of Integrated Risk Prioritization}\label{1235-empirical-evaluation-of-integrated-risk-prioritization}

The research is also expected to provide empirical evidence concerning
the usefulness of cross-layer security information for risk
classification and vulnerability prioritization.

The evaluation will investigate whether information such as:

\begin{itemize}
\tightlist
\item
  dependency position
\item
  native-component involvement
\item
  runtime reachability
\item
  affected software layers
\item
  execution-path participation
\item
  vulnerability-propagation evidence
\item
  and evidence strength
\end{itemize}

provides useful contextual information beyond conventional
vulnerability-severity metrics.

Where machine learning--based classification or prioritization is
implemented, its performance will be compared with simpler rule-based or
severity-based approaches.

This contribution will establish whether ML-based prioritization
provides measurable analytical value within UMSAF rather than assuming
that machine learning inherently improves security prioritization.

\section{Practical Contributions}\label{124-practical-contributions}

\subsection{UMSAF Research Prototype}\label{1241-umsaf-research-prototype}

This dissertation is expected to produce a research prototype
implementing the core UMSAF architecture.

The prototype is expected to support capabilities such as:

\begin{itemize}
\tightlist
\item
  dependency analysis
\item
  SBOM generation
\item
  vulnerability correlation
\item
  package and artifact analysis
\item
  native binary inventory generation
\item
  selected binary analysis and reverse engineering
\item
  serialized-model inspection
\item
  runtime evidence collection
\item
  cross-layer relationship mapping
\item
  vulnerability propagation analysis
\item
  integrated risk assessment
\item
  and security reporting.
\end{itemize}

The prototype will primarily serve as an experimental research artifact
for evaluating the feasibility and effectiveness of the proposed
methodology.

Where feasible and appropriate, the implementation may be structured so
that it can be shared or extended for future research. Public
open-source release may be considered separately based on implementation
maturity, licensing restrictions, security considerations, and
dissertation requirements.

\subsection{Deep-Learning Software Security Guidelines}\label{1242-deep-learning-software-security-guidelines}

The research is expected to produce practical security guidance derived
from the UMSAF design and empirical findings.

Potential guidance areas include:

\begin{itemize}
\tightlist
\item
  secure dependency management
\item
  SBOM generation and maintenance
\item
  software supply-chain awareness
\item
  vulnerability-intelligence integration
\item
  package and artifact integrity verification
\item
  native-component awareness
\item
  secure model-artifact handling
\item
  serialized-model inspection
\item
  runtime monitoring
\item
  cross-layer vulnerability assessment
\item
  vulnerability-propagation analysis
\item
  and risk-based vulnerability prioritization.
\end{itemize}

These recommendations are intended to support software engineers, ML/DL
developers, security practitioners, and DevSecOps teams responsible for
developing and operating ML/DL-enabled software systems.

\subsection{DevSecOps Integration Blueprint}\label{1243-devsecops-integration-blueprint}

The dissertation is expected to propose a practical blueprint for
integrating selected UMSAF capabilities into CI/CD and DevSecOps
workflows.

Potential integration points include:

\begin{itemize}
\tightlist
\item
  automated dependency inventory
\item
  SBOM generation
\item
  dependency auditing
\item
  vulnerability monitoring
\item
  package and artifact integrity checks
\item
  native binary inventory
\item
  selected binary-analysis gates
\item
  serialized-model validation
\item
  runtime security observation
\item
  cross-layer evidence correlation
\item
  vulnerability-propagation assessment
\item
  and risk-based security reporting.
\end{itemize}

Not every UMSAF capability is expected to execute at every CI/CD stage.

For example, computationally intensive reverse engineering, deep binary
analysis, or runtime experimentation may be more appropriate for
targeted security-analysis stages, scheduled assessments, release gates,
or investigation workflows rather than every source-code commit.

The proposed blueprint will therefore consider both \textbf{security
value and operational feasibility} when identifying appropriate
DevSecOps integration points.

\section{Additional and Supporting Contributions}\label{125-additional-and-supporting-contributions}

In addition to the primary theoretical, methodological, empirical, and
practical contributions, the research is expected to provide several
supporting contributions.

\subsection{Serialized Model Security Analysis Methodology}\label{1251-serialized-model-security-analysis-methodology}

The dissertation is expected to provide a systematic methodology for
incorporating serialized ML/DL artifacts and computational structures
into broader software security analysis.

This methodology will examine how model artifacts can be:

\begin{itemize}
\tightlist
\item
  identified
\item
  structurally inspected
\item
  associated with framework functionality
\item
  correlated with computational operations
\item
  related to native and runtime execution
\item
  and incorporated into cross-layer security analysis.
\end{itemize}

This contribution treats serialized models as a distinct
security-relevant component of the ML/DL software ecosystem rather than
merely as application output.

\subsection{ML/DL Ecosystem Threat-Modeling Framework}\label{1252-mldl-ecosystem-threat-modeling-framework}

The seven-layer security model and associated trust boundaries are
expected to provide a reusable threat-modeling foundation for reasoning
about ML/DL software ecosystems.

The model considers threats involving:

\begin{itemize}
\tightlist
\item
  application code
\item
  dependencies and software supply chains
\item
  ML/DL frameworks
\item
  serialized artifacts and computational structures
\item
  native binaries
\item
  runtime and hardware acceleration
\item
  and operating system or infrastructure components.
\end{itemize}

The threat-modeling framework provides the conceptual foundation for
identifying where security assumptions may fail, where untrusted
artifacts may enter the system, and where vulnerabilities or
security-relevant behaviors may cross abstraction and trust boundaries.

\subsection{Cross-Layer Vulnerability Propagation Model}\label{1253-cross-layer-vulnerability-propagation-model}

The dissertation is expected to provide a structured model for
representing potential vulnerability propagation across ML/DL software
layers.

The model will distinguish among:

\begin{itemize}
\tightlist
\item
  component presence
\item
  structural relationships
\item
  runtime reachability
\item
  observed execution
\item
  correlated evidence
\item
  and potential propagation.
\end{itemize}

It will also distinguish the strength of propagation evidence, including
relationships that are:

\begin{itemize}
\tightlist
\item
  directly observed
\item
  strongly correlated
\item
  potential
\item
  or unconfirmed.
\end{itemize}

This contribution is intended to improve the precision with which
cross-layer security relationships are represented and communicated.

\subsection{Runtime Security Evidence Integration}\label{1254-runtime-security-evidence-integration}

The research is expected to contribute a methodology for incorporating
runtime evidence into broader ML/DL security assessment.

Rather than proposing a new general-purpose runtime anomaly-detection
algorithm, the contribution focuses on correlating runtime observations
with:

\begin{itemize}
\tightlist
\item
  dependency information
\item
  vulnerability intelligence
\item
  framework behavior
\item
  serialized-model information
\item
  native binary findings
\item
  hardware-acceleration interactions
\item
  and cross-layer execution relationships.
\end{itemize}

This methodology is intended to provide evidence about whether
components identified through static or structural analysis participate
in actual application execution.

\subsection{Integrated Testing and Reverse Engineering Methodology}\label{1255-integrated-testing-and-reverse-engineering-methodology}

The dissertation is expected to demonstrate how complementary testing
and analysis techniques can be combined when investigating ML/DL
frameworks and applications.

These techniques may include:

\begin{itemize}
\tightlist
\item
  dependency analysis
\item
  vulnerability scanning
\item
  static inspection
\item
  binary reverse engineering
\item
  serialized-model inspection
\item
  runtime analysis
\item
  controlled vulnerability cases
\item
  and specialized testing techniques where applicable.
\end{itemize}

The contribution is the \textbf{integration and correlation} of these
methods within a common ML/DL software-assurance process rather than the
development of replacements for the individual analytical techniques.

\section{Overall Research Contribution}\label{126-overall-research-contribution}

The overarching contribution of this dissertation is the development of
a \textbf{holistic, cross-layer security-assurance approach intended for
ML/DL software ecosystems}, with its primary empirical evaluation
bounded to the controlled TensorFlow/Keras deep-learning case and any
additional workloads actually evaluated.

UMSAF shifts the unit of security analysis from isolated software
artifacts toward the security relationships among the seven formal
layers:

\textbf{Application Layer} ↓ \textbf{Dependency / Software Supply-Chain
Layer} ↓ \textbf{ML/DL Framework Layer} ↓ \textbf{Serialization /
Computational Graph Layer} ↓ \textbf{Native Binary Layer} ↓
\textbf{Runtime / Hardware Acceleration Layer} ↓ \textbf{Operating
System / Infrastructure Layer}

Through this perspective, the dissertation investigates whether
dependency intelligence, SBOM information, vulnerability intelligence,
package and artifact integrity analysis, binary analysis, reverse
engineering, serialized-model inspection, runtime observation,
cross-layer interaction modeling, vulnerability propagation analysis,
and integrated risk assessment can be combined within a common
analytical framework to provide more comprehensive security visibility
and traceability.

The central contribution is therefore \textbf{not any individual
analytical technique in isolation}. Software composition analysis, SBOM
generation, vulnerability scanning, binary reverse engineering,
serialized-model inspection, and runtime analysis already exist as
specialized techniques. The proposed contribution of UMSAF is the
systematic \textbf{integration, correlation, and contextualization of
evidence generated by these techniques across a formal multilayer ML/DL
software ecosystem model}.

The research does not assume that every vulnerability crosses multiple
layers, that every analytical technique provides additional value in
every case, that every vulnerable component is runtime-reachable, or
that UMSAF will outperform all existing security-analysis tools.
Instead, the empirical evaluation will determine \textbf{whether, where,
and to what extent cross-layer correlation provides measurable
improvements in security visibility, vulnerability characterization,
traceability, and risk prioritization}, as well as where its benefits
are limited.

Ultimately, this dissertation aims to establish a \textbf{scalable,
extensible, reproducible, and empirically grounded foundation for
multilayer security assurance of ML/DL software ecosystems},
contributing to the broader goals of improving the transparency,
traceability, security, and trustworthiness of software systems that
incorporate machine learning and deep learning technologies.

\chapter{Limitations and Assumptions}\label{13-limitations-and-assumptions}

The proposed \textbf{Unified Multilayer Security Analysis Framework
(UMSAF)} is designed to provide integrated security visibility across
multiple layers of machine learning (ML) and deep learning (DL) software
ecosystems. However, the scope, experimental environment, available
security evidence, analytical tools, and evolving nature of ML/DL
technologies introduce limitations that must be considered when
interpreting the research findings.

The following limitations and assumptions define the boundaries within
which UMSAF will be designed, implemented, and empirically evaluated.

\section{Limitations}\label{131-limitations}

\subsection{Incomplete Component and Execution-Path Coverage}\label{1311-incomplete-component-and-execution-path-coverage}

ML/DL software ecosystems may contain components and execution paths
that cannot be completely identified through static, structural, binary,
model, or runtime analysis.

Native libraries may be:

\begin{itemize}
\tightlist
\item
  dynamically loaded
\item
  conditionally loaded
\item
  statically linked
\item
  bundled within other software artifacts
\item
  generated or selected at runtime
\item
  hardware-specific
\item
  or invoked only under particular execution conditions.
\end{itemize}

Similarly, framework execution paths may depend on:

\begin{itemize}
\tightlist
\item
  application inputs
\item
  model configuration
\item
  framework settings
\item
  operating system characteristics
\item
  available hardware
\item
  accelerator configuration
\item
  and runtime optimization decisions.
\end{itemize}

Consequently, UMSAF may not identify every software component or
execution path associated with an analyzed ML/DL environment.

The framework therefore aims to provide \textbf{systematic and traceable
security visibility within the analyzed environment rather than
exhaustive coverage of every possible component and execution path}.

\subsection{Platform and Environment Dependency}\label{1312-platform-and-environment-dependency}

Native binaries, framework components, runtime behavior, and
hardware-acceleration mechanisms may vary significantly across operating
systems, processor architectures, framework builds, and deployment
environments.

Relevant variations may include:

\begin{itemize}
\tightlist
\item
  Linux, Windows, and macOS
\item
  x86-64 and other processor architectures
\item
  CPU and GPU execution
\item
  framework build configurations
\item
  compiler and optimization settings
\item
  GPU drivers
\item
  CUDA and related accelerator libraries
\item
  containerized and non-containerized environments
\item
  and cloud or local infrastructure.
\end{itemize}

A vulnerability, binary component, or runtime behavior observed in one
environment may therefore differ from that observed in another
environment.

The experimental results will be interpreted within the context of the
specific software and hardware configurations used during the study.
Generalization to other platforms or deployment environments will
require additional validation.

\subsection{Reverse Engineering Complexity and Scope}\label{1313-reverse-engineering-complexity-and-scope}

ML/DL frameworks such as TensorFlow contain large and complex native
codebases. Exhaustive reverse engineering of all native framework
binaries, functions, libraries, and execution paths would be impractical
within the scope of this dissertation.

UMSAF therefore uses \textbf{targeted reverse engineering} rather than
attempting complete decompilation or exhaustive analysis of the entire
TensorFlow or related native software ecosystem.

Native components may be selected for deeper analysis based on factors
such as:

\begin{itemize}
\tightlist
\item
  dependency relationships
\item
  known vulnerability information
\item
  framework relevance
\item
  application usage
\item
  serialized-model behavior
\item
  runtime observations
\item
  binary characteristics
\item
  and potential security significance.
\end{itemize}

Reverse engineering results may also be affected by:

\begin{itemize}
\tightlist
\item
  stripped symbols
\item
  compiler optimizations
\item
  static linking
\item
  complex control flow
\item
  dynamically generated behavior
\item
  obfuscation
\item
  and incomplete debugging information.
\end{itemize}

Consequently, reverse engineering findings will represent targeted
security evidence rather than a complete reconstruction of all native
framework functionality.

\subsection{Vulnerability Ground-Truth Limitations}\label{1314-vulnerability-ground-truth-limitations}

Public vulnerability-information sources such as CVE, OSV, and NVD
provide important information for identifying known vulnerabilities, but
they do not constitute complete ground truth for all security weaknesses
within ML/DL software ecosystems.

Potential limitations include:

\begin{itemize}
\tightlist
\item
  delayed vulnerability disclosure
\item
  incomplete vulnerability records
\item
  inconsistent package or version identifiers
\item
  differences among vulnerability databases
\item
  incomplete affected-version information
\item
  vulnerability reclassification
\item
  missing vulnerabilities
\item
  and incomplete mappings between vulnerability records and specific
  native components.
\end{itemize}

Furthermore, the presence of a vulnerable package or library version
does not necessarily establish that the vulnerable code path is
reachable or exercised by the selected experimental workload.

UMSAF therefore distinguishes, where possible, among:

\begin{itemize}
\tightlist
\item
  vulnerable-component presence
\item
  structural relationships
\item
  application or framework reachability
\item
  runtime loading
\item
  observed execution
\item
  and potential security impact.
\end{itemize}

Known vulnerabilities and controlled security conditions will be used as
ground truth where sufficiently reliable evidence is available. Results
requiring inference will be identified accordingly.

\subsection{Dynamic Analysis and Fuzzing Scope}\label{1315-dynamic-analysis-and-fuzzing-scope}

UMSAF integrates multiple forms of security analysis, including:

\begin{itemize}
\tightlist
\item
  static and structural analysis
\item
  dependency and software composition analysis
\item
  binary analysis
\item
  reverse engineering
\item
  serialized-model inspection
\item
  runtime behavior analysis
\item
  and cross-layer evidence correlation.
\end{itemize}

Fuzzing and other specialized dynamic testing techniques may be used
selectively where appropriate to support specific experimental
objectives.

However, exhaustive fuzzing of all ML/DL framework APIs, native
components, operators, serialized artifacts, runtime configurations, and
hardware-specific execution paths is outside the scope of this
dissertation.

Fuzzing is therefore treated as a \textbf{complementary validation or
testing technique rather than a primary analytical component of UMSAF}.

\subsection{Serialized Model and Computational Structure Coverage}\label{1316-serialized-model-and-computational-structure-coverage}

Serialized ML/DL artifacts vary significantly by framework, format,
version, and deployment configuration.

Examples may include:

\begin{itemize}
\tightlist
\item
  \texttt{.keras} model artifacts
\item
  TensorFlow SavedModel representations
\item
  framework-specific model formats
\item
  computational graphs
\item
  model metadata
\item
  configuration information
\item
  and other serialized execution structures.
\end{itemize}

Different serialization formats expose different levels of structural
and computational information. Some artifacts may permit detailed
inspection of model architecture and operations, whereas others may
provide more limited visibility.

Framework versions may also modify:

\begin{itemize}
\tightlist
\item
  serialization formats
\item
  metadata structures
\item
  operator representations
\item
  model-loading behavior
\item
  and runtime interpretation.
\end{itemize}

Consequently, the depth of serialized-model and computational-graph
analysis may vary across experimental subjects.

The research does not attempt to provide exhaustive security analysis of
every ML/DL model serialization format.

\subsection{Runtime Observability Limitations}\label{1317-runtime-observability-limitations}

Runtime analysis provides evidence concerning actual software execution
but may not reveal every possible execution path or behavior.

Runtime observations depend on:

\begin{itemize}
\tightlist
\item
  workload inputs
\item
  application configuration
\item
  framework configuration
\item
  selected model
\item
  operating environment
\item
  available hardware
\item
  execution backend
\item
  and instrumentation capabilities.
\end{itemize}

A component that is not observed during a particular execution does not
necessarily imply that it cannot be executed under different conditions.

Runtime instrumentation may also introduce measurement effects,
performance overhead, or changes in timing that can influence observed
behavior.

Accordingly, runtime evidence will be interpreted as evidence of
\textbf{observed execution within defined experimental conditions},
rather than proof of all possible runtime behavior.

\subsection{Cross-Layer Correlation and Propagation Uncertainty}\label{1318-cross-layer-correlation-and-propagation-uncertainty}

A central capability of UMSAF is the correlation of security evidence
across multiple software layers. However, not every relationship can
necessarily be established with the same level of confidence.

Some relationships may be directly observed through:

\begin{itemize}
\tightlist
\item
  dependency metadata
\item
  binary linkage
\item
  framework instrumentation
\item
  serialized-model structures
\item
  or runtime traces.
\end{itemize}

Other relationships may be inferred from multiple evidence sources
without direct runtime confirmation.

UMSAF therefore distinguishes among cross-layer relationships that are:

\begin{itemize}
\tightlist
\item
  \textbf{Directly Observed} --- supported by explicit runtime or
  structural evidence
\item
  \textbf{Strongly Correlated} --- supported by multiple independent
  evidence sources
\item
  \textbf{Potential} --- technically plausible but not directly observed
\item
  \textbf{Unconfirmed} --- insufficient evidence exists to establish the
  relationship.
\end{itemize}

This distinction is particularly important for vulnerability-propagation
analysis because the presence of a vulnerable dependency or component
does not by itself establish that a vulnerability propagates across
multiple layers or produces application-level or system-level impact.

\subsection{Risk Classification and Vulnerability Prioritization Limitations}\label{1319-risk-classification-and-vulnerability-prioritization-limitations}

The effectiveness of integrated risk scoring and machine learning--based
vulnerability prioritization depends on the availability and quality of
suitable features, labels, vulnerability data, and experimental
evidence.

Potential limitations include:

\begin{itemize}
\tightlist
\item
  limited labeled security data
\item
  class imbalance
\item
  incomplete vulnerability ground truth
\item
  uncertainty in cross-layer relationships
\item
  limited numbers of representative vulnerability cases
\item
  feature-selection bias
\item
  model overfitting
\item
  and differences among framework versions and environments.
\end{itemize}

Machine learning--based prioritization will therefore be evaluated
against appropriate simpler baselines, such as rule-based or
vulnerability-severity-based approaches.

The research does not assume that ML-based classification will
necessarily outperform simpler prioritization methods.

\subsection{Framework Evolution and Reproducibility}\label{13110-framework-evolution-and-reproducibility}

ML/DL frameworks, dependencies, model formats, vulnerability databases,
hardware interfaces, and runtime environments evolve rapidly.

Changes may affect:

\begin{itemize}
\tightlist
\item
  package dependency structures
\item
  native binaries
\item
  framework APIs
\item
  serialization formats
\item
  computational operators
\item
  vulnerability mappings
\item
  runtime behavior
\item
  hardware-acceleration paths
\item
  and security characteristics.
\end{itemize}

Consequently, findings derived from specific framework or dependency
versions may require reevaluation as the ecosystem evolves.

To improve reproducibility, the research will document relevant
experimental information such as:

\begin{itemize}
\tightlist
\item
  framework versions
\item
  package versions
\item
  operating system configuration
\item
  hardware configuration
\item
  analysis-tool versions
\item
  model artifacts
\item
  SBOM records
\item
  vulnerability-data timestamps
\item
  and experiment-specific configurations.
\end{itemize}

Nevertheless, exact reproduction may become more difficult as external
packages, repositories, tools, and vulnerability-information sources
change over time.

\subsection{Generalizability of Experimental Results}\label{13111-generalizability-of-experimental-results}

The primary controlled experimental workload is a \textbf{deep
learning--based image classification application implemented as a CNN
with TensorFlow/Keras using the CIFAR-10 dataset}.

This workload provides a reproducible environment for developing and
evaluating UMSAF, but it does not represent all ML/DL applications,
frameworks, architectures, model types, deployment patterns, or
operational environments.

The findings may not generalize directly to:

\begin{itemize}
\tightlist
\item
  all TensorFlow applications
\item
  other deep learning frameworks
\item
  large language models
\item
  generative AI systems
\item
  distributed training environments
\item
  federated learning systems
\item
  embedded ML systems
\item
  specialized accelerators
\item
  or large-scale production deployments.
\end{itemize}

Additional representative workloads, framework configurations, or known
vulnerability cases may be incorporated where necessary to evaluate
specific UMSAF capabilities.

However, the dissertation does not claim exhaustive generalizability
across all ML/DL software ecosystems.

\subsection{Scope of Security Analysis}\label{13112-scope-of-security-analysis}

UMSAF focuses primarily on \textbf{software ecosystem security and
cross-layer software assurance}.

Although related security areas may inform the research, the
dissertation does not attempt to comprehensively address:

\begin{itemize}
\tightlist
\item
  adversarial robustness optimization
\item
  adversarial example defense
\item
  model accuracy improvement
\item
  federated learning security
\item
  cryptographic privacy-preserving machine learning
\item
  hardware side-channel attacks
\item
  formal verification of neural network correctness
\item
  or exhaustive penetration testing of production ML/DL systems.
\end{itemize}

These boundaries are consistent with the defined research scope and
allow the dissertation to concentrate on security relationships among
application code, dependencies, frameworks, serialized artifacts, native
binaries, runtime environments, and supporting infrastructure.

\section{Assumptions}\label{132-assumptions}

The research is conducted under several assumptions necessary to support
the design, implementation, and empirical evaluation of UMSAF.

\subsection{Python-Based Deep-Learning Workload}\label{1321-python-based-deep-learning-workload}

The experimental scope assumes that the primary deep-learning workload
analyzed in this research is \textbf{Python-based} and uses commonly
available Python packaging and framework mechanisms.

This assumption reflects the selected experimental environment involving
TensorFlow/Keras, NumPy, and related dependencies.

It does not imply that all ML/DL systems are Python-based or that UMSAF
could not be extended to other programming-language ecosystems.

\subsection{Availability of Dependency and Package Metadata}\label{1322-availability-of-dependency-and-package-metadata}

The research assumes that sufficient package and dependency metadata is
available to identify and analyze relevant software components.

This information may be obtained from:

\begin{itemize}
\tightlist
\item
  installed package metadata
\item
  package managers
\item
  package repositories
\item
  lock or requirements files
\item
  distribution metadata
\item
  SBOM tools
\item
  and related dependency-analysis mechanisms.
\end{itemize}

The completeness of dependency analysis depends on the availability and
accuracy of this metadata.

\subsection{Availability of Analyzable Native Components}\label{1323-availability-of-analyzable-native-components}

The research assumes that a sufficient subset of native components used
by the selected ML/DL environments can be identified and accessed for
analysis.

These components may include:

\begin{itemize}
\tightlist
\item
  shared libraries
\item
  compiled Python extensions
\item
  package-distributed native binaries
\item
  dynamically linked libraries
\item
  and other accessible native artifacts.
\end{itemize}

The research does not assume that every native component will be
available as a standalone wheel or directly analyzable binary.

Instead, sufficient native-component accessibility is assumed to permit
evaluation of the native binary and reverse engineering capabilities of
UMSAF.

\subsection{Availability and Reliability of Vulnerability Intelligence}\label{1324-availability-and-reliability-of-vulnerability-intelligence}

The research assumes that vulnerability-information sources such as CVE,
OSV, and NVD provide sufficiently reliable identifiers, affected-version
information, and metadata to support experimental vulnerability
correlation.

This assumption does not imply that these databases are complete or
error-free.

Where inconsistencies exist, vulnerability mappings will be verified
against available supporting information before being treated as ground
truth.

\subsection{Availability of Serialized Model Artifacts}\label{1325-availability-of-serialized-model-artifacts}

The research assumes that selected experimental workloads produce or use
serialized model artifacts that can be accessed for analysis.

For the preliminary TensorFlow/Keras CNN workload, the model is saved as
a \texttt{.keras} artifact.

The research assumes that sufficient structural or metadata information
can be extracted from selected serialized artifacts to evaluate the
serialized-model analysis capabilities of UMSAF.

\subsection{Controlled Runtime Observability}\label{1326-controlled-runtime-observability}

The research assumes that selected workloads can be executed within
controlled environments in which sufficient runtime information can be
observed.

This may include information concerning:

\begin{itemize}
\tightlist
\item
  loaded libraries
\item
  framework execution
\item
  native-component usage
\item
  process and thread behavior
\item
  file or resource interactions
\item
  and hardware-acceleration activity where applicable.
\end{itemize}

The research further assumes that instrumentation overhead can be
managed sufficiently to permit meaningful interpretation of runtime
observations.

\subsection{Sufficient Cross-Layer Identifiers and Relationships}\label{1327-sufficient-cross-layer-identifiers-and-relationships}

UMSAF depends on the ability to associate evidence collected by
different analytical modules.

The research therefore assumes that sufficient identifiers or
relationships can be established among:

\begin{itemize}
\tightlist
\item
  packages
\item
  dependencies
\item
  framework components
\item
  serialized artifacts
\item
  native binaries
\item
  runtime components
\item
  vulnerability records
\item
  and infrastructure resources.
\end{itemize}

These relationships may be established using metadata, hashes, package
ownership, binary linkage, API relationships, model structures, runtime
traces, or other available evidence.

Not every relationship is assumed to be directly observable.

\subsection{Representative Value of Controlled Workloads}\label{1328-representative-value-of-controlled-workloads}

The research assumes that the selected controlled and representative
workloads provide sufficient software complexity and framework
interaction to exercise the principal analytical capabilities of UMSAF.

The TensorFlow/Keras CIFAR-10 CNN is not assumed to represent every
ML/DL application.

Instead, it serves as a \textbf{controlled, reproducible experimental
workload} through which dependency analysis, SBOM generation, native
binary analysis, serialized-model inspection, runtime analysis,
cross-layer correlation, and related UMSAF capabilities can be developed
and evaluated.

\subsection{Integrity of the Controlled Experimental Environment}\label{1329-integrity-of-the-controlled-experimental-environment}

Unless an experiment intentionally introduces a modified, vulnerable, or
malicious component, the research assumes that the controlled
experimental environment has not been independently compromised.

This assumption provides a known baseline against which controlled
security conditions, historical vulnerable components, artifact
modifications, and runtime observations can be compared.

\subsection{Legal and Ethical Availability of Research Artifacts}\label{13210-legal-and-ethical-availability-of-research-artifacts}

The research assumes that the software packages, native binaries, model
artifacts, vulnerability information, and related research materials
selected for analysis can be legally and ethically examined within the
context of academic security research.

Reverse engineering and security analysis will be limited to software
and artifacts for which such analysis is permitted or appropriately
justified within the research environment.

\subsection{Tool Reliability}\label{13211-tool-reliability}

The research assumes that the selected analytical tools provide
sufficiently reliable outputs for use as evidence within UMSAF when
their documented limitations are considered.

Potential tools include dependency-analysis utilities, SBOM generators,
vulnerability scanners, Ghidra, binary-inspection tools,
serialized-model inspection mechanisms, and runtime-analysis utilities.

No individual tool is assumed to provide complete or infallible results.

Where practical, findings will be corroborated using multiple evidence
sources or analytical techniques.

\section{Interpretation of Limitations and Assumptions}\label{133-interpretation-of-limitations-and-assumptions}

The limitations and assumptions described above establish the boundaries
within which the results of this dissertation should be interpreted.

UMSAF is not intended to provide exhaustive vulnerability discovery
across every ML/DL framework, model format, native component, hardware
platform, or deployment environment. Nor does the framework assume that
every identified vulnerability is exploitable, runtime-reachable, or
capable of propagating across multiple software layers.

Instead, the research investigates whether security evidence collected
from multiple analytical perspectives can be systematically integrated
and correlated to improve understanding of security conditions within
ML/DL software ecosystems.

Accordingly, conclusions will distinguish among:

\begin{itemize}
\tightlist
\item
  component presence
\item
  known vulnerability association
\item
  structural relationships
\item
  runtime reachability
\item
  observed execution
\item
  cross-layer correlation
\item
  potential vulnerability propagation
\item
  and demonstrated security impact.
\end{itemize}

This distinction is essential to maintaining appropriate evidentiary
boundaries and avoiding unsupported conclusions regarding
exploitability, propagation, or system impact.

The empirical findings will therefore be interpreted within the context
of the selected experimental workloads, software versions, analytical
tools, runtime configurations, available vulnerability intelligence, and
supporting evidence.

These limitations do not invalidate the proposed framework; rather, they
define the conditions under which UMSAF can be evaluated and establish
appropriate boundaries for interpreting its effectiveness,
generalizability, and contribution to \textbf{multilayer security
assurance for ML/DL software ecosystems}.

\chapter{Conclusion}\label{14-conclusion}

Modern \textbf{machine learning (ML) and deep learning (DL) software
systems} operate as complex, multilayer software ecosystems that extend
beyond high-level application code. Python-based ML/DL applications may
depend on extensive third-party software supply chains, ML/DL
frameworks, serialized model artifacts and computational structures,
native C/C++ components, runtime and hardware-acceleration environments,
and supporting operating system and infrastructure components. As a
result, vulnerabilities and security-relevant behaviors may originate,
interact, or manifest across multiple layers and may not be adequately
characterized when individual components or software layers are analyzed
in isolation.

This dissertation proposes the \textbf{Unified Multilayer Security
Analysis Framework (UMSAF)} to address these security visibility and
correlation challenges. UMSAF models the ML/DL software ecosystem across
seven interconnected layers:

\textbf{Application Layer} ↓ \textbf{Dependency / Software Supply-Chain
Layer} ↓ \textbf{ML/DL Framework Layer} ↓ \textbf{Serialization /
Computational Graph Layer} ↓ \textbf{Native Binary Layer} ↓
\textbf{Runtime / Hardware Acceleration Layer} ↓ \textbf{Operating
System / Infrastructure Layer}

Within this architecture, UMSAF integrates \textbf{software composition
and dependency analysis, Software Bill of Materials (SBOM) generation
and analysis, vulnerability intelligence, software supply-chain
analysis, package and artifact integrity analysis, native binary
extraction and characterization, binary reverse engineering,
serialized-model and computational-graph inspection, runtime behavior
analysis, cross-layer interaction modeling, vulnerability-propagation
analysis, and risk assessment}. Rule-based risk scoring is required;
machine-learning-based prioritization remains an optional extension
subject to the data-adequacy decision.

The central premise of UMSAF is that the security significance of an
individual finding may depend not only on the presence of a
vulnerability or security-relevant condition but also on its
relationship to other components within the software ecosystem. A
vulnerable dependency, for example, may have different security
implications depending on whether the affected component is associated
with a relevant native binary, reachable through framework
functionality, loaded during runtime, observed in execution, or
connected to a broader cross-layer propagation path. Similarly,
serialized model artifacts may interact with framework, native, and
runtime components in ways that cannot be fully characterized through
application-source or dependency analysis alone.

UMSAF therefore aims to move beyond isolated security findings by
correlating evidence across software layers and distinguishing among
\textbf{component presence, structural relationships, runtime
reachability, observed execution, cross-layer correlation, and potential
vulnerability propagation}. Where propagation relationships are
investigated, the framework further distinguishes among directly
observed, strongly correlated, potential, and unconfirmed relationships
to maintain appropriate evidentiary boundaries and avoid treating
plausible execution paths as verified security impacts.

The preliminary pipeline provides dependency, SBOM,
vulnerability-snapshot, integrity-check, static, and PE-inventory
evidence for a \textbf{deep learning CNN implemented with
TensorFlow/Keras using CIFAR-10}. A model artifact exists, but model
inspection, runtime tracing, targeted reverse engineering, and
cross-layer integration remain planned. The environment must be rebuilt
and its evidence refreshed before final evaluation; the preliminary run
does not establish generalizability or integrated UMSAF effectiveness.

The dissertation adopts a \textbf{design science and empirical
evaluation methodology} structured around four major phases:

\textbf{Design → Implementation → Experimentation → Evaluation}

Through this methodology, UMSAF will be implemented as a research
artifact and evaluated primarily using the controlled TensorFlow/Keras
deep-learning workload, selected framework and dependency
configurations, known vulnerability information, serialized model
artifacts, native components, and runtime evidence. The evaluation will
compare UMSAF with selected isolated security-analysis approaches and
examine dimensions including \textbf{security visibility, vulnerability
identification and characterization, cross-layer traceability,
vulnerability-propagation analysis, risk prioritization, scalability,
and reproducibility}. Conclusions will be bounded by the workloads and
evidence actually evaluated.

The empirical evaluation does not assume that UMSAF will outperform
every individual security-analysis technique or that every vulnerability
will propagate across multiple software layers. Instead, the research
will determine \textbf{whether, where, and to what extent integrated
cross-layer security analysis provides additional security evidence,
context, and traceability beyond isolated analysis approaches}. Cases in
which cross-layer correlation provides limited or no additional
analytical benefit will also contribute to understanding the
capabilities and boundaries of the proposed framework.

A further objective of the research is to investigate whether integrated
evidence from dependencies, vulnerability intelligence, native binaries,
serialized artifacts, runtime execution, and cross-layer relationships
can support more contextually informed \textbf{risk classification and
vulnerability prioritization}. Machine learning--based prioritization
will be evaluated alongside rule-based and vulnerability-severity-based
approaches only when the documented data-adequacy decision supports
defensible ML experiments; no superior result is assumed.

This dissertation is situated within the broader field of ML/DL software
ecosystem security, but its primary empirical contribution is focused on
a controlled deep-learning software ecosystem. It treats security as a
\textbf{holistic software ecosystem assurance problem} rather than
solely as a problem of adversarial robustness, source-code security,
vulnerable dependencies, model security, native binary analysis, or
runtime monitoring. UMSAF is designed to integrate these perspectives;
conclusions about its effectiveness will be limited to the workloads and
evidence evaluated.

The overarching contribution of this research is the development and
empirical evaluation of a \textbf{scalable, extensible, reproducible,
and holistic cross-layer security-assurance framework}, with its primary
evaluation conducted on the TensorFlow/Keras deep-learning case. By
integrating and correlating evidence across multiple layers, the
research will assess the \textbf{visibility, traceability,
characterization, and prioritization of security risks} in that case and
provide a bounded foundation for future research into broader ML/DL
software assurance.

\clearpage
\addcontentsline{toc}{chapter}{Bibliography}
\printbibliography[title={Bibliography}]

\end{document}
